\documentclass[11pt,a4paper,reqno]{amsart}
\usepackage[T1]{fontenc}
\usepackage{lmodern,microtype,mathtools,amssymb,booktabs,array,tabularx,enumitem}
\usepackage[a4paper,margin=21.5mm,headheight=10pt]{geometry}
\usepackage{xcolor,needspace}
\usepackage[colorlinks=true,linkcolor=black,citecolor=black,urlcolor=black]{hyperref}
\usepackage[nameinlink,noabbrev]{cleveref}
\hypersetup{pdftitle={The cohomology of  finite simple groups is unbounded}}
\numberwithin{equation}{section}
\newtheorem{theorem}{Theorem}[section]
\newtheorem{maintheorem}{Theorem}
\renewcommand{\themaintheorem}{\Alph{maintheorem}}
\crefname{maintheorem}{Theorem}{Theorems}
\Crefname{maintheorem}{Theorem}{Theorems}
\newtheorem{proposition}[theorem]{Proposition}
\newtheorem{conjecture}[theorem]{Conjecture}
\newtheorem{lemma}[theorem]{Lemma}
\newtheorem{corollary}[theorem]{Corollary}
\newtheorem{question}{Question}
\renewcommand{\thequestion}{\Roman{question}}
\theoremstyle{definition}
\newtheorem*{notation}{Notation}
\newtheorem{definition}[theorem]{Definition}
\newtheorem{example}[theorem]{Example}
\theoremstyle{remark}
\newtheorem{remark}[theorem]{Remark}
\theoremstyle{plain}
\newtheorem*{guralnickconjecture}{Conjecture}
\newtheorem*{hdrconjecture}{Question}
\newcommand{\Z}{\mathbb Z}
\newcommand{\C}{\mathbf C}
\newcommand{\M}{\mathcal M}



\newcommand{\eps}{\varepsilon}
\newcommand{\ord}{\operatorname{ord}_v}
\newcommand{\inv}{\operatorname{inv}}
\newcommand{\A}{\mathsf A}
\newcommand{\B}{\mathsf B}
\newcommand{\CC}{\mathsf C}
\newcommand{\gb}[2]{\genfrac{[}{]}{0pt}{}{#1}{#2}_q}
\newcommand{\mus}{\mu_I}
\newcommand{\Part}{\mathsf p}
\setlength{\emergencystretch}{2em}
\setlist[enumerate]{leftmargin=*,itemsep=3pt}






\newcommand{\D}{\mathcal {BS}}
\newcommand{\R}{\mathcal {BS}}
\newcommand{\V}{\mathcal {CBS}}


\definecolor{nicolas}{HTML}{7030A0}
\usepackage{tikz,enumitem,rotating,latexsym,bm,stmaryrd,caption}
\usetikzlibrary{positioning,intersections,decorations,decorations.pathreplacing, arrows.meta}

\usepackage{tikz,tikz-3dplot}

% Author to-do notes, with the signature at the end.
\setlength{\marginparwidth}{21mm}
\setlength{\marginparsep}{3mm}
\usepackage[textsize=small]{todonotes}
\newcommand{\Chris}[1]{\todo[color=gray!15]{#1\hfill\textbf{C}}}
\newcommand{\Nico}[1]{\todo[color=violet!15]{#1\hfill\textbf{N}}}
\newcommand{\Damian}[1]{\todo[color=orange!20]{#1\hfill\textbf{Dn}}}
\newcommand{\David}[1]{\todo[color=green!15]{#1\hfill\textbf{Dv}}}


\title[The   cohomology of  finite simple groups is unbounded]{The   cohomology of     finite simple groups is unbounded}
% Affiliations and email addresses: place each pair after its author.
% Replace the placeholders and uncomment the fields when ready.




 \author{Chris Bowman}
       \address{Department of Mathematics, 
University of York, Heslington, York,  UK}
\email{Chris.Bowman-Scargill@york.ac.uk}





\author{Dami\'an de la Fuente}
\address{LACIM, Université du Québec à Montréal, Montreal, Canada.}
\email{damiandlfa@gmail.com}

\author{Nicol\'as Libedinsky}
 \address{Departamento de Matem\'aticas, Facultad de Ciencias, Universidad de Chile, Chile}
 \email{nlibedinsky@u.uchile.cl}

\author{David Plaza}
\address{Instituto de Matem\'aticas, Universidad de Talca, Talca, Chile}
\email{dplaza@utalca.cl}
\date{}
\hypersetup{pdftitle={The cohomology of  finite simple groups is unbounded},pdfsubject={A complete coefficient calculation with finite partition extraction}}
\allowdisplaybreaks[1]
% Number only referenced equations; retain support for capitalized cleveref references.
\hypersetup{hypertexnames=false}
\usepackage{autonum}
\makeatletter
\autonum@generatePatchedReferenceCSL{Cref}
\makeatother
\begin{document}





















\begin{abstract}
In every fixed positive degree, we show that the cohomology
of finite simple groups, algebraic groups and quantum groups,
with coefficients in simple modules, can have arbitrarily
large dimension. In particular, we disprove Guralnick's conjecture. 
 \end{abstract}
\maketitle
% \tableofcontents


\vspace{-0.6cm}
\section{Introduction}



\subsection{} 
% Unbounded cohomology is a persistent phenomenon in topology and geometry, even in a fixed degree and with coefficients of rank one:
Cohomology in topology and geometry is rarely bounded, even in a single degree with coefficients of rank one. A smooth projective complex curve of genus $g$ has first cohomology of dimension $2g$ in singular, de Rham and $\ell$-adic cohomology, and of dimension $g$ for the structure sheaf.
In 1984, Guralnick conjectured \cite[Conjecture~2]{Guralnick} that the first cohomology of finite groups behaves very differently:
\begin{guralnickconjecture}[Guralnick]
There is an absolute constant $C$ such that
\[
\dim_k H^1(G,V)\le C
\]
for every finite group $G$, every field $k$, and every faithful
absolutely irreducible finite-dimensional $kG$-module $V$.
\end{guralnickconjecture}


 



Guralnick's conjecture concerns degree one. 
In higher degrees the analogous statement fails for arbitrary finite groups: Guralnick, Kantor, Kassabov and Lubotzky showed that   the dimensions of $H^n(G,V)$ are unbounded  for every $n\ge2$ as $G$ ranges over {\em all finite groups} and $V$ over faithful absolutely irreducible modules \cite[Theorems~12.3 and~12.6]{GKKL2007}.
Their examples use wreath products to build large cohomology, 
and they pose the following question concerning {\em simple} groups:

\begin{question}\cite[Question~12.1]{GKKL2007} \label{question:higher-degree-bound}
  For which $n\geq 1$    is there an absolute constant $C_n$ such that
\[
\dim_k H^n(G,V)\le C_n
\]
for every finite \textbf{simple} group $G$, every field $k$, and every faithful
absolutely irreducible finite-dimensional $kG$-module $V$.
\end{question}



We show that no such constant $C_n$ exists for any $n\geq1$, answering
Question~\ref{question:higher-degree-bound}. The case $n=1$ disproves
Guralnick's conjecture. Before stating our results, we recall the
developments that sustained interest in the conjecture for more than
forty years.
 
\subsection{}
 
 
Despite its audacity, the conjecture looked safe for almost two decades.
Guralnick and Hoffman even suggested that $C=2$ might suffice
\cite[Conjecture 2]{GuralnickHoffman}. Only in 2003 did Scott find modules with
$\dim H^1=3$, for $\mathrm{PSL}_6$ in large defining characteristic
\cite{Scott2003} (see also \cite{BrayWilson}), and six years later $3$ was still the largest
known value \cite[Introduction]{CPS09}. 
%
Meanwhile, a sustained programme of positive results gave bounds
increasingly independent of the module. Early bounds depended on
$\dim V$ \cite{AschbacherGuralnick,GuralnickHoffman}.
For finite groups of Lie type in defining characteristic, Cline,
Parshall and Scott bounded $\dim H^1$ in terms of the root system
alone, uniformly in the characteristic \cite{CPS09}.
In the algebraic group setting, Parshall and Scott obtained
analogous bounds for $\operatorname{Ext}^1$ between arbitrary
irreducible rational modules
\cite[Theorem~5.1]{PS-BoundingExt}.
Explicit bounds for finite groups of Lie type were obtained by
Guralnick and Tiep in cross characteristic \cite{GT11} and by
Parker and Stewart in defining characteristic \cite{PS14}.
Moreover, a universal bound for finite simple groups would imply
one for all finite groups with faithful absolutely irreducible
coefficients \cite[Lemma~5.2]{GKKL2007}.





At the 2012 AIM workshop \emph{Cohomology bounds and growth rates}, Scott
and Sprowl, using a new algorithm for computing individual
Kazhdan--Lusztig basis elements \cite{ScottSprowl}, reached $5$ for
$\mathrm{PSL}_7$ (see \cite[\S4]{Luebeck}), and then L\"ubeck, running
the Kazhdan--Lusztig algorithm at scale, found\[
\dim H^1(\mathrm{PSL}_N,V)\;=\;3,\;16,\;469,\;36672
\qquad (N=6,7,8,9)
\]
in sufficiently large defining characteristic, and $3{,}537{,}142$ in
type $E_6$ \cite[\S4.1 and Theorem~4.7]{Luebeck}.
%
  In light of these examples Guralnick, Hodge, Parshall, and  Scott concluded their AIM workshop saying that Guralnick's conjecture {\em `could no longer be believed'} \cite{GHPS}.
The theory has advanced since, with Guralnick and Tiep's bounds in terms
of sectional rank \cite{GTSectional,GTSectionalII} and the
characteristic-free bounds of \cite{BBDNPPW,BNPPSS15}; but the question
of a uniform bound has remained open for a further fourteen years. Our
first main result settles it, thus proving that Guralnick's conjecture is false. 
 
\begin{maintheorem}
\label{thm:A}
\label{thm:guralnick}
\label{thm:guralnick-counterexample}
\label{cor:unbounded-first-cohomology}
For every integer $r\ge2$ and every sufficiently large prime $p$,
the finite
simple group $\mathrm{PSL}_{2r+1}(\mathbb F_p)$
has a faithful absolutely
irreducible module $V_{r,p}$ over $\mathbb F_p$ such that
\begin{equation}\label{eq:unbounded-first-cohomology}
\dim_{\mathbb F_p}
H^1\big(\mathrm{PSL}_{2r+1}(\mathbb F_p),V_{r,p}\big)\;=\;r.
\end{equation}
%
The module $V_{r,p}$ is obtained by 
first restricting    the simple  
$\mathrm{SL}_{2r+1}(\overline{\mathbb F}_p)$-module labelled by the weight
\[
\lambda_{r,p}\;=\;
2\varpi_{r-1}+(p-2)(\varpi_r+\varpi_{r+1})+2\varpi_{r+2}
\]
to the finite subgroup $\mathrm{SL}_{2r+1}(\mathbb F_p)$,
 and then descending to the projective quotient. 
 % (Here the $\varpi_i$ are the fundamental weights.) 
 % Here the $\varpi_i$ are the fundamental weights; the details and the proof are given in  \Cref{sec:cohomology}.
 \end{maintheorem}
 
The computed examples reached dimension in the millions, but they did not lead to a  
uniform family that   could be proven unbounded. 
Theorem~\ref{thm:A} supplies such a family with
an explicit highest weight and an exact dimension formula.
 %
Theorem~\ref{thm:A} goes beyond unboundedness:
 it implies that every nonnegative integer
occurs as the dimension of first cohomology of a finite simple
group with faithful absolutely irreducible coefficients.
This solves an inverse problem in the spirit of constructing
varieties with prescribed Hodge numbers
\cite[Theorem~1]{Schreieder2015}.



 
\subsection{}
Our approach to Theorem~\ref{thm:A} is through Kazhdan--Lusztig combinatorics.
Let $W$ be the affine Weyl group of $\mathrm{SL}_N$, acting on weights by the
dot action $x\cdot\lambda$ with translations scaled by $p$, and let $w_0$ be
the longest element of its finite Weyl subgroup. We write
$$h_{y,w}(v)=v^{\ell(w)-\ell(y)}P_{y,w}(v^{-2})$$ for the Kazhdan--Lusztig
polynomials in Soergel's normalization, and $\mu(y,w)$ for the coefficient of
$v$ in $h_{y,w}$. For $p$ large and $x\cdot0$ $p$-restricted,
\begin{equation}
  \mu(w_0,w_0x)=\dim H^1(\mathrm{SL}_N(\mathbb F_p),L(x\cdot0))  ,
\end{equation}
by results of
Andersen and of Bendel, Nakano and Pillen (see Section~\ref{sec:cohomology}). As $H^1\ne0$ forces the centre to act trivially, $L(x\cdot0)$ is not
faithful for $\mathrm{SL}_N(\mathbb F_p)$; it descends to a faithful module
for the simple group $\mathrm{PSL}_N(\mathbb F_p)$ with the same $H^1$, so showing that these coefficients are unbounded would disprove Guralnick's conjecture.
This combinatorial formulation is also what allowed L\"ubeck to attack the
conjecture by computer: his values are coefficients $\mu(w_0,w_0x)$, computed by machine for $N\le9$. The following is the technical heart of the paper.
 


 

 \begin{maintheorem}
\label{thm:B}
\label{thm:mu}
\label{thm:main}
Let $r\ge2$, put $N=2r+1$, and let $W$ be the affine Weyl group of type
$\widetilde A_{2r}$, realized by affine permutations of period $N$.
For the affine permutation with window
\[
\beta_r=
[\,N+1,\dots,N+r-1,\;1-r,\;r+1,\;N+r,\;2-r,\dots,0\,]
\]
one has
\[
\mu(w_0,w_0\beta_r)\;=\;r.
\]
More generally, the elements $y_0,\dots,y_{r-1}$ of
Section~\ref{sec:second-chain}, with $y_0=\beta_r$, satisfy
\[
h_{w_0,w_0y_s}(v)
\;=\;r\,v^{s+1}+O(v^{s+3})
\qquad(0\le s\le r-1).
\]
\end{maintheorem}


\begin{remark}
Scott and Xi constructed unbounded $\mu(y,w)$ with both endpoints varying, so that they govern $\mathrm{Ext}^1$ between two simple modules rather than $\mathrm{H}^1$.
They asked whether the coefficients can still
be unbounded when the lower endpoint is required to be $w_0$ (as in
the formula for first cohomology above) and proposed explicit candidates
\cite[Section~4]{ScottXi}.
\Cref{thm:B} gives an affirmative answer using a different family of upper endpoints,
namely the elements $w_0\beta_r$. 
\end{remark}









The higher-degree coefficients in Theorem~\ref{thm:B} have consequences
for the cohomology of finite simple groups in every degree. Taking
$s=n-1$ in that theorem gives a coefficient of $v^n$ equal to $r$ whenever $1\leq n\leq r$.
This is the dimension of the corresponding rational cohomology group
in sufficiently large characteristic. The comparison with finite groups
in Section~\ref{subsec:higher-cohomology} then gives the following result.

 
 
\begin{maintheorem}
\label{thm:C}
\label{thm:higher}
Let $n\ge1$ and $r\ge\max\{2,n\}$.
For every sufficiently large prime $p$ and every sufficiently large
integer $a$, the finite simple group
$\mathrm{PSL}_{2r+1}(\mathbb F_{p^a})$ has a faithful absolutely
irreducible module $V$ over $\overline{\mathbb F}_p$ such that
\[
\dim_{\overline{\mathbb F}_p}
H^n\big(\mathrm{PSL}_{2r+1}(\mathbb F_{p^a}),V\big)\;\ge\;r.
\]
\end{maintheorem}

 
 
\subsection{}  
 The natural home for our results is the quantum group setting. Let $U_\zeta$ be Lusztig's quantum group of $\mathfrak{sl}_N(\mathbb C)$ at a primitive root of unity $\zeta$ of odd order $\ell>N$, and let $\cdot_\ell$ denote the dot action of the affine Weyl group with translations scaled by $\ell$. Here Lusztig's character formula holds \cite{Tanisaki2004}, and its cohomological interpretation \cite{CPS-abstractKL,PS-BoundingExt} expresses $h_{w_0,w_0x}$ as the generating function of the dimensions of $H^j(U_\zeta,L_\zeta(x\cdot_\ell0))$. The corresponding statement for algebraic groups in characteristic $p$ is deduced from the quantum case: Lusztig's character formula was proved for large $p$ by comparison with the quantum group \cite{AJS}. Our results for finite groups (Theorems~\ref{thm:guralnick} and~\ref{thm:C}) therefore pass through the quantum group, and Theorem~\ref{thm:D} below states the quantum result directly.  


 
Parshall and Scott bounded the dimensions of cohomology groups for quantum groups in terms of the degree and the root system \cite[Corollary~6.7]{PS-BoundingExt}, and asked whether the dependence on the root system could be removed \cite[Remark~6.6]{PS-BoundingExt}. Theorem~\ref{thm:D} shows that it cannot, even in a fixed positive degree.





\begin{maintheorem}
\label{thm:D}
\label{thm:quantum}
Let $r\ge2$, put $N=2r+1$, and let $\ell>N$ be odd.
Let $U_\zeta$ be Lusztig's quantum group of $\mathfrak{sl}_N(\mathbb C)$
at a primitive $\ell$th root of unity $\zeta$.
For $0\le s\le r-1$, put $\lambda_{r,s}=y_s\cdot_\ell0$.
Then the irreducible modules $L_\zeta(\lambda_{r,s})$ satisfy
\[
\dim_{\mathbb C}
H^{s+1}\big(U_\zeta,L_\zeta(\lambda_{r,s})\big)\;=\;r.
\]
Moreover, $H^j\big(U_\zeta,L_\zeta(\lambda_{r,s})\big)=0$ for
$0\le j<s+1$, so the prescribed dimension occurs in the first
nonzero cohomological degree.
\end{maintheorem}

% The main theorems in this paper give further instances of the unboundedness
% illustrated by the opening examples. For the smooth projective curves
% of genus $g$ discussed there, the cohomology dimensions reflect their
% geometry. In the complex case, for instance, the underlying surface
% has $g$ handles. Each contributes two independent classes to first
% singular cohomology with rational coefficients, so its dimension is $2g$.
% It would be interesting to trace the unbounded growth of cohomology
% dimensions in the families constructed here to structural features of
% their groups and modules.


\subsection{} 
Our group-theoretic constructions require  the characteristic to be sufficiently large
relative to the rank, and therefore leave  open the possibility that
Guralnick's conjecture holds in a given fixed characteristic $p>0$. 
For the same reason, the bounds sought in
Question~\ref{question:higher-degree-bound} may still hold in every
degree, with constants depending on the degree and the characteristic.
This leads us to formulate the following fixed-characteristic version of
Question~\ref{question:higher-degree-bound}.
 
% \begin{question}\label{question:fixed-characteristic-first}
% Fix a prime $p$. Does there exist a constant $C_p$ such that
% \[
% \dim H^1(G,V)\leq C_p
% \]
% for every finite group $G$ and every faithful absolutely irreducible
% module $V$ in characteristic~$p$?
% \end{question}

% We also ask the fixed-characteristic version of
% Question~\ref{question:higher-degree-bound}.

\begin{question}\label{question:fixed-characteristic-higher}
Fix a prime $p$ and an integer $n\geq 1$. Does there exist a constant
$C_{n,p}$ such that
\[
\dim H^n(G,V)\leq C_{n,p}
\]
for every finite \textbf{simple} group $G$ and every absolutely irreducible
module $V$ in characteristic $p$?
\end{question}
 
Keeping the characteristic fixed while the rank grows takes us
beyond the regime where ordinary Kazhdan--Lusztig polynomials control
cohomology. The modular phenomena reflected in the
$p$-canonical basis, which is substantially harder to compute
than the ordinary Kazhdan--Lusztig basis, suggest that extending
our approach to these questions would require considerably
deeper control of modular representation theory.
 
% Thus Theorem~\ref{thm:B} implies Theorem~\ref{thm:D} directly, while
% Theorems~\ref{thm:guralnick} and~\ref{thm:C} require two further steps:
% the dictionary in characteristic $p$, valid for large $p$ by reduction to
% the quantum case, and the comparison between algebraic and finite groups
% \cite{BNPcompare,CPSvdK1977,CPS1983}. The second step is where exactness
% can be lost: in degree one the comparison of \cite{BNPcompare} is an
% isomorphism, while in higher degrees we only use injectivity of
% restriction, producing the inequality in Theorem~\ref{thm:C}.
 


\subsection*{AI statement}
This project began while CB, NL, and DP were participating in the 2025 ICERM Semester Program on Categorification and Computation in Algebraic Combinatorics. 
  Our starting point was the observation that, within the parallelepiped of alcoves corresponding to p-restricted weights, Lübeck’s examples lie in the regions farthest from the identity. 
  Direct computations using Sage within this region were impractical in higher ranks, and the available values showed no recognizable pattern to guide a proof. 
  At this point DF joined the project, using MAGMA calculations to conduct a more detailed investigation of this region.
  This led us to ask whether more tractable examples might exist closer to the identity.

We turned to GPT-6 Astra and asked it to explore this region. From the data it provided, we identified the elements $\beta_r$ as a promising family exhibiting a uniform growth pattern. 
 However, the obstacle remained: although these elements are closer to the identity, and we could compute individual cases in small ranks, the usual Kazhdan--Lusztig algorithm gave no way to directly calculate the entire infinite family.



We therefore asked GPT-6 Astra to find a nearby auxiliary element
$X_0$ whose Kazhdan--Lusztig basis element could be computed more
directly, and from which $\beta_r$ could be reached using only simple
reflections of the finite Weyl group. 
 The strategy of the proof was ours: we asked for a computable element in a prescribed region, together with a factorization of $\beta_r$ through simple reflections along which the Kazhdan–Lusztig coefficients remain manageable. Following our instructions, the model produced a proof of Theorem B.
 All arguments were checked, rewritten, and completed by the authors, who take full responsibility for the results.


\smallskip
\noindent{\bf Structure of the paper.}
Our strategy is to estimate the relevant Kazhdan--Lusztig polynomials at $X_0$ and then propagate these estimates along two successive chains leading to $\beta_r$, using multiplication formulas in the spherical Hecke module.
In \Cref{sec:argument}, we introduce the notation and preliminaries, describe these chains, and deduce
Theorems~\ref{thm:A}, \ref{thm:B}, \ref{thm:C} and~\ref{thm:D} from the coefficient estimates and multiplication formulas proved in the subsequent sections.
\Cref{sec:parabolic} establishes the initial estimates at $X_0$, and \Cref{sec:corrections} develops the tools common to both chains.
Finally, \Cref{sec:first-multiplication-proof,sec:second-multiplication-proof} establish the multiplication formulas and coefficient estimates along the first and second chains, respectively.

\smallskip
\noindent{\bf Acknowledgments.}
This work began during the \emph{Categorification and Computation in Algebraic Combinatorics} programme
at the Institute for Computational and Experimental Research in Mathematics (ICERM). We thank ICERM for its hospitality and for a productive environment. This  work was supported by the National Science Foundation under Grant No.\ DMS-2424556 while CB, NL and DP were in residence at ICERM.
DF acknowledges support from a LACIM postdoctoral fellowship at the Université du Québec à Montréal.
































































































































































































































































































































































































































































































































































































































































































































































































































































\section{The main theorems}
\label{sec:argument}\label{sec:conventions}\label{sec:chains}

In this section we introduce the affine Weyl group and its Hecke
algebra and state the identities and coefficient estimates needed
for the Kazhdan--Lusztig calculation. Taking these results as given,
we prove \Cref{thm:B} and use it to deduce the cohomological results
in \Cref{thm:A,thm:C,thm:D}. The proofs of the identities and
estimates are deferred to the subsequent sections.

\subsection{The affine Weyl group and its Hecke algebra}
Fix $r\geq1$, put $N=2r+1$, and let $W$ be the affine Weyl group
of type $\widetilde A_{2r}$ with simple reflections $S=\{s_0,\ldots,s_{N-1}\}$ and Coxeter graph
\begin{center}
\begin{tikzpicture}[scale=0.8]
    \node[draw,circle,inner sep=2pt,label=above:{$s_0$}] (s0) at (0,2) {};
    \node[draw,circle,inner sep=2pt,label=below:{$s_1$}] (s1) at (-3,0) {};
    \node[draw,circle,inner sep=2pt,label=below:{$s_2$}] (s2) at (-1.7,0) {};
    \node[draw,circle,inner sep=2pt,label=below:{$s_{N-2}$}] (sn2) at (1.7,0) {};
    \node[draw,circle,inner sep=2pt,label=below:{$s_{N-1}$}] (sn1) at (3,0) {};

    \draw (s0) -- (s1);
    \draw (s0) -- (sn1);
    \draw (s1) -- (s2);
    \draw (sn2) -- (sn1);
    \draw (s2) -- (-0.5,0);
    \draw (0.5,0) -- (sn2);
    \node at (0,0) {$\cdots$};
\end{tikzpicture}
\end{center}
We write $\ell$ for the Coxeter length and $\leq$ for the Bruhat order. Let $I=\{s_1,\ldots,s_{N-1}\}$ and write $w_I$ for the longest
element of the finite parabolic subgroup $W_I$ of type $A_{2r}$ (this is called $w_0$ in the Introduction).
%
We use the standard realization of this Coxeter system by affine
permutations \cite[Section~8.3]{BB}. The abstract simple reflections
correspond to the periodic adjacent transpositions described below.
% 
% \begin{definition}\label{def:windows}
Under this realization, the elements of $W$ are the bijections
$w:\mathbb Z\to\mathbb Z$ satisfying
\[
 w(i+N)=w(i)+N\quad(i\in\mathbb Z),\qquad
 \sum_{i=1}^Nw(i)=\frac{N(N+1)}2.
\]
Such an affine permutation is determined by its \emph{window}
$[w(1),\ldots,w(N)]$. Products are compositions:
$(xy)(i)=x(y(i))$. The simple reflections have windows
\[
\begin{aligned}
 s_i&=[1,\ldots,i-1,i+1,i,i+2,\ldots,N] &&(1\leq i<N),\\
 s_0&=[0,2,\ldots,N-1,N+1].
\end{aligned}
\]
Thus right multiplication by $s_i$, $1\leq i<N$, exchanges positions
$i$ and $i+1$, while right multiplication by $s_0$ replaces the first and
last entries by $w(N)-N$ and $w(1)+N$, respectively.
Empty intervals in windows are omitted.
The subgroup $W_I$ is the ordinary symmetric group on
$\{1,\ldots,N\}$. We write $ {}^I W$ for the minimal representatives
of $W_I\backslash W$; they are characterized by
\begin{equation}\label{eq: char representatives}
   w^{-1}(1)<\cdots<w^{-1}(N).  
\end{equation}
\color{black}More generally, for a proper subset $P\subsetneq S$ we write $W_P$ for the
(finite) subgroup it generates and $w_P$ for its longest element. For
$Q\subseteq P\subsetneq S$, we write ${}^QW_P$ for the set of minimal coset
representatives of  $W_Q\backslash W_P$. Every $u\in W_P$ factors
uniquely as $u=ab$ with $a\in W_Q$, $b\in{}^QW_P$ and
$\ell(u)=\ell(a)+\ell(b)$ \cite[Proposition~2.4.4]{BB}.
 


Let $\mathcal H=\mathcal H(W)$ be the Hecke algebra over
$\mathbb Z[v,v^{-1}]$ with standard basis $\{H_w\}_{w\in W}$ and relations
\[
 H_s^2=1+(v^{-1}-v)H_s,\qquad
 H_xH_y=H_{xy}\quad\text{if }\ell(xy)=\ell(x)+\ell(y).
\]
Its bar involution sends $v$ to $v^{-1}$ and $H_w$ to
$H_{w^{-1}}^{-1}$.
 For every $w\in W$, we let $C_w\in\mathcal H$ denote the  unique element  
satisfying
\[
 \overline{C_w}=C_w,\qquad
 C_w\in H_w+\sum_{y<w}v\mathbb Z[v]H_y.
\]
The elements $\{C_w\}_{w\in W}$ form the canonical basis of
Kazhdan and Lusztig \cite{KL}.
The \emph{Kazhdan--Lusztig polynomials} $h_{y,w}$ in this normalization are
the change of basis coefficients defined by
\begin{equation}\label{eq:KL-coefficients}
 C_w=\sum_{y\leq w}h_{y,w}H_y.
\end{equation}
Thus $h_{w,w}=1$ and $h_{y,w}\in v\mathbb Z[v]$ for $y<w$;
we put $h_{y,w}=0$ when $y\not\leq w$.
It is easy to check that $C_s=H_s+v$.
The canonical basis element of $w_I$ is the symmetrizer \cite[Proposition~2.9]{Soergel}
\[
 C_{w_I}=\sum_{u\in W_I}v^{\ell(w_I)-\ell(u)}H_u.
\]
Pairing the terms indexed by $u$ and $us$, and using the quadratic
Hecke relation, gives
\begin{equation}\label{eq:finite-symmetrizer-eigenrelation}
 C_{w_I}H_s=v^{-1}C_{w_I}\qquad(s\in I).
\end{equation}

% To compute the coefficients indexed by $w_I$, we work in the spherical module.
% Lemma~\ref{lem:mu} will identify the polynomials $h_{w_I,w_Ix}$ with a single coordinate in its standard basis.
% 
% \begin{definition}\label{def:spherical}\label{def:spherical-module}
The \emph{spherical module} is the right ideal $\M=C_{w_I}\mathcal H$. 
Its standard basis is $\{M_x \mid x\in {}^IW\}$ where
\[
 M_x=C_{w_I}H_x.
\]
The elements
\[
 \C_x=C_{w_Ix}\qquad(x\in{}^I W)
\]
lie in $\M$ and form its Kazhdan--Lusztig basis
\cite[Theorem~3.1 and Proposition~3.4]{Soergel}. In particular, both
families are $\mathbb Z[v,v^{-1}]$-bases of $\M$, and
\[
 \C_x\in M_x+\sum_{y<x}v\mathbb Z[v]M_y.
\]
% \end{definition}

For an element $X$ of a free module with a distinguished basis, we write $[Y]X$ for the coefficient of a basis element $Y$ in $X$.
We define
\[
 \eps(F)=[M_e]F \quad(F\in\M),\qquad h_x=\eps(\C_x) \quad (x\in{}^I W),
\]
and
\begin{equation}\label{eq:mu-definition}
 \mu(y,w)=[v]h_{y,w}\quad(y,w\in W),\qquad
 \mus(c,x)=\mu(w_Ic,w_Ix)\quad(c,x\in{}^I W).
\end{equation}

{\color{black}
\begin{lemma}\label{lem:mu}
For $c,x\in{}^I W$, we have
\[
[M_c]\C_x=h_{w_Ic,w_Ix}.
\]
In particular, $h_x=h_{w_I,w_Ix}$ and
$\mus(e,x)=[v]h_x$.
Moreover, $h_e=1$, while $h_x\in v\mathbb Z[v]$ for $x\ne e$,
and all exponents in $h_x$ have parity $\ell(x)$.
\end{lemma}
}
\begin{proof}\color{black}
This is standard; see    \cite[Theorem 3.1 \& Proposition 3.4]{Soergel} for the spherical module identification and  \cite{KL}  for the parity of polynomial exponents. \end{proof}


The chains considered in \Cref{sec:first-chain} below use only right multiplication by
generators in $I$.
The following rule will turn their Hecke algebra multiplications into scalar recurrences.
Put $\eta=v+v^{-1}$.

\begin{lemma}\label{lem:scalar}
For $s\in I$ and $F\in\M$,
\begin{equation}\label{eq:scalar}
 \eps(FH_s)=v^{-1}\eps(F),\qquad
 \eps(FC_s)=(v+v^{-1})\eps(F).
\end{equation}
% \begin{equation}\label{eq:scalar}
%  \eps(FH_s)=v^{-1}\eps(F),\qquad
%  \eps(FC_s)=\eta\eps(F),\qquad \eta=v+v^{-1}.
% \end{equation}
\end{lemma}
\begin{proof}
Right multiplication by $H_s$ preserves the span of the standard
basis elements indexed by each $(I,I)$ double coset. If
$c\in{}^I W$ and $c\ne e$, then $W_IcW_I\ne W_I$, so $M_cH_s$
cannot contribute to $M_e$. The remaining basis element satisfies
$M_eH_s=v^{-1}M_e$ by
\eqref{eq:finite-symmetrizer-eigenrelation}. This proves the first
identity. 
The second follows from $C_s=H_s+v$.
\end{proof}

By Lemma~\ref{lem:mu}, Theorem~\ref{thm:mu} amounts to computing
$[v]h_{\beta_r}$. We do this along the two chains defined below.

\subsection{The first chain}\label{sec:first-chain}
To compute $[v]h_{\beta_r}$, we introduce an auxiliary element $X_0$
and connect it to $\beta_r$ by two chains of finite simple reflections.
We first obtain coefficient estimates along the first chain, then use
the second chain to determine the desired coefficient.

\begin{definition}\label{def:endpoints}
We define $X_0$ and, for completeness, recall the definition of $\beta_r$
from the introduction:
\begin{align}
 X_0&=[1-r,r+1,N+1,\ldots,N+r,2-r,\ldots,0],\notag\\
 \beta_r&=[N+1,\ldots,N+r-1,1-r,r+1,N+r,2-r,\ldots,0].
 \label{eq:endpoints}
\end{align}
\end{definition}

For $r=5$, \Cref{figure1,figure2} show the affine string diagrams of $\beta_r$ and $X_0$, respectively.


 


\begin{lemma}\label{lem:endpoints}
Both elements in Definition~\ref{def:endpoints} lie in $ {}^I W$, and
\[
 \beta_r=X_0(s_2\cdots s_r)(s_1\cdots s_{r-1}).
\]
Every prefix of $\beta_r$ lies in $ {}^I W$ and each simple reflection increases the length.
Consequently, $\ell(\beta_r)=\ell(X_0)+2r-2$.
\end{lemma}
\begin{proof}
Each  window in \eqref{eq:endpoints} contains precisely the values
\[
 1-r,\ldots,0,\qquad r+1,\qquad N+1,\ldots,N+r.
\]
The residues $1,\ldots,r$ have been shifted by $N$, those in
$r+2,\ldots,N$ by $-N$, and $r+1$ is unchanged. Hence the residues
are distinct and the window sum is $N(N+1)/2$.
It follows that both $X_0$ and $\beta_r$ are elements of $W$. 

The upper entries $N+1,\ldots,N+r$ occur in increasing positions,
as do the lower entries $1-r,\ldots,0$. The inverse images of the
residues $1,\ldots,r$ therefore increase in $[1-N,0]$, the inverse
image of $r+1$ lies in $[1,N]$, and those of $r+2,\ldots,N$ increase
in $[N+1,2N]$. Thus both inverse windows are increasing.
By \eqref{eq: char representatives} we conclude that $X_0$ and $\beta_r$ belong to ${}^IW$. 

Right multiplication by $s_2\cdots s_r$ moves $r+1$ past
$N+1,\ldots,N+r-1$; multiplication by $s_1\cdots s_{r-1}$ then
moves $1-r$ past the same entries. This gives the window of $\beta_r$.
Each swap exchanges an increasing adjacent numerical pair, so it
increases Coxeter length. %\David{This argument requires some length formula. We have some at the beginning of 3.1 but it is to ugly to move here. }
 It also preserves the position order of
the upper and lower entries, so every intermediate element remains
in $ {}^I W$. For $r=1$, both products are empty and
$X_0=\beta_1=s_0$.
\end{proof}

\begin{figure}[!htbp]
  \centering
  \begin{tikzpicture}[
    x=0.55cm,          % Compact horizontal spacing
    y=1.1cm,           % Height proportion
    >=Stealth,
    thick,
 scale=0.85 ]

    \clip (-5.8, -2.1) rectangle (17.8, 5.8);

    \foreach \i in {-20,...,30} {
      \fill[gray!35] (\i, 0) circle (1.5pt);
      \fill[gray!35] (\i, 5) circle (1.5pt);
    }

    \foreach \k in {-2, -1, 0, 1, 2} {
      \foreach \i/\w in {
        1/12, 2/13, 3/14, 4/15, 5/-4,
        6/6, 7/16, 8/-3, 9/-2, 10/-1, 11/0
      } {
        \pgfmathsetmacro{\src}{\i + \k*11}
        \pgfmathsetmacro{\tgt}{\w + \k*11}
       
        \ifnum \k=0 \else
          \draw[gray!45, dashed, ->, looseness=0.8] (\src, 0) to[out=60, in=240] (\tgt, 5);
        \fi
      }
    }

    \foreach \i in {-5,...,0} {
      \node[below=3pt, font=\tiny, color=gray!80] at (\i, 0) {$\i$};
    }
    \foreach \i in {12,...,17} {
      \node[below=3pt, font=\tiny, color=gray!80] at (\i, 0) {$\i$};
    }
    \foreach \j in {-5,...,-1, 1, 2, 3, 4, 5, 7, 8, 9, 10, 11, 17} {
      \node[above=3pt, font=\tiny, color=gray!80] at (\j, 5) {$\j$};
    }

    \foreach \i/\w in {
      1/12, 2/13, 3/14, 4/15, 5/-4,
      6/6, 7/16, 8/-3, 9/-2, 10/-1, 11/0
    } {
      \draw[black, line width=1.4pt, ->, looseness=0.8] (\i, 0) to[out=60, in=240] (\w, 5);
     
      \node[below=3pt, font=\tiny\bfseries, color=black] at (\i, 0) {$\i$};
     
      \node[above=3pt, font=\tiny\bfseries, color=black] at (\w, 5) {$\w$};
    }

    \draw[decorate, decoration={brace, mirror, raise=14pt, amplitude=6pt}, thick, black]
      (0.6, 0) -- (11.4, 0)
      node[midway, yshift=-26pt, font=\small, color=black]
      {$\beta_5 = [12, 13, 14, 15, -4, 6, 16, -3, -2, -1, 0]$};

  \end{tikzpicture}
 
  
  
  \vspace{-0.6cm}\caption{The affine string diagram of $\beta_5$ in type
  $\widetilde A_{10}$. Solid black strands show its window  and dashed
  gray strands indicate periodic translates.   } 
  \label{figure1}
\end{figure}





\vspace{-0.5cm}
\begin{figure}[htbp]
   \centering
  \begin{tikzpicture}[
    x=0.55cm,          % Compact horizontal spacing
    y=1.1cm,           % Height proportion
    >=Stealth,
    thick,  scale=0.85
  ]

    % --- Clip everything to visible range ---
    \clip (-5.8, -2.1) rectangle (17.8, 6.75);

    % --- Background guide dots ---
    \foreach \i in {-20,...,30} {
      \fill[gray!35] (\i, 0) circle (1.5pt);
      \fill[gray!35] (\i, 5) circle (1.5pt);
    }

    % --- Complete Shifted Gray Lines (All Affine Shifts) ---
    % Window mapping for X_0:
    % 1->-4, 2->6, 3->12, 4->13, 5->14, 6->15, 7->16, 8->-3, 9->-2, 10->-1, 11->0
    \foreach \k in {-2, -1, 0, 1, 2} {
      \foreach \i/\w in {
        1/-4, 2/6, 3/12, 4/13, 5/14, 
        6/15, 7/16, 8/-3, 9/-2, 10/-1, 11/0
      } {
        \pgfmathsetmacro{\src}{\i + \k*11}
        \pgfmathsetmacro{\tgt}{\w + \k*11}
        
        % Draw gray strands for non-core windows (k != 0)
        \ifnum \k=0 \else
          \draw[gray!45, densely dotted, ->, looseness=0.8] (\src, 0) to[out=60, in=240] (\tgt, 5);
        \fi
      }
    }

    % --- Top and Bottom Axis Labels (Gray range) ---
    \foreach \i in {-5,...,0} {
      \node[below=3pt, font=\tiny, color=gray!80] at (\i, 0) {$\i$};
    }
    \foreach \i in {12,...,17} {
      \node[below=3pt, font=\tiny, color=gray!80] at (\i, 0) {$\i$};
    }
    \foreach \j in {-5,...,-5, 1, 2, 3, 4, 5, 7, 8, 9, 10, 11, 17} {
      \node[above=3pt, font=\tiny, color=gray!80] at (\j, 5) {$\j$};
    }

    % --- Core Window Strands (Solid black) ---
    \foreach \i/\w in {
      1/-4, 2/6, 3/12, 4/13, 5/14, 
      6/15, 7/16, 8/-3, 9/-2, 10/-1, 11/0
    } {
      \draw[black, line width=1.4pt, ->, looseness=0.8] (\i, 0) to[out=60, in=240] (\w, 5);
      
      % Bottom Input Node & Label
      \fill[black] (\i, 0) circle (2.5pt);
      \node[below=3pt, font=\tiny\bfseries, color=black] at (\i, 0) {$\i$};
      
      % Top Output Label & Node
      \fill[black] (\w, 5) circle (2.5pt);
      \node[above=3pt, font=\tiny\bfseries, color=black] at (\w, 5) {$\w$};
    }

    % --- Bottom Bracket ---
    \draw[decorate, decoration={brace, mirror, raise=14pt, amplitude=6pt}, thick, black] 
      (0.6, 0) -- (11.4, 0) 
      node[midway, yshift=-28pt, font=\small, color=black] 
      {$X_0 = [-4, 6, 12, 13, 14, 15, 16, -3, -2, -1, 0]$};

    % --- Overbraces: the value sets A_5, B_5, C_5 (lower, middle, upper) ---
    \begin{scope}[red]
    \draw[decorate, decoration={brace, raise=12pt, amplitude=5pt}, thick]
      (-4.4, 5) -- (0.4, 5) node[midway, yshift=24pt, font=\small] {$\A$};
    \draw[decorate, decoration={brace, raise=12pt, amplitude=4pt}, thick]
      (5.6, 5) -- (6.4, 5) node[midway, yshift=23pt, font=\small] {$\B$};
    \draw[decorate, decoration={brace, raise=12pt, amplitude=5pt}, thick]
      (11.6, 5) -- (16.4, 5) node[midway, yshift=24pt, font=\small] {$\CC$};
    \end{scope}

  \end{tikzpicture}

  \vspace{-0.8cm}
  \caption{The affine string diagram for $X_0$, with
  $N=11$.
  Compare $X_0(s_2s_3s_4s_5)(s_1s_2s_3s_4)$ with the diagram in  \Cref{figure1} (see also  \Cref{lem:endpoints}). 
  To see that this is an element of $^I W$ simply note that the strands 
  that connect to the top vertices $1,\dots ,11$ are non-crossing. 
  The labels $\sf A, B, C$ are used in \Cref{sec:parabolic}.
  }
  \label{figure2}
\end{figure}





\begin{definition}\label{def:chains}
For $0\leq k\leq r-1$, put
 $   x_k=X_0(s_2\cdots s_{k+1})$.  
Equivalently, $x_k$ can be expressed in terms of its window as follows 
\begin{equation}\label{eq:chains}
 x_k=[1-r,N+1,\ldots,N+k,r+1,
      N+k+1,\ldots,N+r,2-r,\ldots,0].
\end{equation}
Empty products are the identity. Hence $x_0=X_0$ and
$x_{k+1}=x_ks_{k+2}$ for $0\leq k\leq r-2$.
The factorizations are length additive by Lemma~\ref{lem:endpoints}.
\end{definition}

% \begin{definition}\label{def:chains}
% For $0\leq k\leq r-1$, put
% \begin{equation}\label{eq:first-chain-product}
%  x_k=X_0(s_2\cdots s_{k+1}),\qquad h_{x_{k}}.
% \end{equation}
% Equivalently, $x_k$ can be expressed in terms of its window as follows 
% \begin{equation}\label{eq:chains}
%  x_k=[1-r,N+1,\ldots,N+k,r+1,
%       N+k+1,\ldots,N+r,2-r,\ldots,0].
% \end{equation}
% Empty products are the identity. Hence $x_0=X_0$ and
% $x_{k+1}=x_ks_{k+2}$ for $0\leq k\leq r-2$.
% The factorizations are length additive by Lemma~\ref{lem:endpoints}.
% \end{definition}

Our aim is to estimate $h_{x_{r-1}}$. The starting point is
$h_{X_0}$, whose initial coefficients are given by the following
proposition. Its proof is deferred to Subsection~\ref{sec:seedproof}.
For the estimates below, we write $f=O(v^d)$ when
$f\in v^d\mathbb Z[[v]]$.
\begin{proposition}\label{prop:seed}
Let $\Part(d)$ be the number of partitions of $d$, with
$\Part(0)=1$. For $r\geq2$,
\begin{equation}\label{eq:seed}
 h_{X_0}=\sum_{d=0}^{r-1}\Part(d)v^{2d+1}+O(v^{2r+1}).
\end{equation}
\end{proposition}
\begin{proof}
     See subsection~\ref{sec:seedproof}.
\end{proof}
The coefficient estimate along this chain will be proved together
with the following multiplication formulas in \Cref{sec:first-multiplication-proof}. Their correction terms are indexed by
the following family of affine permutations.

\begin{definition}\label{def:gamma}
For $1\leq j\leq r$  we let 
\[
 \gamma_j=[j+1,N+1,\ldots,N+j,j+2,\ldots,N-j,1-j,\ldots,0].
\]
 \end{definition}
As in the proof of Lemma~\ref{lem:endpoints}, the inverse windows of the 
$\gamma_j$ 
are increasing, so $\gamma_j\in{}^I W$. An example is depicted in \cref{gamma3}.

\begin{figure}[htbp]
  \centering
  % ---- choose j here (N=11, so 1 <= j <= 5) and paste the matching window list ----
  %  j=1: 1/2,2/12,3/3,4/4,5/5,6/6,7/7,8/8,9/9,10/10,11/0
  %  j=2: 1/3,2/12,3/13,4/4,5/5,6/6,7/7,8/8,9/9,10/-1,11/0
  %  j=3: 1/4,2/12,3/13,4/14,5/5,6/6,7/7,8/8,9/-2,10/-1,11/0
  %  j=4: 1/5,2/12,3/13,4/14,5/15,6/6,7/7,8/-3,9/-2,10/-1,11/0
  %  j=5: 1/6,2/12,3/13,4/14,5/15,6/16,7/-4,8/-3,9/-2,10/-1,11/0
  \def\Wlist{1/4,2/12,3/13,4/14,5/5,6/6,7/7,8/8,9/-2,10/-1,11/0}%  gamma_3
  \begin{tikzpicture}[
    x=0.55cm,          % Compact horizontal spacing
    y=1.1cm,           % Height proportion
    >=Stealth,
    thick,  scale=0.85
  ]

    % --- Clip everything to visible range ---
    \clip (-5.8, -2.1) rectangle (17.8, 5.8);

    % --- Background guide dots ---
    \foreach \i in {-20,...,30} {
      \fill[gray!35] (\i, 0) circle (1.5pt);
      \fill[gray!35] (\i, 5) circle (1.5pt);
    }

    % --- Complete Shifted Gray Lines (All Affine Shifts) ---
    \foreach \k in {-2, -1, 0, 1, 2} {
      \foreach \i/\w in \Wlist {
        \pgfmathsetmacro{\src}{\i + \k*11}
        \pgfmathsetmacro{\tgt}{\w + \k*11}
        \ifnum \k=0 \else
          \draw[gray!45, densely dotted, ->, looseness=0.8] (\src, 0) to[out=60, in=240] (\tgt, 5);
        \fi
      }
    }

    % --- Bottom axis labels (gray): the values outside 1..11 ---
    \foreach \i in {-5,...,0} {
      \node[below=3pt, font=\tiny, color=gray!80] at (\i, 0) {$\i$};
    }
    \foreach \i in {12,...,17} {
      \node[below=3pt, font=\tiny, color=gray!80] at (\i, 0) {$\i$};
    }

    % --- Top axis labels (gray): every visible value that is NOT an output of the window ---
    \foreach \v in {-5,...,17} {
      \def\found{0}
      \foreach \i/\w in \Wlist { \ifnum\v=\w \xdef\found{1}\fi }
      \ifnum\found=0
        \node[above=3pt, font=\tiny, color=gray!80] at (\v, 5) {$\v$};
      \fi
    }

    % --- Core Window Strands (Solid black) ---
    \foreach \i/\w in \Wlist {
      \draw[black, line width=1.4pt, ->, looseness=0.8] (\i, 0) to[out=60, in=240] (\w, 5);
      \fill[black] (\i, 0) circle (2.5pt);
      \node[below=3pt, font=\tiny\bfseries, color=black] at (\i, 0) {$\i$};
      \fill[black] (\w, 5) circle (2.5pt);
      \node[above=3pt, font=\tiny\bfseries, color=black] at (\w, 5) {$\w$};
    }

    % --- Bottom Bracket ---
    \draw[decorate, decoration={brace, mirror, raise=14pt, amplitude=6pt}, thick, black]
      (0.6, 0) -- (11.4, 0)
      node[midway, yshift=-28pt, font=\small, color=black]
      {$\gamma_3 = [4, 12, 13, 14, 5, 6, 7, 8, -2, -1, 0]$};

  \end{tikzpicture}

  \vspace{-0.8cm}
  \caption{The affine string diagram for $\gamma_3$, with $r=5$ (so $N=11$).
   }
   \label{gamma3}
\end{figure}


{\color{black}
\begin{proposition}\label{prop:first-multiplication}
For $r\geq2$ and $0\leq k\leq r-2$,
\begin{equation}\label{eq:first-multiplication}
 \C_{x_k}C_{s_{k+2}}
 =
 \begin{cases}
  \C_{x_{k+1}}+\C_e,&k=0,\\
  \C_{x_{k+1}}+\C_{x_{k-1}}+\C_{\gamma_k},&k>0.
 \end{cases}
\end{equation}
\end{proposition}
\begin{proof}
This proposition and Lemma~\ref{lem:scalar-comparison} are proved
simultaneously in \Cref{sec:first-multiplication-proof}.
\end{proof}
}


For $a,b\geq0$, let $d_{a,b}$ be the number of partitions of $a+b$
with at least $a$ parts, and put
\begin{equation}\label{eq:A-definition}
 A_{a,b}=\sum_{i=0}^{a}d_{i,b},\qquad A_{-1,b}=0.
\end{equation}
These numbers describe the finite approximation to each coefficient
$h_{x_a}$ in the following lemma.

\begin{lemma}\label{lem:scalar-comparison}
For $r\geq1$ and $0\leq a\leq r-1$,
\begin{equation}\label{eq:pascal-triangle}
 h_{x_a}=\sum_{b=0}^{r-1-a}A_{a,b}v^{a+1+2b}
       +O(v^{2r+1-a}).
\end{equation}
In particular,
\[
 h_{x_a}=(a+1)v^{a+1}+O(v^{a+3}),\qquad
 h_{x_{r-1}}=rv^r+O(v^{r+2}).
\]
\end{lemma}
\begin{proof}
For $r=1$, one has $x_0=s_0$ and $h_{s_0}=v$, so the assertion
holds. For $r\geq2$, the proof of this lemma is intertwined with
that of Proposition~\ref{prop:first-multiplication}: the multiplication
identities and the approximation \eqref{eq:pascal-triangle} are proved
simultaneously in \Cref{sec:first-multiplication-proof}.
\end{proof}

\subsection{The second chain}\label{sec:second-chain}
The first chain gives $h_{x_{{r-1}}}=rv^r+O(v^{r+2})$. We now complete
the path from $x_{r-1}$ to $\beta_r$. The indices below are chosen so
that each successive multiplication decreases the index by one.
We will use the estimate for $h_{x_{{r-1}}}$ to determine the coefficient
of $v$ at the endpoint.

\begin{definition}\label{def:second-chain}
For $0\leq s\leq r-1$, put
\begin{equation}\label{eq:second-chain-product}
 y_s=x_{r-1}(s_1\cdots s_{r-1-s}).
\end{equation}
Its window is
\begin{equation}\label{eq:second-chain}
% \begin{aligned}
   y_s=[N+1,\ldots,N+r-1-s,1-r,N+r-s,\ldots,N+r-1, 
      r+1,N+r,2-r,\ldots,0].  
% \end{aligned}
\end{equation}
Thus $y_{r-1}=x_{r-1}$, $y_0=\beta_r$, and
\[
 y_t=y_{t+1}s_{r-1-t}\qquad(0\leq t\leq r-2).
\]
The factorizations are length additive by Lemma~\ref{lem:endpoints}.
\end{definition}

{\color{black}
Our starting value is $h_{x_{r-1}}$, and our goal is to estimate
$h_{\beta_r}$. The multiplication formulas again separate the cover
from the remaining corrections. Their identity coefficients will be high
enough that the leading term can be transported along the chain.

\begin{proposition}
\label{prop:second-multiplication}
For $r\geq2$ and every $0\leq t\leq r-2$, there exists an element
$\mathcal E_t\in\M\subseteq\mathcal H$ such that
\begin{equation}\label{eq:second-multiplication}
 \C_{y_{t+1}}C_{s_{r-1-t}}
 =\C_{y_t}+\C_{y_{t+2}}+\mathcal E_t,
\end{equation}
where the term $\C_{y_{t+2}}$ is omitted when $t=r-2$, and
\[
 \eps(\mathcal E_t)=O(v^{t+3}).
\]
\end{proposition}
The proof is given in Section~\ref{sec:second-multiplication-proof}.
We first obtain a coarse estimate from the first-chain endpoint and
positivity. It already places the cover contribution beyond the degree
needed at each step.

\begin{lemma}\label{lem:second-recurrence}
Let $r\geq2$. For $0\leq s\leq r-1$,
\begin{equation}\label{eq:second-coarse-bound}
 h_{y_s}=O(v^{s+1}).
\end{equation}
Moreover,
\begin{equation}\label{eq:outer}
 h_{y_t}=\eta h_{y_{t+1}}+O(v^{t+3})\qquad(0\leq t\leq r-2).
\end{equation}
\end{lemma}
\begin{proof}
The known path from $x_{r-1}$ to $y_s$ has $r-1-s$ finite
length-increasing steps. By Corollary~\ref{cor:weak-valuation}, each step lowers the smallest
possible exponent by at most one. Since
$h_{x_{r-1}}=rv^r+O(v^{r+2})$ by Lemma~\ref{lem:scalar-comparison},
successive applications give
\[
 h_{y_s}=O\bigl(v^{\,r-(r-1-s)}\bigr)=O(v^{s+1}).
\]
This proves \eqref{eq:second-coarse-bound} independently of the
second-chain multiplication formulas.
Now apply $\eps$ to Proposition~\ref{prop:second-multiplication}.
The scalar rule \eqref{eq:scalar} gives $\eta h_{y_{t+1}}$ on the left.
When the cover is present, \eqref{eq:second-coarse-bound} gives
$h_{y_{t+2}}=O(v^{t+3})$; the remainder satisfies the same bound by the
proposition. Absorbing both contributions proves \eqref{eq:outer}.
\end{proof}

}

\subsection{Proof of Theorem B}\label{sec:scalarfinish}
We are now in position to prove \Cref{thm:B}.

\begin{proof}[Proof of Theorem \ref{thm:B}]
Let $r\geq2$ and write $e_t=h_{y_t}-\eta h_{y_{t+1}}$, so that
$e_t=O(v^{t+3})$ by \eqref{eq:outer}. Iteration gives
\begin{equation}\label{eq:second-telescope}
 h_{y_s}=\eta^{r-1-s}h_{x_{r-1}}
       +\sum_{t=s}^{r-2}\eta^{t-s}e_t
     =\eta^{r-1-s}h_{x_{r-1}}+O(v^{s+3}).
\end{equation}
The sum is empty when $s=r-1$. Each summand has order at least
$(t+3)-(t-s)=s+3$. By Lemma~\ref{lem:scalar-comparison},
$h_{x_{r-1}}=rv^r+O(v^{r+2})$. Since
$\eta^{r-1-s}=v^{-(r-1-s)}(1+v^2)^{r-1-s}$, substitution in
\eqref{eq:second-telescope} gives
\[
 h_{w_I,w_Iy_s}=h_{y_s}=rv^{s+1}+O(v^{s+3}),
\]
where the first equality is Lemma~\ref{lem:mu}. For $s=0$ we have
$y_0=\beta_r$, so taking the coefficient of $v$ yields
$\mu(w_I,w_I\beta_r)=r$.
\end{proof}

{\color{black}
\begin{example}\label{ex:mu-A4}
Here $r=2$, $N=5$, and the path is
\[
 \underbrace{[-1,3,6,7,0]}_{X_0=x_0}\xrightarrow{s_2}\underbrace{[-1,6,3,7,0]}_{x_1=y_1}
 \xrightarrow{s_1}\underbrace{[6,-1,3,7,0]}_{y_0=\beta_2}.
\]
The initial coefficient begins $h_{x_{0}}=v+v^3+O(v^5)$, so
$h_{x_1}=\eta h_{x_0}-1=2v^2+O(v^4)$. Along the second chain,
\[
 h_{\beta_2}=\eta h_{x_1}+O(v^3)=2v+O(v^3).
\]
Thus $\mu(w_I,w_I\beta_2)=2$.
\end{example}
}
\begin{example}\label{rex:roadmap}
Let $r=5$. \Cref{tab:roadmap} lists the multiplication formulas along the
chains $x_0=X_0,\ldots,x_4$ and $y_4=x_4,\ldots,y_0=\beta_5$; all
multiplicities are one. The column Cover is the previous element of the
chain (Lemma~\ref{lem:covers}), and the column Extra is the only other
term that occurs: $\C_e$ at the first step, then $\C_{\gamma_k}$
(\Cref{def:gamma}) on the first chain and $\C_{\zeta_t}$, for certain
elements $\zeta_t$ found by computer, on the second.

Applying $\eps$ to a row, the identity coefficient of the new element is
$(v+v^{-1}) h_x$ minus those of the two correction terms (Lemma~\ref{lem:scalar}).
On the first chain the terms of degree $k$ cancel, since
$h_{x_k}=(k+1)v^{k+1}+\cdots$ and $h_{\gamma_k}=v^k+\cdots$ (\Cref{lem:gaussian-products}), and
$h_{x_{k+1}}$ begins with $(k+2)v^{k+2}$. 
On the second chain both
correction terms are of order at least $t+3$, so the leading coefficient
$5$ moves down one degree at each step, giving $[v]h_{\beta_5}=5$. The
proof does not need the second-chain formulas exactly, only that the
correction terms are of this order
(Proposition~\ref{prop:second-multiplication}). 
\end{example}

\providecommand{\nocover}{\multicolumn{1}{l@{\;\phantom{+}\;}}{}}
\begin{table}[htbp]
\centering\small
\renewcommand{\arraystretch}{1.2}
\begin{tabular}{@{}l l @{\;=\;} l @{\;+\;} l @{\;+\;} l c@{}}
\toprule
Step & Product & New & Cover & Extra & New $h$ starts\\\midrule
$x_0\to x_1$ & $\C_{x_0}C_{s_2}$ & $\C_{x_1}$    & \nocover   & $\C_e$          & $2v^2$\\
$x_1\to x_2$ & $\C_{x_1}C_{s_3}$ & $\C_{x_2}$    & $\C_{x_0}$ & $\C_{\gamma_1}$ & $3v^3$\\
$x_2\to x_3$ & $\C_{x_2}C_{s_4}$ & $\C_{x_3}$    & $\C_{x_1}$ & $\C_{\gamma_2}$ & $4v^4$\\
$x_3\to x_4$ & $\C_{x_3}C_{s_5}$ & $\C_{x_4}$    & $\C_{x_2}$ & $\C_{\gamma_3}$ & $5v^5$\\\midrule
$y_4\to y_3$ & $\C_{y_4}C_{s_1}$ & $\C_{y_3}$    & \nocover   & $\C_{\zeta_3}$  & $5v^4$\\
$y_3\to y_2$ & $\C_{y_3}C_{s_2}$ & $\C_{y_2}$    & $\C_{y_4}$ & $\C_{\zeta_2}$  & $5v^3$\\
$y_2\to y_1$ & $\C_{y_2}C_{s_3}$ & $\C_{y_1}$    & $\C_{y_3}$ & $\C_{\zeta_1}$  & $5v^2$\\
$y_1\to y_0$ & $\C_{y_1}C_{s_4}$ & $\C_{\beta_5}$ & $\C_{y_2}$ & $\C_{\zeta_0}$ & $5v$\\
\bottomrule
\end{tabular}
\caption{The eight multiplication formulas for $r=5$
(\Cref{rex:roadmap}).}
\label{tab:roadmap}
\end{table}
 


% {\color{black}
% \begin{example}\label{ex:mu-A6}
% For $r=3$ and $N=7$, the first chain moves $4$ past $8,9$, and the
% second moves $-2$ past the same two entries. The rows of the table
% follow the path
% \[
%  x_0,\ x_1,\ x_2=y_2,\ y_1,\ y_0=\beta_3.
% \]
% The required coefficients are
% \begin{center}\small
% \begin{tabular}{@{}lll@{}}
% \toprule
% Window & Next generator & Required part of $h$\\\midrule
% $[-2,4,8,9,10,-1,0]$ & $s_2$ & $v+v^3+2v^5+O(v^7)$\\
% $[-2,8,4,9,10,-1,0]$ & $s_3$ & $2v^2+3v^4+O(v^6)$\\
% $[-2,8,9,4,10,-1,0]$ & $s_1$ & $3v^3+O(v^5)$\\
% $[8,-2,9,4,10,-1,0]$ & $s_2$ & $3v^2+O(v^4)$\\
% $[8,9,-2,4,10,-1,0]$ & --- & $3v+O(v^3)$\\\bottomrule
% \end{tabular}
% \end{center}
% The only relevant Gaussian correction is
% $h_{\gamma_{1}}=v+v^3+O(v^5)$. Hence
% \begin{align*}
%  h_{x_2}&=\eta h_{x_1}-h_{x_0}-h_{\gamma_1}\\
%     &=(2v+5v^3)-(v+v^3)-(v+v^3)+O(v^5)
%       =3v^3+O(v^5).
% \end{align*}
% The second chain transfers this leading coefficient to degree one.
% Consequently $\mu(w_I,w_I\beta_3)=3$.
% \end{example}
% }


\subsection{First cohomology of finite groups}\label{sec:cohomology}
To prove Theorem~\ref{thm:guralnick}, we interpret the coefficient
computed in Theorem~\ref{thm:mu} as the dimension of a first
cohomology group. The comparison below applies when
$\lambda=x\cdot0$ is $p$-restricted. We will check
these conditions for $x=\beta_r$.

Let $p$ be a prime number and let $G=\operatorname{SL}_N$ over $\overline{\mathbb F}_p$ and let
$\Phi$ be its root system. Write $\alpha_i,\varpi_i$ for the simple
roots and fundamental weights, $\theta$ for the highest root, and
$\rho=\sum_i\varpi_i$. We identify $W$ with
$p\mathbb Z\Phi\rtimes W_I$, where $\mathbb Z\Phi$ is the root lattice: $s_i$ is the finite simple reflection for $i\geq1$, and $s_0$ is the affine reflection in $\langle\lambda,\theta^\vee\rangle=p$.
For $\beta\in p\mathbb Z\Phi$, let $t_\beta$ denote the translation by $\beta$.
We use the dot action
\[
 w\cdot\lambda=w(\lambda+\rho)-\rho.
\]
For a dominant weight $\lambda$, let $L(\lambda)$ be the simple
rational $G$-module of highest weight $\lambda$.
We say a weight $\lambda$ is
$p$-restricted if $0\leq\langle\lambda,\alpha_i^\vee\rangle<p$
for every simple root $\alpha_i$. Here $H^1(G,-)$ denotes rational
cohomology; for the finite groups below, we use ordinary group
cohomology.

The first comparison is Andersen's interpretation of the spherical
KL coefficients; the precise degree-one statement is included in
\cite[Theorem~4.5]{Luebeck}, following \cite{Andersen}.
The second is the finite-group comparison
\cite[Theorem~7.5]{BNPcompare}, in the special case recorded in
\cite[Theorem~4.6]{Luebeck}.

\begin{proposition}
\label{prop:cohomology-input}
Assume Lusztig's character formula for type $A_{N-1}$ in
characteristic $p$, and suppose $p\geq5(N-1)$.
If $x\in{}^I W$ and $\lambda=x\cdot0$ is $p$-restricted, then
\begin{equation}\label{eq:cohomology-comparison}
 [v]h_x
 =\dim H^1(G,L(\lambda))
 =\dim H^1(\operatorname{SL}_N(\mathbb F_p),L(\lambda)).
\end{equation}
For fixed $N$, these assumptions on $p$ hold for every sufficiently
large prime, by \cite{AJS,Fiebig}; see also
\cite[Theorem~3.1]{Luebeck}.
\end{proposition}

We now check that the element $\beta_r$ gives a weight to which this
comparison applies.
%The restriction $r\geq2$ keeps all fundamental weight indices in the following formulas between $1$ and $N-1$.

\begin{lemma}\label{lem:highest-weight}
Let $r\geq2$, $N=2r+1$, and $p>2$.
With $\nu=\varpi_r+\varpi_{r+1}$, the window of $\beta_r$ gives
\begin{equation}\label{eq:translation-highest-weight}
 \beta_r=t_{p\nu}(s_rs_{r+1}s_r),\qquad
 \lambda_{r,p}  \coloneqq  \beta_r\cdot0
 =p\nu-2(\alpha_r+\alpha_{r+1}).
\end{equation}
Equivalently,
\begin{equation}\label{eq:restricted-highest-weight}
 \lambda_{r,p} =2\varpi_{r-1}
 +(p-2)(\varpi_r+\varpi_{r+1})+2\varpi_{r+2}.
\end{equation}
This is a nonzero $p$-restricted weight in
$\mathbb Z\Phi$.
\end{lemma}
\begin{proof}
Use the usual sum-zero coordinates for type $A_{N-1}$.
If $w(j)=a_j+m_jN$, with $1\leq a_j\leq N$, the affine action
on a vector $\mu$ is
\[
 (w\mu)_{a_j}=\mu_j+pm_j.
\]
The residue permutation of $\beta_r$ is the transposition of $r$
and $r+2$, namely $s_rs_{r+1}s_r$.
Its shifts, indexed by value residue, are
\[
 \nu=(\underbrace{1,\ldots,1}_{r},0,
       \underbrace{-1,\ldots,-1}_{r})
 =\varpi_r+\varpi_{r+1}.
\]
This proves the factorization in
\eqref{eq:translation-highest-weight}.
The finite transposition sends $\rho$ to
$\rho-2(\alpha_r+\alpha_{r+1})$, giving its weight formula.
Expanding the simple roots in fundamental weights yields
\eqref{eq:restricted-highest-weight}.
Its nonzero coefficients are $2,p-2,p-2,2$, all between $0$
and $p-1$. Finally $\nu$ is an integral sum-zero vector, so
$\nu\in\mathbb Z\Phi$; the same is true of $\lambda_{r,p} $.
\end{proof}

\begin{proof}[Proof of Theorem~\ref{thm:guralnick}]
Choose $p$ satisfying Proposition~\ref{prop:cohomology-input}.
Lemma~\ref{lem:highest-weight} verifies its weight hypotheses
for $x=\beta_r$. Equations~\eqref{eq:cohomology-comparison} and
Theorem~\ref{thm:main} give
\[
 \dim H^1(\operatorname{SL}_N(\mathbb F_p),L(\lambda_{r,p} ))
 =[v]h_{\beta_r}=r.
\]
By Steinberg's restriction theorem \cite{Steinberg}, the restriction of
$L(\lambda_{r,p} )$ to $\operatorname{SL}_N(\mathbb F_p)$ is nontrivial and
absolutely irreducible, since $\lambda_{r,p} $ is nonzero and
$p$-restricted.
Its central character is trivial because $\lambda_{r,p} \in\mathbb Z\Phi$.
It therefore descends to a nontrivial absolutely irreducible module
$\overline V$ for $\operatorname{PSL}_N(\mathbb F_p)$.
Let $Z$ be the center of $\mathrm{SL}_N(\mathbb F_p)$.
Since $Z$ acts trivially and has order prime to $p$,
its positive-degree cohomology with these coefficients
vanishes. Inflation therefore gives isomorphisms
in every degree
\cite[Proposition~1.8.11]{Sharifi}.
In particular, it identifies the first cohomology groups,
so $\dim H^1(\operatorname{PSL}_N(\mathbb F_p),\overline V)=r$.
Finally $\operatorname{PSL}_N(\mathbb F_p)$ is simple for $N\geq5$;
the kernel of its nontrivial representation on $\overline V$ is therefore
trivial, proving faithfulness.

Since $G$ is split over $\mathbb F_p$, $L(\lambda_{r,p})$ admits an
$\mathbb F_p$-form, so $\overline V=V_{r,p}\otimes_{\mathbb F_p}\overline{\mathbb F}_p$
for an $\mathbb F_p\mathrm{PSL}_N(\mathbb F_p)$-module $V_{r,p}$, which is
therefore faithful and absolutely irreducible. As the cohomology of finite
groups commutes with extension of scalars,
$\dim_{\mathbb F_p}H^1(\mathrm{PSL}_N(\mathbb F_p),V_{r,p})=r$.
This proves (1.1) for the module $V_{r,p}$ of
Theorem~\ref{thm:guralnick}.
\end{proof}





\subsection{Higher cohomology of finite simple groups}
\label{subsec:higher-cohomology}

Recall the notation and affine dot action of
Subsection~\ref{sec:cohomology}. Let $k=\overline{\mathbb F}_p$
and let $G=\mathrm{SL}_N$ be the algebraic group over $k$,
with its standard $\mathbb F_p$-structure.
We write $H^j(G,-)$ for rational cohomology and
$H^j(G(\mathbb F_q),-)$ for ordinary cohomology of the
finite group $G(\mathbb F_q)$, where $q$ is a power of $p$.
All dimensions below are taken over $k$.

Much of Subsection~\ref{sec:cohomology} generalises to all
cohomological degrees. Indeed, Andersen's theorem, in the
formulation of \cite[Theorem~4.5]{Luebeck}, gives
\begin{equation}\label{eq:higher-rational-kl}
h_x(v)
=\sum_{j\geq0}\dim_k H^j(G,L(x\cdot0))v^j
=\sum_{j\geq0}\dim_k
\operatorname{Ext}^j_G(k,L(x\cdot0))v^j
\end{equation}
for $x\in{}^I W$ such that $x\cdot0$ is $p$-restricted, provided $p\geq N$ and Lusztig's character formula holds in characteristic $p$.
The Ext groups are computed in the category of rational
$G$-modules. The polynomial denoted $m_{1,x}$ in that
reference is our $h_x$.

To prove Theorem~\ref{thm:C}, we will restrict the action
to a sufficiently large finite subgroup $G(\mathbb F_q)$.
Restriction need not preserve the cohomology dimension,
but eventual injectivity ensures that it cannot decrease.
We first verify the weight hypotheses for every element
of the second chain, rather than only its endpoint $\beta_r$.

\begin{lemma}\label{lem:second-chain-weights}
Let $r\geq2$, $N=2r+1$, $0\leq s\leq r-1$, and $p>N$.
In the sum-zero coordinates of type $A_{N-1}$, put
$\lambda_{r,s}=y_s\cdot0$. Then
\[
\lambda_{r,s}=
\bigl(
\underbrace{p,\ldots,p}_{r-1-s},
\underbrace{p-1,\ldots,p-1}_{s},
p-2,\ 0,\ -p+s+2,
\underbrace{-p,\ldots,-p}_{r-1}
\bigr).
\]
In particular, $\lambda_{r,s}$ is a nonzero $p$-restricted weight in the root lattice.
For $s=0$ this recovers the weight $\lambda_{r,p}$ of
Lemma~\ref{lem:highest-weight}.
\end{lemma}

\begin{proof}
Set $b=r-1-s$ and
\[
\nu=(\underbrace{1,\ldots,1}_{r},0,
       \underbrace{-1,\ldots,-1}_{r}).
\]
The window in \eqref{eq:second-chain} has translation part
$p\nu$, so write $y_s=t_{p\nu}\sigma_s$ with $\sigma_s\in S_N$.
Reading the positions of its residues gives
\[
\bigl(\sigma_s^{-1}(1),\ldots,\sigma_s^{-1}(N)\bigr)
=\bigl(1,\ldots,b,\ b+2,\ldots,r,\ r+2,\ r+1,\ b+1,
        \ r+3,\ldots,N\bigr),
\]
where empty intervals are omitted.
Since $\rho_i=r+1-i$, we have
\[
(y_s\cdot0)_i=p\nu_i+i-\sigma_s^{-1}(i).
\]
This proves the displayed coordinate formula.
Its coordinates sum to zero and are integers, so it lies
in the root lattice. Its consecutive coordinate differences
are nonnegative and belong to
\[
\{0,1,2,s+2,p-2,p-s-2\}.
\]
As $s\leq r-1$ and $p>2r+1$, all these differences are
less than $p$. They are the simple-coroot pairings,
proving dominance and $p$-restrictedness.
The coordinate $p-2$ is positive, so the weight is nonzero.
\end{proof}

We also use eventual injectivity of restriction.

\begin{lemma}\label{lem:eventual-injectivity}
Let $M$ be a finite-dimensional rational $G$-module over $k$
admitting an $\mathbb F_p$-form. Then, for each fixed degree
$j$ and all sufficiently large $a\in\mathbb Z_{>0}$,
restriction induces an injection
\[
H^j(G,M)\longrightarrow H^j(G(\mathbb F_{p^a}),M).
\]
\end{lemma}

Here the source is rational cohomology of the algebraic
group $G$ and the target is ordinary group cohomology
of the finite group $G(\mathbb F_{p^a})$.
The assertion over $\mathbb F_p$ is stated explicitly in \cite[Theorem~13.2]{Petrov2025}, as a consequence of \cite[Theorem~6.6]{CPSvdK1977} and \cite[Theorem~2.1]{CPS1983}. The stated version over $k$ follows by extension of scalars.

\begin{proof}[Proof of Theorem~\ref{thm:C}]
Fix $n\geq1$ and $r\geq\max\{2,n\}$, and set
$s=n-1$ and $N=2r+1$.
Choose $p>N$  large enough so that
\eqref{eq:higher-rational-kl} applies in type $A_{N-1}$,
as in Subsection~\ref{sec:cohomology}.
Put $L=L(\lambda_{r,s})$.

By Lemma~\ref{lem:second-chain-weights}, the weight
$\lambda_{r,s}$ is nonzero and $p$-restricted.
Thus \eqref{eq:higher-rational-kl} and
Theorem~\ref{thm:B} give
\[
\dim_k H^n(G,L)=[v^n]h_{y_{n-1}}(v)=r.
\]
The simple highest-weight module $L$ admits an
$\mathbb F_p$-form, so Lemma~\ref{lem:eventual-injectivity}
applies. Keeping $p$ and $L$ fixed, we obtain, for every
sufficiently large $a\in\mathbb Z_{>0}$,
\[
\dim_k H^n(G(\mathbb F_q),L)\geq r,
\qquad \mathrm{with}\  q=p^a.
\]
The required threshold may depend on $p$, $r$ and $n$.





We now repeat the argument in the proof of
Theorem~\ref{thm:A}, with $q$ in place of $p$.
Since $\lambda_{r,s}$ is nonzero and $p$-restricted, it is also $q$-restricted. By Steinberg's restriction theorem, $L(\lambda_{r,s})$ restricts to a nontrivial irreducible $k$-module for $\mathrm{SL}_N(\mathbb F_q)$. Moreover, since $\lambda_{r,s}$ lies in the root lattice, the center acts trivially, so this module descends to a nontrivial irreducible $k$-module $V$ for $\mathrm{PSL}_N(\mathbb F_q)$.
This module is faithful because $\mathrm{PSL}_N(\mathbb F_q)$
is simple ($N=2r+1\geq5$), and absolutely irreducible
because $k$ is algebraically closed.

As explained in that proof, inflation across the center,
whose order is prime to $p$, is an isomorphism in every
degree. Hence
\[
H^n(\mathrm{PSL}_N(\mathbb F_q),V)
\cong H^n(\mathrm{SL}_N(\mathbb F_q),L),
\]
and therefore
\[
\dim_k H^n(\mathrm{PSL}_N(\mathbb F_q),V)\geq r.
\]
Finally, with $n$ fixed, let $r$ tend to infinity.
This excludes an absolute bound in degree $n$.



\end{proof}

\begin{remark}
The prime $p$ is allowed to depend on $r$.
Thus Theorem~\ref{thm:C} excludes a bound depending only
on the cohomological degree, but does not exclude a bound
depending on both the degree and the characteristic.
\end{remark}

\begin{corollary}\label{cor:algebraic-cohomology}
For every $n\geq1$ and $r\geq\max\{2,n\}$, and every
sufficiently large prime $p$, the algebraic group
$G=\mathrm{SL}_{2r+1}$ over $k=\overline{\mathbb F}_p$
has a simple rational module $L$ such that
\[
\dim_k H^n(G,L)=r.
\]
In particular, these dimensions are unbounded in every
fixed positive degree.
\end{corollary}

\begin{proof}
This is the rational cohomology calculation in the
proof of Theorem~\ref{thm:C}.
\end{proof}

\subsection{Quantum groups}\label{subsec:quantum}

In the quantum case, the polynomials computed above determine
the cohomology directly.
%
Let $N\geq2$, let $\ell>N$ be odd, and let
$\zeta\in\mathbb C$ be a primitive $\ell$th root of unity.
Let $U_\zeta$ be Lusztig's quantum enveloping algebra
(with divided powers) of $\mathfrak{sl}_N(\mathbb C)$.
We work in the category of finite-dimensional integrable
type~1 modules.
We define cohomology in this category by
\[
H^n(U_\zeta,M)=
\operatorname{Ext}^n_{U_\zeta}(\mathbb C,M).
\]
For a dominant weight $\lambda$, we write $L_\zeta(\lambda)$
for the irreducible module of highest weight $\lambda$.
We use the same affine action as in
Subsection~\ref{sec:cohomology}, with $\ell$ in place of $p$:
the generators $s_i$ for $1\leq i<N$ act as the finite
simple reflections, and
\[
s_0(u)=u-\bigl(\langle u,\theta^\vee\rangle-\ell\bigr)\theta.
\]
These extend to an action of $W$ by affine transformations, and we write
$w\cdot_\ell\lambda=w(\lambda+\rho)-\rho$ for the corresponding dot action.
%
For $x\in{}^IW$, the cohomological Kazhdan--Lusztig formula
\cite[(2.3.1)]{Ko2019} gives, in our notation,
\begin{equation}\label{eq:quantum-cohomology}
\textstyle \sum_{n\geq0}
\dim_{\mathbb C}H^n
\bigl(U_\zeta,L_\zeta(x\cdot_\ell0)\bigr)v^n
=h_{w_I,w_Ix}(v)=h_x(v),
\end{equation}
where the second equality follows from Lemma~\ref{lem:mu}.

We explain how to obtain this identity from Ko's formula.
In type $A$, every positive integer $\ell$ is KL-good
\cite[\S2.2]{Ko2019}, and $\ell>N$ ensures that $-2\rho$
lies in the interior of the antidominant fundamental alcove. Ko's dot action is our $\cdot_\ell$; the difference lies
in the choice of Coxeter generators. 
We therefore take $\lambda=-2\rho$, $y=w_I$, and $w=xw_I$
in \cite[(2.3.1)]{Ko2019}. These choices give
\[
w_I\cdot_\ell(-2\rho)=0,
\qquad
(xw_I)\cdot_\ell(-2\rho)=x\cdot_\ell0.
\]
Both weights are dominant because $x\in{}^IW$.
Since the standard module of highest weight $0$ is $\mathbb C$,
the left-hand side of Ko's formula is the cohomology series
in \eqref{eq:quantum-cohomology}.

For the right-hand side, Ko uses the Coxeter generators
associated with the antidominant fundamental alcove.
Conjugation by $w_I$ identifies this Coxeter system with
ours and sends $(w_I,xw_I)$ to $(w_I,w_Ix)$.
Thus, in our normalization, Ko's right-hand side is
$h_{w_I,w_Ix}(v)$, proving \eqref{eq:quantum-cohomology}.


\begin{proof}[Proof of Theorem~\ref{thm:quantum}]
Fix $r\geq2$ and put $N=2r+1$.
For $0\leq s\leq r-1$, put
\[
\lambda_{r,s}=y_s\cdot_\ell0.
\]
The calculation in Lemma~\ref{lem:second-chain-weights}
applies with $\ell$ in place of $p$, since it uses only
the inequality $p>N$, not primality.
In particular, $\lambda_{r,s}$ is a nonzero dominant
weight. By Lemma \ref{lem:endpoints}, we have $y_s\in{}^IW$, and applying \eqref{eq:quantum-cohomology} and
Theorem~\ref{thm:B} gives
\[\textstyle 
\sum_{n\geq0}
\dim_{\mathbb C}H^n
\bigl(U_\zeta,L_\zeta(\lambda_{r,s})\bigr)v^n
=h_{y_s}(v)
=rv^{s+1}+O(v^{s+3}).
\]
Comparing coefficients we obtain that
\[
\dim_{\mathbb C}H^{s+1}
\bigl(U_\zeta,L_\zeta(\lambda_{r,s})\bigr)=r,
\]
and that the cohomology vanishes in every degree
$0\leq n<s+1$.
\end{proof}











































































































































































































































































































































































































































































































































































































































































































































































































\section{The parabolic construction of the initial coefficient}
\label{sec:parabolic}

In this section we write the canonical basis element of $X_0$ as a combination of normalized parabolic products, read off its identity coefficient, and prove one further coefficient estimate. 
The final subsection computes the polynomials $h_{\gamma_j}$. 
These results are the inputs to \Cref{sec:first-multiplication-proof}, where the multiplication formulas of \Cref{prop:first-multiplication} and the coefficient estimates of \Cref{lem:scalar-comparison} for the first chain $x_0=X_0,\ldots,x_{r-1}$ are proved together.


For large $r$ the element $\C_{X_0}$ is too long to calculate using the Kazhdan--Lusztig algorithm, and in any case we need it for all $r$ at once.  This leads us to build it as a combination of
products of symmetrizers $C_{w_P}$ of finite parabolic subgroups. 
Such products are bar-invariant by construction.
We show (Proposition~\ref{prop:canonical-lift}) that
\[
 \C_{X_0}=\R_r-\R_{r-1}+(v+v^{-1})\sum_{s=2}^{r-1}\V_s,
\]
where $\R_j$ and $\V_s$ are the products of Definition~\ref{def:products}.
By Lemma~\ref{lem:criterion} it is enough to see
that the right-hand side equals $M_{X_0}$ plus terms in $v\mathbb Z[v]M_c$; 
so we need the lowest power of $v$ in each standard coefficient. 
These coefficients are, up to explicit powers of $v$, 
constant on orbits of words in three letters
(Subsections~\ref{sec:words} and~\ref{sec:averaging}), and one computation per orbit suffices
(Proposition~\ref{prop:lowesttable}). It turns out that $\R_r-\R_{r-1}$ has the right coefficient at $X_0$ but a spurious
$-v^{-1}$ at certain orbit maxima, and the terms $(v+v^{-1})\V_s$ cancel exactly this.

Having the formula, we apply $\eps$. All factors after the first lie in $W_I$, so each acts by a scalar
\eqref{eq:eps-pi}, and the result is a sum of Gaussian binomial coefficients. The partition numbers appear in this
way in Subsection~\ref{sec:seedproof}.

\subsection{Products of symmetrizers}\label{sec:products}

We put $q=v^2$. For $n\geq1$ write
\begin{equation}\label{eq:q-integers}
 [n]_q=1+q+\cdots+q^{n-1},\qquad
 [n]_q!=\prod_{i=1}^n[i]_q,\qquad [0]_q!=1,
 \qquad \gb nk=\frac{[n]_q!}{[k]_q!\,[n-k]_q!}
\end{equation} for $0\leq k\leq n$
be the Gaussian binomial coefficients, which we set to zero outside this range.
We write $f\sim v^b$ to mean $f\in v^b\bigl(1+v^2\mathbb Z[v^2]\bigr)$.

 For $Q\subseteq I$, we define the  \emph{position blocks} of $Q$ to be 
 the orbits of $W_Q$ on 
$\{1,\ldots,N\}$. 
For a finite parabolic subset $P\subseteq S$ put
\[
 C_{w_P}=\sum_{u\in W_P}v^{\ell(w_P)-\ell(u)}H_u,\qquad
 \pi_P=\sum_{u\in W_P}v^{\ell(w_P)-2\ell(u)}.
\]
As for $I$ in \eqref{eq:finite-symmetrizer-eigenrelation}, pairing $u$ with $us$ gives
\begin{equation}\label{eq:parabolic-eigenrelation}
 C_{w_P}H_s=H_sC_{w_P}=v^{-1}C_{w_P}\qquad(s\in P),
\end{equation}
and $\pi_P$ is bar-invariant (replace $u$ by $w_Pu$). If $T\subseteq P$, splitting $W_P$ into cosets
$W_Tu$ with $u\in{}^TW_P$ gives
\begin{equation}\label{eq:parabolic-decomposition}
 C_{w_P}=C_{w_T}\sum_{u\in{}^TW_P}v^{\ell(w_P)-\ell(w_T)-\ell(u)}H_u.
\end{equation}
For $Q\subseteq I$, the Poincar\'e polynomial of $S_m$ is $[m]_q!$, so
\begin{equation}\label{eq:pi-blocks}
 \pi_Q=\prod_{B}\Pi_{|B|},\qquad \Pi_m:=v^{-\binom m2}[m]_q!,
\end{equation}
the product running over the position blocks $B$ of $Q$. In particular $\pi_Q\sim v^{-\ell(w_Q)}$.
Lemma~\ref{lem:scalar} gives $\eps(FH_u)=v^{-\ell(u)}\eps(F)$ for $u\in W_I$, hence
\begin{equation}\label{eq:eps-pi}
 \eps(F\,C_{w_Q})=\pi_Q\,\eps(F)\qquad(F\in\M,\ Q\subseteq I).
\end{equation}

\begin{lemma}\label{lem:criterion}
Let $F\in\M$ be bar-invariant, and suppose $F-M_x\in\mathcal L\coloneqq\bigoplus_{c\in{}^IW}v\mathbb Z[v]M_c$
for some $x\in{}^IW$. Then $F=\C_x$.
\end{lemma}
\begin{proof}
Since $\C_x\in M_x+\mathcal L$, the element $F-\C_x$ is bar-invariant and lies in $\mathcal L$.
By \cite[Claim~2.4]{Soergel} it is zero.
\end{proof}

\begin{definition}\label{def:products}
Let $r\geq2$ and $N=2r+1$. For a sequence $(P_1,\ldots,P_k)$ of finite parabolic subsets of $S$, put $P_0=I$ and
\begin{equation}\label{eq:D}
 \D(P_1,\ldots,P_k)=
 \frac{C_{w_I}\,C_{w_{P_1}}\cdots C_{w_{P_k}}}
 {\prod_{i=1}^{k}\pi_{P_{i-1}\cap P_i}}.
\end{equation}
Let $E=\{s_3,\ldots,s_{N-1}\}$ and, for $1\leq s<N$ and $1\leq j\leq r$,
\[
 P_s=S\setminus\{s_s\},\qquad
 J_j=S\setminus\{s_j,s_{j+1}\},\qquad
 K_j=I\setminus\{s_1,s_{j+1}\}.
\]
We set
\[
 \R_j=\D(J_j,K_j,E)\quad(1\leq j\leq r),\qquad
 \V_s=\D(P_s,E)\quad(2\leq s\leq r-1).
\]
\end{definition}

\begin{lemma}[Overlaps cancel]\label{lem:overlaps}
For finite parabolic subsets $P_1,\ldots,P_k$ the element $\D(P_1,\ldots,P_k)$ lies in $\M$ and is bar-invariant, and
\[
 \D(P_1,\ldots,P_k)\,H_s=v^{-1}\,\D(P_1,\ldots,P_k)\qquad(s\in P_k).
\]
\end{lemma}

  \begin{proof}
Put $F_i=\D(P_1,\ldots,P_i)$, so $F_0=C_{w_I}$, which satisfies the claim for $P_0=I$ by
\eqref{eq:finite-symmetrizer-eigenrelation}. Suppose $F_{k-1}\in\M$ and $F_{k-1}H_s=v^{-1}F_{k-1}$ for
$s\in P_{k-1}$. Let $T=P_{k-1}\cap P_k$. Then
$F_{k-1}C_{w_T}=\pi_TF_{k-1}$, since $C_{w_T}=\sum_{u\in W_T}v^{\ell(w_T)-\ell(u)}H_u$ and $H_u$ acts on $F_{k-1}$ by $v^{-\ell(u)}$.
Using \eqref{eq:parabolic-decomposition} with $T\subseteq P_k$,
\[
 F_k=\frac{F_{k-1}C_{w_{P_k}}}{\pi_T}
 =F_{k-1}\sum_{u\in{}^TW_{P_k}}v^{\ell(w_{P_k})-\ell(w_T)-\ell(u)}H_u .
\]
This lies in $F_{k-1}\mathcal H\subseteq\M$, and the eigenrelation for $s\in P_k$ follows from
\eqref{eq:parabolic-eigenrelation}. Bar-invariance holds because the numerator in \eqref{eq:D}
is a product of bar-invariant elements and each $\pi_P$ is bar-invariant.
\end{proof}



\begin{remark}
$\D(P_1,\ldots,P_k)$ is the character of the singular Bott--Samelson bimodule of
 \cite{Williamson-singular} attached to the sequence
$I\supset I\cap P_1\subset P_1\supset P_1\cap P_2\subset\cdots\subset P_k$. We only use the algebraic
properties of Lemma~\ref{lem:overlaps}. In general $\D(P_1,\ldots,P_k)$ is bar-invariant but not a canonical
basis element. The point of the next subsections is that $\R_r-\R_{r-1}+(v+v^{-1})\sum\V_s$ is.
\end{remark}

\begin{proposition}\label{prop:canonical-lift}
For $r\geq2$ we have
\[
 \C_{X_0}=\R_r-\R_{r-1}+(v+v^{-1})\sum_{s=2}^{r-1}\V_s .
\]
\end{proposition}
The proof is given in Subsection~\ref{sec:lowest-table}, after the expansions of $\R_j$ and $\V_s$ have been
computed.

\subsection{Words, windows and orbits}\label{sec:words}

Elements of ${}^IW$ whose window lies in $[1-N,2N]$ are encoded by words in three letters, and the sets on which
right multiplication by $C_{w_Q}$ acts turn out to be sets of rearrangements of a word. We use the standard length
formula \cite[Section~8.3]{BB}
\begin{equation}\label{eq:affine-window-length}
 \ell(w)=\sum_{1\leq i<j\leq N}
 \left|\left\lfloor\frac{w(j)-w(i)}N\right\rfloor\right|.
\end{equation}

\begin{definition}\label{def:words}
Let $t\in\{0,\ldots,r\}$. We define an
\emph{admissible word of type $t$} to be a word
$\omega=\omega_1\cdots\omega_N$ in the letters $\A,\B,\CC$ 
such that  $\A$  and $\CC$ appear  exactly $t$ times and 
$\B$ appears exactly  $N-2t$ times. 
We write $d_{\omega,i}(\mathsf X)$ for the number of times $\mathsf X$
appears in  
$\omega_1,\ldots,\omega_i$. We define the window $c_\omega=[c_\omega(1),\ldots,c_\omega(N)]$ as follows:
\[
 c_\omega(i)=
 \begin{cases}
  d_{\omega,i}(\A)-t & \text{if }\omega_i=\A,\\
  d_{\omega,i}(\B)+t & \text{if }\omega_i=\B,\\
  d_{\omega,i}(\CC)+N & \text{if }\omega_i=\CC.
 \end{cases}
\]
\end{definition}

We will use admissible words to describe permutations like those in 
\Cref{figure1,figure2} where the entries $\A,\B,\CC$ record whether an entry is non-positive, positive but less than or equal to $N$, or strictly greater than $N$ respectively.  We have that 
the entries of $c_\omega$ at the letters $\A,\B,\CC$ are, in increasing order of position 
  \begin{equation}\label{eq:lists}
   A_t=(1-t,\ldots,0),\qquad
   B_t=(t+1,\ldots,N-t),\qquad
   C_t=(N+1,\ldots,N+t).
  \end{equation}
% Shifting $A_t$ up by $N$ and $C_t$ down by $N$ turns the union of the three tuples into $\{1,\ldots,N\}$.
%
 By construction, $c_\omega\in W$ (to see this note that the entries of $c_\omega$ have distinct residues modulo $N$ and sum to $N(N+1)/2$).


\begin{lemma}\label{lem:words}
If $\omega$ is an admissible word of type $t$, then $c_\omega\in{}^IW$. Conversely, every $x\in{}^IW$
whose window lies in $[1-N,2N]$ equals $c_\omega$ for a unique admissible word $\omega$.
\end{lemma}
\begin{proof}
Let $p_1<\cdots<p_t$, $q_1<\cdots<q_{N-2t}$ and $u_1<\cdots<u_t$ be the positions of $\A$, $\B$ and $\CC$ in $\omega$.
We have that 
\[
 c_\omega^{-1}(a)=u_a-N,\qquad
 c_\omega^{-1}(t+b)=q_b,\qquad
 c_\omega^{-1}(N-t+a)=p_a+N
\]
for $1\leq a\leq t$, $1\leq b\leq N-2t$. The groups are increasing 
and   lie in $[1-N,0]$, $[1,N]$ and $[N+1,2N]$ respectively; therefore 
$c_\omega^{-1}(1)<\cdots<c_\omega^{-1}(N)$. Hence $c_\omega\in{}^IW$.

We now suppose that  $x\in{}^IW$ has window in $[1-N,2N]$.
We let $u$ (respectively  $t$) be the total number  of entries
of $x$ that are
strictly greater than $N$ (respectively non-positive). 
By assumption $x^{-1}$ is increasing on $1,\ldots,N$;
therefore these entries are $N+1,\ldots,N+u$ and $1-t,\ldots,0$. 
The remaining $N-u-t$ entries 
are given by $u+1,\ldots,N-t$ in increasing order.
The entries sum to
$N(N+1)/2+N(u-t)$ and $x\in W$ by assumption, and  so $u=t$. Replacing the three lists by $\CC$, $\A$, $\B$ gives the
unique word.
\end{proof}

\begin{lemma}\label{lem:equal-letters}
Let $\omega$ be admissible and $s_i\in I$ with $\omega_i=\omega_{i+1}$. Then $c_\omega s_i=s'c_\omega$ for some $s'\in I$,
$\ell(c_\omega s_i)=\ell(c_\omega)+1$, and $M_{c_\omega}H_{s_i}=v^{-1}M_{c_\omega}$.
\end{lemma}
\begin{proof}
Equal adjacent letters give $c_\omega(i+1)=c_\omega(i)+1$, and both values lie in one interval $[kN+1,(k+1)N]$.
Hence $c_\omega s_ic_\omega^{-1}$, which exchanges the values $c_\omega(i)$ and $c_\omega(i)+1$ (and their translates by $N\mathbb Z$),
is a simple reflection $s'\in I$, and $c_\omega(i)<c_\omega(i+1)$ gives $\ell(c_\omega s_i)=\ell(c_\omega)+1=\ell(s'c_\omega)$.
So $M_{c_\omega}H_{s_i}=C_{w_I}H_{s'c_\omega}=C_{w_I}H_{s'}H_{c_\omega}=v^{-1}M_{c_\omega}$ by
\eqref{eq:finite-symmetrizer-eigenrelation}.
\end{proof}

\begin{lemma}\label{lem:word-length}
Let $\omega$ be an admissible word of type $t$, and let $\iota(\omega)$ be its inversion number for the letter order
$\CC<\B<\A$,
\[
 \iota(\omega)=\#\{(i,j):1\leq i<j\leq N,\ \omega_i>\omega_j\}.
\]
Then
$\ell(c_\omega)=2t(N-t)-\iota(\omega)$.
\end{lemma}
\begin{proof}
Let $\omega_0=\A^t\B^{N-2t}\CC^t$. 
Its window is increasing, and in \eqref{eq:affine-window-length} only the pairs
(an entry of $A_t$, an entry of $C_t$) contribute, each by one. So $\ell(c_{\omega_0})=t^2$, while
$\iota(\omega_0)=t^2+2t(N-2t)$.
Suppose $\sigma'$ is obtained from $\sigma$ by one of the interchanges $\A\B\to\B\A$, $\A\CC\to\CC\A$ or
$\B\CC\to\CC\B$ in positions $i,i+1$. The two letters keep their counts, so $c_{\sigma'}=c_\sigma s_i$, and
$c_\sigma(i)<c_\sigma(i+1)$ because $\A$, $\B$, $\CC$ entries are $\leq0$, in $[1,N]$ and $>N$ respectively.
Hence $\ell(c_{\sigma'})=\ell(c_\sigma)+1$, whereas $\iota(\sigma')=\iota(\sigma)-1$.
Every admissible word of type $t$ is reached from $\omega_0$ by such interchanges, so
$\ell(c_\omega)+\iota(\omega)=t^2+\iota(\omega_0)=2t(N-t)$.
\end{proof}

\begin{definition}\label{def:orbit}
For $x\in{}^IW$ and $Q\subseteq I$ put $\mathcal O_Q(x)={}^IW\cap W_IxW_Q$. It consists of the minimal
representatives of the cosets $W_Ixq$ with $q\in W_Q$; it is finite, depends only on $W_IxW_Q$, and for fixed $Q$
these sets partition ${}^IW$.
\end{definition}

\begin{lemma}\label{lem:word-double-cosets}
Let $Q\subseteq I$ and 
$\omega,\omega'$ be admissible words of  type $t$.  
We write $\omega'\sim \omega$ if $\omega,\omega'$ differ only    by permutation of the 
letters within the position blocks of $Q$. 
We have that 
$W_Ic_\omega W_Q=W_Ic_{\omega'}W_Q$ if and only if $\omega'\sim \omega$;    hence  $\mathcal O_Q(c_\omega)=\{c_{\omega'}\mid \omega'\sim \omega\}$. 
Let $\omega_Q$ and $\omega^Q$ be obtained from $\omega$ by sorting within each position block of $Q$ in the orders
$\A<\B<\CC$ and $\CC<\B<\A$ respectively. Then $c_{\omega_Q}$ and $c_{\omega^Q}$ are the unique minimum and   maximum   elements of $\mathcal O_Q(c_\omega)$ in the Bruhat order. 
\end{lemma}
 \begin{proof}
Say that position $i$ of $x\in W$ \emph{carries} the letter $\A$, $\B$ or $\CC$ if $x(i)\leq0$,
$x(i)\in[1,N]$ or $x(i)>N$. Left multiplication by $W_I$ moves each value inside its interval $[kN+1,(k+1)N]$,
so it preserves the letters carried; right multiplication by $\tau\in W_Q$ rearranges them within position blocks of $Q$.
Since $c_{\omega'}$ carries the word $\omega'$, equality of double cosets forces $\omega'$ to be a rearrangement of $\omega$.

Conversely let $\omega'$ be obtained from $\omega$ by a permutation $\tau\in W_Q$ of positions. Then $c_\omega \tau$ and
$c_{\omega'}$ carry the same word, and the entries carrying each letter form the same set $A_t$, $B_t$ or $C_t$, inside one
interval $[kN+1,(k+1)N]$. Matching them residue by residue defines $\sigma\in W_I$ with $\sigma c_\omega \tau=c_{\omega'}$.
As $c_{\omega'}\in{}^IW$ (Lemma~\ref{lem:words}), it is the minimal representative of $W_Ic_\omega \tau$.
Every element of $W_Ic_\omega W_Q\cap{}^IW$ has window in $[1-N,2N]$ by the same observation, so it is some $c_{\omega'}$ by
Lemma~\ref{lem:words}. This proves the first two assertions.

An interchange of letters at positions $i,i+1$ of one block is right multiplication by $s_i\in Q$, and by the proof of
Lemma~\ref{lem:word-length} it changes the length by one and the inversion number by one in the opposite direction.
Sorting $\omega'$ within blocks by interchanges $\A\B\to\B\A$, $\A\CC\to\CC\A$, $\B\CC\to\CC\B$ ends at $\omega^Q$, and each
step raises the length by one; sorting in the reverse direction ends at $\omega_Q$ and lowers it by one. So
$c_{\omega_Q}\leq c_{\omega'}\leq c_{\omega^Q}$ in Bruhat order.
\end{proof}



\vspace{-0.5cm}
\begin{figure}[htbp]
  \centering
  \begin{tikzpicture}[
    x=0.55cm,          % Compact horizontal spacing
    y=1.1cm,           % Height proportion
    >=Stealth,
    thick,  scale=0.9
  ]
 
    \clip (-5.8, -2.6) rectangle (17.8, 6.75);
 
    \foreach \i in {-20,...,30} {
      \fill[gray!35] (\i, 0) circle (1.5pt);
      \fill[gray!35] (\i, 5) circle (1.5pt);
    }
 
    \foreach \k in {-2, -1, 0, 1, 2} {
      \foreach \i/\w in {
        1/12, 2/13, 3/14, 4/15, 5/6,
        6/-4, 7/16, 8/-3, 9/-2, 10/-1, 11/0
      } {
        \pgfmathsetmacro{\src}{\i + \k*11}
        \pgfmathsetmacro{\tgt}{\w + \k*11}
 
        \ifnum \k=0 \else
          \draw[gray!45, densely dotted, ->, looseness=0.8] (\src, 0) to[out=60, in=240] (\tgt, 5);
        \fi
      }
    }
 
    \foreach \i in {-5,...,0} {
      \node[below=3pt, font=\tiny, color=gray!80] at (\i, 0) {$\i$};
    }
    \foreach \i in {12,...,17} {
      \node[below=3pt, font=\tiny, color=gray!80] at (\i, 0) {$\i$};
    }
    \foreach \j in {-5,...,-5, 1, 2, 3, 4, 5, 7, 8, 9, 10, 11, 17} {
      \node[above=3pt, font=\tiny, color=gray!80] at (\j, 5) {$\j$};
    }
 
    \foreach \i/\w in {
      1/12, 2/13, 3/14, 4/15, 5/6,
      6/-4, 7/16, 8/-3, 9/-2, 10/-1, 11/0
    } {
      \draw[black, line width=1.4pt, ->, looseness=0.8] (\i, 0) to[out=60, in=240] (\w, 5);
 
      \fill[black] (\i, 0) circle (2.5pt);
      \node[below=3pt, font=\tiny\bfseries, color=black] at (\i, 0) {$\i$};
 
      \fill[black] (\w, 5) circle (2.5pt);
      \node[above=3pt, font=\tiny\bfseries, color=black] at (\w, 5) {$\w$};
    }
 
    \draw[decorate, decoration={brace, mirror, raise=14pt, amplitude=4pt}, thick, black]
      (0.6, 0) -- (6.4, 0);
      \draw[decorate, decoration={brace, mirror, raise=14pt, amplitude=4pt}, thick, black]
      (6.6, 0) -- (11.4, 0);

\draw[decorate, decoration={brace, mirror, raise=24pt, amplitude=4pt}, thick, green!70!black]
      (0.6, 0) -- (4.4, 0);
 \draw[decorate, decoration={brace, mirror, raise=24pt, amplitude=4pt}, thick, green!70!black]
      (4.6, 0) -- (6.4, 0);
      \draw[decorate, decoration={brace, mirror, raise=24pt, amplitude=4pt}, thick, green!70!black]
      (6.6, 0) -- (8.4, 0);
 \draw[decorate, decoration={brace, mirror, raise=24pt, amplitude=4pt}, thick, green!70!black]
      (8.6, 0) -- (11.4, 0);

\begin{scope}[black]
    \draw[decorate, decoration={brace, raise=12pt, amplitude=5pt}, thick]
      (-4.4, 5) -- (0.4, 5) node[midway, yshift=24pt, font=\small] {$\mathsf{A}$};
    \draw[decorate, decoration={brace, raise=12pt, amplitude=4pt}, thick]
      (5.6, 5) -- (6.4, 5) node[midway, yshift=23pt, font=\small] {$\mathsf{B}$};
    \draw[decorate, decoration={brace, raise=12pt, amplitude=5pt}, thick]
      (11.6, 5) -- (16.4, 5) node[midway, yshift=24pt, font=\small] {$\mathsf{C}$};
    \end{scope}
 
  \end{tikzpicture}
 
  \vspace{-1.4cm}
  \caption{The element $c_{\omega^Q}=[12,13,14,15,6,-4,16,-3,-2,-1,0]$ for $N=11$, with $\omega^Q=\CC^4\B\A\,\CC\A^4$. It is the longest element of
  $\mathcal O_Q(c_\omega)$ in two cases: for $c_\omega=X_0$ (Figure~\ref{figure2}) and $Q$ with position blocks
  $\{1,\ldots,6\},\{7,\ldots,11\}$ (black braces), and for $c_\omega=\beta_5$ (Figure~\ref{figure1}) and $Q$ with blocks
  $\{1,\ldots,4\},\{5,6\},\{7,8\},\{9,10,11\}$ (green braces). In each case $c_{\omega^Q}$ is obtained by permuting the bottom nodes inside the braces.}
  \label{figure-omegaQ}
\end{figure}

\medskip
We shall use the sets $\mathcal O_E(x)$ for $E=\{s_3,\ldots,s_{N-1}\}$, whose position blocks are $\{1\}$, $\{2\}$ and
$\{3,\ldots,N\}$. By Lemma~\ref{lem:word-double-cosets} such a set is determined by the type of $\omega$ and its first two
letters, it consists of the rearrangements of the suffix $\omega_3\cdots\omega_N$, and its longest element has the suffix
sorted in the order $\CC<\B<\A$. In every product $\R_j$ or $\V_s$ the first two letters are $\A$ or $\B$
(Lemma~\ref{lem:initial-supports}), so only four prefixes occur.

\begin{definition}\label{def:final-maxima}
A triple $(t,a,b)$ with $t\in\{0,\ldots,r\}$ and $a,b\in\{0,1\}$ is \emph{feasible} if $t\geq a+b$ and
$N-2t-2+a+b\geq0$. For such a triple put
\[
 \omega(t,a,b)=\A^a\B^{1-a}\,\A^b\B^{1-b}\,\CC^t\B^{N-2t-2+a+b}\A^{t-a-b},
 \qquad z(t,a,b)=c_{\omega(t,a,b)}\in{}^IW .
\]
\end{definition}
The pairs $(a,b)=(0,0),(0,1),(1,0),(1,1)$ correspond to the prefixes $\B\B$, $\B\A$, $\A\B$, $\A\A$. After the prefix,
the remaining letters are $t$ letters $\CC$, $N-2t-2+a+b$ letters $\B$ and $t-a-b$ letters $\A$; feasibility says
that these counts are nonnegative, so $\omega(t,a,b)$ is admissible of type $t$. Its suffix is sorted in the order
$\CC<\B<\A$, so $z(t,a,b)$ is the longest element of $\mathcal O_E(z(t,a,b))$ by
Lemma~\ref{lem:word-double-cosets}.

\begin{lemma}\label{lem:orbit-maximal-lengths}
For a feasible triple $(t,a,b)$,
\[
 \lambda(t,a,b)\coloneqq\ell(z(t,a,b))=2t(N-t)-2t-(a+b)(N-2t)+b(1-a).
\]
Moreover $(0,0,0)$ and $(r,1,0)$ are feasible, and
\begin{equation}\label{eq:orbit-maxima-special-cases}
 z(0,0,0)=e,\qquad z(r,1,0)=X_0.
\end{equation}
\end{lemma}
\begin{proof}
The suffix of $\omega(t,a,b)$ is sorted, so every inversion has its first index equal to $1$ or $2$.
The first letter forms inversions with the $t$ letters $\CC$, and also with the $N-2t$ letters $\B$ if it is $\A$:
this contributes $t+a(N-2t)$. The second letter precedes all $t$ letters $\CC$; if it is $\A$ it also precedes the
$N-2t-(1-a)$ letters $\B$ after it (the first letter is a $\B$ exactly when $a=0$). It contributes
$t+b(N-2t)-b(1-a)$. These sets of inversions are disjoint, so
$\iota(\omega(t,a,b))=2t+(a+b)(N-2t)-b(1-a)$, and Lemma~\ref{lem:word-length} gives the length.

Finally $\omega(0,0,0)=\B^N$ and $\omega(r,1,0)=\A\B\,\CC^r\A^{r-1}$ give, by the replacement rule of
Definition~\ref{def:words}, the elements $e$ and $X_0$ of \eqref{eq:endpoints}.
\end{proof}

\begin{corollary}\label{lem:endpoint-lengths}
$\ell(X_0)=2r^2-1$ and $\ell(\beta_r)=2r^2+2r-3$.
\end{corollary}
\begin{proof}
By \eqref{eq:orbit-maxima-special-cases} and Lemma~\ref{lem:orbit-maximal-lengths},
$\ell(X_0)=\lambda(r,1,0)=2r(r+1)-2r-1=2r^2-1$. The length-additive path of Lemma~\ref{lem:endpoints}
from $X_0$ to $\beta_r$ has $2r-2$ steps.
\end{proof}

% \begin{example}\label{ex:orbit-X0}
% Let $N=11$ and $r=5$. The word $\omega=\A\B\,\CC^5\A^4$ is admissible of type $5$, and the replacement rule gives
% \[
%  \begin{array}{c|ccccccccccc}
%   i&1&2&3&4&5&6&7&8&9&10&11\\\hline
%   \omega_i&\A&\B&\CC&\CC&\CC&\CC&\CC&\A&\A&\A&\A\\
%   c_\omega(i)&-4&6&12&13&14&15&16&-3&-2&-1&0
%  \end{array}
% \]
% (the letters $\A$ receive $A_5=(-4,\ldots,0)$, the $\B$ receives $B_5=(6)$ and the letters $\CC$ receive
% $C_5=(12,\ldots,16)$), so $c_\omega=X_0$ is the element drawn in Figure~\ref{figure2}, and $\omega=\omega(5,1,0)$.
% The prefix $\B\B$ is not feasible for $t=5$, since $N-2t-2=-1$; the three feasible triples of type $5$ are
% \[
%  \begin{array}{cllc}
%   (t,a,b) & \omega(t,a,b) & z(t,a,b) & \ell\\\hline
%   (5,1,0) & \A\B\,\CC^5\A^4 & [-4,6,12,13,14,15,16,-3,-2,-1,0]=X_0 & 49\\
%   (5,0,1) & \B\A\,\CC^5\A^4 & [6,-4,12,13,14,15,16,-3,-2,-1,0]=X_0s_1 & 50\\
%   (5,1,1) & \A\A\,\CC^5\B\A^3 & [-4,-3,12,13,14,15,16,6,-2,-1,0] & 48
%  \end{array}
% \]
% The orbit $\mathcal O_E(X_0)$ consists of the $\binom94=126$ elements attached to the words $\A\B\,\omega'$ with
% $\omega'$ a rearrangement of $\CC^5\A^4$. Each exchange $\A\CC\to\CC\A$ in the suffix raises the length by one
% (Lemma~\ref{lem:word-length}), so the lengths run from $29$, at $\A\B\,\A^4\CC^5$, up to $49$, at $X_0$.
% \end{example}

\begin{example}\label{ex:orbit-X0}
Let $N=11$ and $r=5$. The word $\omega=\A\B\,\CC^5\A^4$ is admissible of type $5$, and the replacement rule gives
\begin{center}
\small
\begin{tabular}{@{}l*{11}{c}@{}}
\toprule
$i$ & $1$ & $2$ & $3$ & $4$ & $5$ & $6$ & $7$ & $8$ & $9$ & $10$ & $11$\\\midrule
$\omega_i$ & $\A$ & $\B$ & $\CC$ & $\CC$ & $\CC$ & $\CC$ & $\CC$ & $\A$ & $\A$ & $\A$ & $\A$\\
$c_\omega(i)$ & $-4$ & $6$ & $12$ & $13$ & $14$ & $15$ & $16$ & $-3$ & $-2$ & $-1$ & $0$\\\bottomrule
\end{tabular}
\end{center}
(the letters $\A$ receive $A_5=(-4,\ldots,0)$, the $\B$ receives $B_5=(6)$ and the letters $\CC$ receive
$C_5=(12,\ldots,16)$), so $c_\omega=X_0$ is the element drawn in Figure~\ref{figure2}, and $\omega=\omega(5,1,0)$.
The prefix $\B\B$ is not feasible for $t=5$, since $N-2t-2=-1$; the three feasible triples of type $5$ are
\begin{center}
\small
\begin{tabular}{@{}cllc@{}}
\toprule
$(t,a,b)$ & $\omega(t,a,b)$ & $z(t,a,b)$ & $\ell$\\\midrule
$(5,1,0)$ & $\A\B\,\CC^5\A^4$ & $[-4,6,12,13,14,15,16,-3,-2,-1,0]=X_0$ & $49$\\
$(5,0,1)$ & $\B\A\,\CC^5\A^4$ & $[6,-4,12,13,14,15,16,-3,-2,-1,0]=X_0s_1$ & $50$\\
$(5,1,1)$ & $\A\A\,\CC^5\B\A^3$ & $[-4,-3,12,13,14,15,16,6,-2,-1,0]$ & $48$\\\bottomrule
\end{tabular}
\end{center}
The orbit $\mathcal O_E(X_0)$ consists of the $\binom94=126$ elements attached to the words $\A\B\,\omega'$ with
$\omega'$ a rearrangement of $\CC^5\A^4$. Each exchange $\A\CC\to\CC\A$ in the suffix raises the length by one
(Lemma~\ref{lem:word-length}), so the lengths run from $29$, at $\A\B\,\A^4\CC^5$, up to $49$, at $X_0$.
\end{example}


\subsection{Averaging over a parabolic subgroup}\label{sec:averaging}

We now describe right multiplication by $C_{w_Q}$ on the orbits $\mathcal O_Q(x)$: on each orbit the result is a scalar
multiple of a fixed weighted sum of standard basis elements, so its coefficient at the maximum determines all others.

\begin{lemma}\label{lem:parabolic-double-cosets}
Let $Q\subseteq I$, let $\omega$ be admissible, and put $y=c_{\omega_Q}$, $z=c_{\omega^Q}$ and
\[
 L=\{s_i\in Q:\ \omega_Q\text{ has equal letters in positions }i,i+1\}.
\]
\begin{enumerate}[label=\textup{(\roman*)}]
\item $L=\{s\in Q:\ ysy^{-1}\in I\}$, and for $s\in L$ we have $ys=s'y$ with $s'\in I$ and $\ell(ys)=\ell(y)+1$.
\item The map $u\mapsto yu$ is a bijection ${}^LW_Q\to\mathcal O_Q(c_\omega)$, and $\ell(yu)=\ell(y)+\ell(u)$.
\item $\ell(z)-\ell(y)=\ell(w_Q)-\ell(w_L)$, and this is the number of pairs of positions in a common block of $Q$
carrying different letters.
\end{enumerate}
\end{lemma}
\begin{proof}
(i) If $s_i\in L$, Lemma~\ref{lem:equal-letters} gives the second statement. If $\omega_Q$ has different letters at
$i,i+1$, then $ys_iy^{-1}$ exchanges $y(i)$ and $y(i+1)$ (and their translates), which lie in different intervals
$[kN+1,(k+1)N]$, so it is not in $W_I$.
(ii) is the standard parametrization of the cosets in a double coset, see \cite[Lemmas~2.12 and~2.14]{EK}.
(iii) By Lemma~\ref{lem:word-length}, $\ell(z)-\ell(y)=\iota(\omega_Q)-\iota(\omega^Q)$. Inversions between different
blocks are the same in both words; inside a block, $\omega_Q$ is sorted in the order $\A<\B<\CC$ and so has an inversion
for every pair of different letters, and $\omega^Q$ has none. Finally $W_Q$ and $W_L$ are products of symmetric groups on the blocks of $Q$
and on the runs of equal letters of $\omega_Q$, so $\ell(w_Q)-\ell(w_L)$ is the same count.
\end{proof}

\begin{lemma}\label{lem:averaging}
Let $\omega$ be admissible, $x=c_\omega$ and $Q\subseteq I$. Let $y,z$ be the minimum and maximum of
$\mathcal O=\mathcal O_Q(x)$ and $L$ as in Lemma~\ref{lem:parabolic-double-cosets}. Then for every $c\in\mathcal O$
\begin{equation}\label{eq:double-coset-averaging}
 M_cC_{w_Q}=\pi_L\,v^{\ell(y)-\ell(c)}\sum_{d\in\mathcal O}v^{\ell(z)-\ell(d)}M_d .
\end{equation}
\end{lemma}
\begin{proof}
For $s\in L$, Lemma~\ref{lem:parabolic-double-cosets}(i) and \eqref{eq:finite-symmetrizer-eigenrelation} give
$M_yH_s=C_{w_I}H_{s'}H_y=v^{-1}M_y$, hence $M_yC_{w_L}=\pi_LM_y$.
By \eqref{eq:parabolic-decomposition} and Lemma~\ref{lem:parabolic-double-cosets}(ii),
\[
 M_yC_{w_Q}=\pi_L\sum_{u\in{}^LW_Q}v^{\ell(w_Q)-\ell(w_L)-\ell(u)}M_{yu}
 =\pi_L\sum_{d\in\mathcal O}v^{\ell(z)-\ell(d)}M_d,
\]
because $\ell(z)-\ell(yu)=\ell(w_Q)-\ell(w_L)-\ell(u)$ by (iii). Finally, if $c=yu$ then $M_c=M_yH_u$ and
$H_uC_{w_Q}=v^{-\ell(u)}C_{w_Q}$, so $M_cC_{w_Q}=v^{\ell(y)-\ell(c)}M_yC_{w_Q}$.
\end{proof}

\begin{corollary}\label{cor:lowest}
In the setting of Lemma~\ref{lem:averaging}, let $f_c$ $(c\in{}^IW)$ be Laurent polynomials, almost all zero, and put
$X=\bigl(\sum_cf_cM_c\bigr)C_{w_Q}$. Then for every $d\in\mathcal O$
\begin{equation}\label{eq:averaging-sum}
 [M_d]X=\pi_L\,v^{\ell(y)}\Bigl(\sum_{c\in\mathcal O}f_c\,v^{-\ell(c)}\Bigr)v^{\ell(z)-\ell(d)},
\end{equation}
and in particular
\begin{equation}\label{eq:double-coset-coefficient-ratios}
 [M_d]X=v^{\ell(z)-\ell(d)}\,[M_z]X .
\end{equation}
Suppose moreover that $f_c\sim v^{A-\ell(c)}$ for every $c\in\mathcal O$ with $f_c\neq0$, that $u$ is the unique
longest such $c$, and that $G=X/\pi_D\in\M$ for some $D\subseteq Q$. Then, with $g=\ell(z)-\ell(u)$,
\begin{equation}\label{eq:lowest-averaging}
 [M_d]G\sim v^{A-\ell(d)-(\ell(w_Q)-\ell(w_D))+2g}\qquad(d\in\mathcal O).
\end{equation}
In particular every coefficient of $G$ on $\mathcal O$ is nonzero.
\end{corollary}
\begin{proof}
$C_{w_Q}$ lies in the subalgebra generated by the $H_s$, $s\in Q$, and each $M_cH_s$ is supported on $\mathcal O_Q(c)$.
These sets partition ${}^IW$, so only $c\in\mathcal O$ contribute on $\mathcal O$, and
\eqref{eq:averaging-sum} follows from Lemma~\ref{lem:averaging}; \eqref{eq:double-coset-coefficient-ratios} is immediate.

For \eqref{eq:lowest-averaging}: each nonzero summand $f_cv^{-\ell(c)}$ is $\sim v^{A-2\ell(c)}$, with a uniquely
smallest exponent at $c=u$, so the sum in \eqref{eq:averaging-sum} is $\sim v^{A-2\ell(u)}$. Now use $\pi_L\sim v^{-\ell(w_L)}$,
$\ell(y)-\ell(w_L)=\ell(z)-\ell(w_Q)$ (Lemma~\ref{lem:parabolic-double-cosets}(iii)) and $\ell(u)=\ell(z)-g$:
this gives $[M_d]X\sim v^{A-\ell(d)-\ell(w_Q)+2g}$. Finally $\pi_D\sim v^{-\ell(w_D)}$, and
$[M_d]G$ is a Laurent polynomial because $G\in\M$, so dividing gives \eqref{eq:lowest-averaging}.
\end{proof}

\subsection{The expansions of $\R_j$ and $\V_s$, and proof of Proposition~\ref{prop:canonical-lift}}\label{sec:lowest-table}
 We now compute the expansions of the (singular) Bott--Samelson elements 
$\R_j$ and $\V_s$ in the basis $M_c$, using the combinatorics of admissible   words from the previous section. We   use these to prove Proposition~\ref{prop:canonical-lift}.
\begin{lemma}\label{lem:initial-supports}
Let $1\leq j\leq r$ and $1\leq s<N$. Then
\begin{equation}\label{eq:initial-expansions}
 \D(J_j)=\sum_{\omega}v^{\,j(N-j-1)-\ell(c_\omega)}M_{c_\omega},\qquad
 \D(P_s)=\sum_{\omega'}v^{\,s(N-s)-\ell(c_{\omega'})}M_{c_{\omega'}},
\end{equation}
where $\omega$ runs over the admissible words with letters in $\{\A,\B\}$ in positions $\leq j$, the letter $\B$ in
position $j+1$ and letters in $\{\B,\CC\}$ in positions $\geq j+2$, and $\omega'$ runs over the admissible words with
letters in $\{\A,\B\}$ in positions $\leq s$ and in $\{\B,\CC\}$ in positions $>s$.
\end{lemma}
\begin{proof}
Let $P\subsetneq S$ and $L=I\cap P$. By \eqref{eq:parabolic-decomposition}, $C_{w_P}=C_{w_L}\sum_{u\in{}^LW_P}v^{\ell(w_P)-\ell(w_L)-\ell(u)}H_u$,
and $C_{w_I}C_{w_L}=\pi_LC_{w_I}$ by \eqref{eq:finite-symmetrizer-eigenrelation}. The left descents of $u\in W_P$ lie in $P$,
so ${}^LW_P=W_P\cap{}^IW$, and
\begin{equation}\label{eq:BSP}
 \D(P)=\sum_{u\in W_P\cap{}^IW}v^{\ell(w_P)-\ell(w_L)-\ell(u)}M_u .
\end{equation}
For $P=J_j$ and $P=P_s$ the pair $(W_P,W_L)$ is $(S_{N-1},S_j\times S_{N-j-1})$ and $(S_N,S_s\times S_{N-s})$, so
$\ell(w_P)-\ell(w_L)$ equals $j(N-j-1)$ and $s(N-s)$ respectively.

It remains to identify $W_P\cap{}^IW$. As permutations of $\mathbb Z$, $W_{P_s}$ consists of the $x$ with $x([s+1,s+N])=[s+1,s+N]$, i.e.
\begin{equation}\label{eq:Ps-window}
 x(i)\in[s+1-N,s]\ \ (1\leq i\leq s),\qquad x(i)\in[s+1,s+N]\ \ (s<i\leq N),
\end{equation}
and $W_{J_j}$ consists of the $x\in W_{P_j}$ with $x(j+1)=j+1$.
If $x\in{}^IW$ satisfies \eqref{eq:Ps-window}, its window lies in $[1-N,2N]$, so $x=c_{\omega'}$ by
Lemma~\ref{lem:words}, and \eqref{eq:Ps-window} says that $\omega'$ has no $\CC$ in positions $\leq s$ and no $\A$ in
positions $>s$. Conversely, let $\omega'$ be such a word, of type $t$. Then $t\leq\min(s,N-s)$, and the letters $\B$
in positions $\leq s$ are the first $s-t$ of them, so $c_{\omega'}$ has values in $[1-t,s]$ in positions $\leq s$ and
in $[s+1,N+t]$ in positions $>s$. Thus $c_{\omega'}$ satisfies \eqref{eq:Ps-window}, and it lies in ${}^IW$ by
Lemma~\ref{lem:words}. For $J_j$ take $s=j$: then $c_\omega(j+1)=j+1$ exactly when $\omega_{j+1}=\B$, because this $\B$ is the
$(j-t+1)$-st one and has value $t+(j-t+1)=j+1$, whereas a $\CC$ there has value $>N$.
\end{proof}

\begin{example}\label{ex:initial-N11}
Let $N=11$ and $j=5$. The product $\D(J_5)$ is the sum over words with letters in $\{\A,\B\}$ in positions $1,\ldots,5$,
a $\B$ in position $6$, letters in $\{\B,\CC\}$ in positions $7,\ldots,11$, and equally many $\A$'s and $\CC$'s. There are
$\sum_t\binom5t^2=\binom{10}5=252$ of them. The coefficient of $M_c$ is $v^{25-\ell(c)}$: it equals $v^{25}$ at $c=e$ (the word
$\B^{11}$) and $1$ at the longest element, $\A^5\B\,\CC^5$, of length $25$.
\end{example}

We can now compute the shape of the standard expansions of $\R_j$ and $\V_s$. Recall that $f\sim v^b$ means $f\in v^b(1+v^2\mathbb Z[v^2])$.

\begin{proposition}\label{prop:lowesttable}
The expansions of $\R_j$ $(1\leq j\leq r)$ and $\V_s$ $(2\leq s\leq r-1)$ in the standard basis are
\begin{align}
 \R_j&=\sum_{\substack{(t,a,b)\ \mathrm{feasible}\\a+b\leq t\leq j+a-1}}
 v^{d_j(t,a,b)}\,\rho_{j;t,a,b}(v^2)
 \sum_{c\in\mathcal O_E(z(t,a,b))}v^{\lambda(t,a,b)-\ell(c)}M_c,\label{eq:R-standard-expansion}\\
 \V_s&=\sum_{\substack{(t,a,b)\ \mathrm{feasible}\\a+b\leq t\leq s-2+a+b}}
 v^{d'_s(t,a,b)}\,\rho'_{s;t,a,b}(v^2)
 \sum_{c\in\mathcal O_E(z(t,a,b))}v^{\lambda(t,a,b)-\ell(c)}M_c,\label{eq:V-standard-expansion}
\end{align}
where $\rho_{j;t,a,b}(q),\rho'_{s;t,a,b}(q)\in\mathbb Z[q]$ have constant term $1$, and the exponents are
% \[
%  \begin{array}{c|cc}
%   (a,b)&d_j(t,a,b)&d'_s(t,a,b)\\\hline
%   (0,0)&N-2j&2(N-s-t)\\
%   (0,1)&1&N-2t-1\\
%   (1,0)&0&N-2t\\
%   (1,1)&2-(N-2j)&2(s-t).
%  \end{array}
% \]
\begin{center}
\begin{tabular}{@{}cll@{}}
\toprule
$(a,b)$ &\quad  $d_j(t,a,b)$ \quad  & $d'_s(t,a,b)$\\\midrule
$(0,0)$ & $N-2j$ & $2(N-s-t)$\\
$(0,1)$ & $1$ & $N-2t-1$\\
$(1,0)$ & $0$ & $N-2t$\\
$(1,1)$ & $2-(N-2j)$ & $2(s-t)$\\\bottomrule
\end{tabular}
\end{center}
In other words, on each orbit $\mathcal O_E(z(t,a,b))$ the coefficient at $c$ is $v^{\lambda(t,a,b)-\ell(c)}$ times the coefficient at the maximum $z(t,a,b)$, and the latter is $\sim v^{d_j(t,a,b)}$, respectively $\sim v^{d'_s(t,a,b)}$.
\end{proposition}
\begin{proof}
Throughout, an orbit $\mathcal O_E(z(t,a,b))$ is determined by $t$ and the first two letters, and
$E$ does not move them.

\emph{The case of $\R_j$.} By Lemma~\ref{lem:initial-supports}, $\D(J_j)$ has coefficient
$v^{j(N-j-1)-\ell(c)}$ at each $M_c$ in its support. The position blocks of $K_j$ are $\{1\}$, $\{2,\ldots,j+1\}$ and
$\{j+2,\ldots,N\}$. For a word of type $t$ with first letter of indicator $a$, these blocks carry respectively the
first letter, $t-a$ letters $\A$ and $j-t+a$ letters $\B$, and $t$ letters $\CC$ and $N-j-1-t$ letters $\B$. The orbit
$\mathcal O_{K_j}(c)$ meets the support of $\D(J_j)$ exactly when the second block contains a $\B$ (to fill the fixed
position $j+1$), i.e.\ $a\leq t\leq j+a-1$. In such an orbit the unique longest element of the support has each block sorted in
the order $\CC<\B<\A$ subject to that fixed $\B$; the $t-a$ letters $\A$ before it give a length gap $g=t-a$ from the maximum.
Since $\ell(w_{K_j})-\ell(w_{J_j\cap K_j})=j-1$, Corollary~\ref{cor:lowest} (with $D=J_j\cap K_j$) gives, for every $c$ with
these block counts,
\[
 [M_c]\D(J_j,K_j)\sim v^{A_K-\ell(c)},\qquad A_K=j(N-j-1)-(j-1)+2(t-a).
\]
In particular every such $c$ occurs.

Now apply $C_{w_E}$ and divide by $\pi_{K_j\cap E}$. Fix $\mathcal O_E(z(t,a,b))$ and write $k=t-a-b$ and $u=j-t+a+b-1$. The indices
occurring in $\D(J_j,K_j)$ and lying in this orbit have the following letter counts:
\begin{center}
\begin{tabular}{@{}lccc@{}}
\toprule
Positions & $\A$ & $\B$ & $\CC$\\\midrule
$3,\ldots,j+1$ & $k$ & $u$ & $0$\\
$j+2,\ldots,N$ & $0$ & $N-j-1-t$ & $t$\\\bottomrule
\end{tabular}
\end{center}
They exist exactly when all counts are nonnegative and $t\leq j+a-1$, that is, when $a+b\leq t\leq j+a-1$ (and the triple
is feasible). Sorting each of the two blocks in the order $\CC<\B<\A$ gives the unique longest of them, and its length gap from
$z(t,a,b)$ is $g=k(N-j-1)+ut$: each of the $k$ letters $\A$ precedes all $N-j-1$ letters of the second block, and each of the $u$ letters
$\B$ precedes the $t$ letters $\CC$. Moreover $\ell(w_E)-\ell(w_{K_j\cap E})=(j-1)(N-j-1)$, and
$\R_j=\D(J_j,K_j)C_{w_E}/\pi_{K_j\cap E}\in\M$ by Lemma~\ref{lem:overlaps}. Corollary~\ref{cor:lowest} therefore gives
\[
 [M_{z(t,a,b)}]\R_j\sim v^{\,A_K-\lambda(t,a,b)-(j-1)(N-j-1)+2(k(N-j-1)+ut)} .
\]
By Lemma~\ref{lem:orbit-maximal-lengths} this exponent simplifies to $d_j(t,a,b)=(1-a-b)(N-2j)+b(1+a)$, which gives the first column of the table.
By \eqref{eq:double-coset-coefficient-ratios}, $[M_c]\R_j=v^{\lambda(t,a,b)-\ell(c)}[M_{z(t,a,b)}]\R_j$ on each orbit, and the definition of $\sim$ provides
the polynomial $\rho_{j;t,a,b}$ with $\rho_{j;t,a,b}(0)=1$. All indices of $\R_j$ lie in the orbits listed, since the support of
$\D(J_j)$ is moved only by $C_{w_{K_j}}$ and then by $C_{w_E}$, each of which stays within orbits.
Summing over the orbits proves \eqref{eq:R-standard-expansion}.

\emph{The case of $\V_s$.} The product $\V_s=\D(P_s)C_{w_E}/\pi_{P_s\cap E}$ has only one averaging step. By
Lemma~\ref{lem:initial-supports}, $\D(P_s)$ has coefficient $v^{s(N-s)-\ell(c)}$, and its support consists of the words with
letters in $\{\A,\B\}$ in positions $\leq s$ and in $\{\B,\CC\}$ in positions $>s$. Since $s\geq2$, the first two letters are
$\A$ or $\B$. In $\mathcal O_E(z(t,a,b))$ the indices of the support have the following letter counts, where $k=t-a-b$:
\begin{center}
\begin{tabular}{@{}lccc@{}}
\toprule
Positions & $\A$ & $\B$ & $\CC$\\\midrule
$3,\ldots,s$ & $k$ & $s-2-k$ & $0$\\
$s+1,\ldots,N$ & $0$ & $N-s-t$ & $t$\\\bottomrule
\end{tabular}
\end{center}
All counts are nonnegative exactly when $a+b\leq t\leq s-2+a+b$ (note $N-s-t\geq0$ since $t\leq s$ and $s\leq r-1$). The longest of them has each block
sorted, with length gap $g=k(N-s)+(s-2-k)t$ from $z(t,a,b)$, and $\ell(w_E)-\ell(w_{P_s\cap E})=(s-2)(N-s)$.
Corollary~\ref{cor:lowest} gives $[M_{z(t,a,b)}]\V_s\sim v^{\,s(N-s)-\lambda(t,a,b)-(s-2)(N-s)+2g}=v^{\,2(N-s)-\lambda(t,a,b)+2g}$, which simplifies, again by
Lemma~\ref{lem:orbit-maximal-lengths}, to the exponent $d'_s(t,a,b)$ of the table. The rest is as for $\R_j$.
\end{proof}

\begin{example}\label{ex:lowesttable-N11}
Let $N=11$, $r=5$, $j=5$. In the standard basis $\R_5=\D(J_5,K_5,E)$ has $11\,131$ terms. By
Proposition~\ref{prop:lowesttable} they lie in $18$ orbits $\mathcal O_E(z(t,a,b))$, namely those with
\[
 0\leq t\leq4\ (\B\B),\qquad 1\leq t\leq4\ (\B\A),\qquad 1\leq t\leq5\ (\A\B),\qquad 2\leq t\leq5\ (\A\A),
\]
and on each of them the coefficient of $M_c$ is $v^{d_5+\lambda-\ell(c)}\rho(v^2)$, with $d_5=1,1,0,1$ for the four prefixes.
So $18$ polynomials describe all $11\,131$ coefficients.

The orbit of $X_0=z(5,1,0)$ (Example~\ref{ex:orbit-X0}) has $t=5$. Since $\R_4$ and the $\V_s$ only occur for $t\leq4$, none of them meets this orbit, so there $\C_{X_0}$ agrees with $\R_5$. As $[M_{X_0}]\C_{X_0}=1$ and $d_5(1,0)=0$, we get
$\rho_{5;5,1,0}=1$, and Corollary~\ref{cor:lowest} gives
\[
 [M_c]\,\C_{X_0}=v^{\,49-\ell(c)}\qquad(c\in\mathcal O_E(X_0)),
\]
since $\ell(X_0)=2r^2-1=49$. The orbit has $126$ elements with $\ell(c)$ from $29$ to $49$: the coefficient is $1$ at $X_0$, then $v$ at the elements of length
$48$, and so on down to $v^{20}$ at $\A\B\,\A^4\CC^5$.
\end{example}

\begin{proof}[Proof of Proposition~\ref{prop:canonical-lift}]
Put $F=\R_r-\R_{r-1}+(v+v^{-1})\sum_{s=2}^{r-1}\V_s$. By Lemma~\ref{lem:overlaps}, $F\in\M$ is bar-invariant, so by
Lemma~\ref{lem:criterion} it suffices to show $F-M_{X_0}\in\bigoplus_{c\in{}^IW}v\mathbb Z[v]M_c$.

By Proposition~\ref{prop:lowesttable}, $F$ is supported on the orbits $\mathcal O_E(z)$, $z=z(t,a,b)$, and on each of them
$[M_c]F=v^{\lambda(t,a,b)-\ell(c)}[M_z]F$. Since $z$ is the unique longest element of its orbit (Lemma~\ref{lem:word-double-cosets}),
the exponent is positive for $c\neq z$. It therefore suffices to show
\begin{equation}\label{eq:canonical-lift-maximal-coefficients}
 [M_z]F\in
 \begin{cases}
  1+v\mathbb Z[v],&z=X_0,\\
  v\mathbb Z[v],&z\neq X_0,
 \end{cases}
\end{equation}
for the maxima $z=z(t,a,b)$, where $z(r,1,0)=X_0$ by \eqref{eq:orbit-maxima-special-cases}.

By Proposition~\ref{prop:lowesttable}, the coefficients of $\R_j$ and $\V_s$ at $z(t,a,b)$ are $v^{d_j}\rho_j(v^2)$ and
$v^{d'_s}\rho'_s(v^2)$ with $\rho(0)=1$, so only the lowest exponents matter. For $\R_r$, $\R_{r-1}$ and $(v+v^{-1})\V_s$
(whose exponents are those of $\V_s$ lowered by one) they are, in the columns $(a,b)=(0,0),(0,1),(1,0),(1,1)$, that is $\B\B,\B\A,\A\B,\A\A$:
\begin{center}
\renewcommand{\arraystretch}{1.2}
\begin{tabular}{@{}lcccc@{}}
\toprule
 & $\B\B$ & $\B\A$ & $\A\B$ & $\A\A$\\
\midrule
$\R_r$     & $1$ & $1$ & $0$ & $1$\\
$\R_{r-1}$ & $3$ & $1$ & $0$ & $-1$\\
$(v+v^{-1})\V_s$ & $2(N-s-t)-1$ & $N-2t-2$ & $N-2t-1$ & $2(s-t)-1$\\
\bottomrule
\end{tabular}
\end{center}
The supports are $a+b\leq t\leq r+a-1$ for $\R_r$, $t\leq r+a-2$ for $\R_{r-1}$ and $t\leq s-2+a+b$ for $\V_s$.
In these ranges all entries are positive except the zeros in the column $\A\B$ and the $-1$ in the column $\A\A$, which occurs for
$\R_{r-1}$ when $2\leq t\leq r-1$ and for $(v+v^{-1})\V_s$ when $t=s$. Hence:
\begin{itemize}
\item for $\B\B$ and $\B\A$ the coefficient lies in $v\mathbb Z[v]$;
\item for $\A\B$ the coefficients of $\R_r$ and $\R_{r-1}$ both lie in $1+v\mathbb Z[v]$ and cancel for $1\leq t\leq r-1$; at $t=r$ only
$\R_r$ contributes, and this is $z=X_0$;
\item for $\A\A$, the term $-\R_{r-1}$ contributes $-v^{-1}+v\mathbb Z[v]$ at $z(t,1,1)$ for $2\leq t\leq r-1$, and the term
$(v+v^{-1})\V_t$ contributes $v^{-1}+v\mathbb Z[v]$ at the same maximum. These cancel to leave $v\mathbb Z[v]$, as do all other contributions.
\end{itemize}
This proves \eqref{eq:canonical-lift-maximal-coefficients}, and Lemma~\ref{lem:criterion} gives $F=\C_{X_0}$.
For $r=2$ all ranges involving $2\leq s\leq r-1$ are empty.
\end{proof}

\begin{example}\label{ex:canonical-lift}
For the smallest case in which the correction term is present, $r=3$ and $N=7$, we have
$\C_{X_0}=\R_3-\R_2+(v+v^{-1})\V_2$. The coefficients at four of the maxima are
\begin{center}
\renewcommand{\arraystretch}{1.2}
\begin{tabular}{@{}lcccc@{}}
\toprule
$z$ & $\R_3$ & $\R_2$ & $(v+v^{-1})\V_2$ & $\C_{X_0}$\\
\midrule
$z(3,1,0)=X_0$ & $1$ & $0$ & $0$ & $1$\\
$z(2,1,0)$ & $1+2v^2+v^4$ & $1$ & $0$ & $2v^2+v^4$\\
$z(2,1,1)$ & $v+2v^3+2v^5+v^7$ & $v^{-1}+v+v^3$ & $v^{-1}+v$ & $v+v^3+2v^5+v^7$\\
$z(0,0,0)=e$ & $v+2v^3+4v^5+O(v^7)$ & $v^3+2v^5+O(v^7)$ & $O(v^9)$ & $v+v^3+2v^5+O(v^7)$\\
\bottomrule
\end{tabular}
\end{center}
At $X_0$ only $\R_3$ contributes. At $z(2,1,0)$, and likewise at $z(1,1,0)$, the constants of $\R_3$ and $\R_2$ cancel.
At $z(2,1,1)$ the term $-\R_2$ would leave a coefficient $-v^{-1}$, and it is removed by $(v+v^{-1})\V_2$.
The last row is the coefficient of $M_e$, that is $h_{X_0}$; it begins $v+v^3+2v^5$, with coefficients
$\Part(0),\Part(1),\Part(2)=1,1,2$ as in Proposition~\ref{prop:seed}.
\end{example}

\subsection{Coefficient estimates}\label{sec:seedproof}

We now read off coefficients of $\C_{X_0}$ and of $\C_{\gamma_j}$ from the products above. The identity coefficient is
$h_{X_0}=\eps(\C_{X_0})$, and in each product $\R_j$ and $\V_s$ every factor after the initial one lies in $I$, so by
\eqref{eq:eps-pi} it acts on $\eps$ by a scalar. The computation reduces to Gaussian binomial coefficients, and this
proves Proposition~\ref{prop:seed}. We then record one further coefficient of $\C_{X_0}$, needed in \Cref{lem:kappa-vanishing}, and the identity coefficient of $\C_{\gamma_j}$.

\begin{proof}[Proof of Proposition~\ref{prop:seed}]
By Lemma~\ref{lem:initial-supports}, $\eps(\D(J_j))=v^{j(N-j-1)}$ and $\eps(\D(P_s))=v^{s(N-s)}$ (the coefficient at $M_e$, the word $\B^N$).
The blocks of $K_j,J_j\cap K_j,E,K_j\cap E$ have sizes $(1,j,N-j-1)$, $(1,j-1,1,N-j-1)$, $(1,1,N-2)$, $(1,1,j-1,N-j-1)$,
and those of $E$, $P_s\cap E$ have sizes $(1,1,N-2)$, $(1,1,s-2,N-s)$. So \eqref{eq:eps-pi} and \eqref{eq:pi-blocks} give
\begin{align}
 \eps(\R_j)&=v^{j(N-j-1)}\,\frac{\Pi_j}{\Pi_{j-1}}\,\frac{\Pi_{N-2}}{\Pi_{j-1}\Pi_{N-j-1}}
 =v^{N-2j}\,[j]_q\gb{N-2}{j-1},\\[4pt]
 \eps(\V_s)&=v^{s(N-s)}\,\frac{\Pi_{N-2}}{\Pi_{s-2}\Pi_{N-s}}
 =v^{2(N-s)}\gb{N-2}{s-2}.
\end{align}
In particular $(v+v^{-1})\,\eps(\V_s)=O(v^{2r+3})$ for $2\leq s\leq r-1$. Applying $\eps$ to
Proposition~\ref{prop:canonical-lift} and using the Gaussian identity
\[
 [r]_q\gb{2r-1}{r-1}-q\,[r-1]_q\gb{2r-1}{r-2}
 =\gb{2r-1}{r-1}+(1+q)\,q^{r}\gb{2r-1}{r-2},
\]
we obtain
\begin{align}
 h_{X_0}
 &=v\biggl([r]_q\gb{2r-1}{r-1}-q\,[r-1]_q\gb{2r-1}{r-2}\biggr)+O(v^{2r+3})\\[4pt]
 &=v\gb{2r-1}{r-1}+v\,(1+q)\,q^{r}\gb{2r-1}{r-2}+O(v^{2r+3})\\[4pt]
 &=v\gb{2r-1}{r-1}+O(v^{2r+1}).
\end{align}
Finally $\gb{2r-1}{r-1}$ counts partitions in a $(r-1)\times r$ box by size \cite[Proposition~1.7.3]{StanleyEC1}, so
\[
 \gb{2r-1}{r-1}\equiv\sum_{d=0}^{r-1}\Part(d)\,q^{d}\pmod{q^{r}}.
\]
Substituting $q=v^2$ proves \eqref{eq:seed}.
\end{proof}

\begin{lemma}\label{lem:kappa-seed-coefficient}
Let $r\geq2$, $N=2r+1$, and put $z=z(1,1,0)=c_{\A\B\CC\B^{N-3}}=[0,2,N+1,3,\ldots,N-1]$. For $1\leq j\leq r$,
\begin{equation}\label{eq:kappa-seed-R}
 [M_z]\R_j=[j]_q\gb{N-3}{j-1}.
\end{equation}
Consequently
\begin{equation}\label{eq:kappa-seed-coefficient}
 [M_z]\C_{X_0}=O(v^{N-3}).
\end{equation}
\end{lemma}
\begin{proof}
Let $y=c_{\A\B^{N-2}\CC}$ $(=s_0)$, the minimum of $\mathcal O_E(z)$, which consists of the type-one words with prefix $\A\B$.
The maximum is $z$, of length $N-2$, since moving $\CC$ one position to the left across a $\B$ raises the length by one.
Here $L=\{s_3,\ldots,s_{N-2}\}$ and Lemma~\ref{lem:averaging} gives $[M_z](M_yC_{w_E})=\pi_L$.

Fix $j$ and put $H=K_j\cap E$. In $\D(J_j)$ only the type-one words with first letter $\A$ can contribute to $M_z$,
because $K_j$ and $E$ preserve the type and the first letter. By Lemma~\ref{lem:initial-supports} these words are
$\A\B^{p-2}\CC\B^{N-p}$ with $j+2\leq p\leq N$; they form the orbit $\mathcal O_H(y)$ and the longest of them has length $N-j-1$.
Their part of $\D(J_j)$ is therefore, by Lemma~\ref{lem:averaging} for $Q=H$,
\[
 F_j=\sum_{c\in\mathcal O_H(y)}v^{j(N-j-1)-\ell(c)}M_c=\frac{v^{(j-1)(N-j-1)}}{\pi_{H\cap L}}\,M_yC_{w_H}.
\]
We claim $F_jH_s=v^{-1}F_j$ for every $s\in K_j$. For $s\in H$ this follows from \eqref{eq:parabolic-eigenrelation}, since $F_j$ is a multiple of
$M_yC_{w_H}$. The only generator of $K_j$ outside $H$ is $s_2$ (present when $j\geq2$), and in every word of the support
positions $2$ and $3$ both carry $\B$, so $M_cH_{s_2}=v^{-1}M_c$ by Lemma~\ref{lem:equal-letters}. Hence $F_jC_{w_{K_j}}=\pi_{K_j}F_j$.
The normalization of $\R_j$ is $\pi_{J_j\cap K_j}\pi_H$, and $C_{w_H}C_{w_E}=\pi_HC_{w_E}$ by \eqref{eq:parabolic-eigenrelation}, so
\[
 [M_z]\R_j
 =v^{(j-1)(N-j-1)}\frac{\pi_{K_j}}{\pi_{J_j\cap K_j}}\,\frac{\pi_L}{\pi_{H\cap L}}
 =v^{(j-1)(N-j-1)}\frac{\Pi_j}{\Pi_{j-1}}\,\frac{\Pi_{N-3}}{\Pi_{j-1}\Pi_{N-j-2}} .
\]
Here we used \eqref{eq:pi-blocks}: the blocks of $K_j$ and $J_j\cap K_j$ have sizes $(1,j,N-j-1)$ and $(1,j-1,1,N-j-1)$, while $W_L\cong S_{N-3}$ and
$W_{H\cap L}\cong S_{j-1}\times S_{N-j-2}$. The two quotients equal $v^{-(j-1)}[j]_q$ and $v^{-(j-1)(N-j-2)}\gb{N-3}{j-1}$, and their powers of $v$
cancel against $v^{(j-1)(N-j-1)}$.
This proves \eqref{eq:kappa-seed-R}.

Put $G_0=\gb{2r-2}{r-1}$ and $G_1=\gb{2r-2}{r-2}$. Then $[r-1]_qG_0=[r]_qG_1$, so
$G_1=[r-1]_qG_0/[r]_q$ and $[r]_q-[r-1]_q=q^{r-1}$ give
\[
 [M_z](\R_r-\R_{r-1})=[r]_qG_0-[r-1]_qG_1=q^{r-1}(G_0+G_1).
\]
This is divisible by $q^{r-1}=v^{N-3}$. For $2\leq s\leq r-1$, the case $(t,a,b)=(1,1,0)$ of Proposition~\ref{prop:lowesttable} gives
$[M_z]\V_s=O(v^{N-2})$, so $(v+v^{-1})[M_z]\V_s=O(v^{N-3})$ as well. Proposition~\ref{prop:canonical-lift} now proves
\eqref{eq:kappa-seed-coefficient}. For $r=2$ the sum over $s$ is empty and $[M_z]\C_{X_0}=q(2+q)=2v^2+v^4$.
\end{proof}

\begin{lemma}\label{lem:gaussian-products}
Let $r\geq2$, $1\leq j\leq r$ and $Q=\{s_2,\ldots,s_{N-1}\}$. Then
\begin{equation}\label{eq:gaussian-products}
 \D(P_{N-j},Q)=\C_{\gamma_j},\qquad h_{\gamma_j}=v^j\gb{N-1}{j}.
\end{equation}
\end{lemma}
\begin{proof}
By Lemma~\ref{lem:initial-supports}, the coefficient at $M_c$ in $\D(P_{N-j})$ is $v^{j(N-j)-\ell(c)}$ whenever $c$ is in its support, and
the admissible words there have only $\A,\B$ in the first $N-j$ positions and only $\B,\CC$ in the last $j$.

The position blocks of $Q$ are $\{1\}$ and $\{2,\ldots,N\}$, so a set $\mathcal O_Q(c)$ meeting the support is determined by its
type $u\leq j$ and its first letter ($\B$ or $\A$). Its maximum is obtained by sorting the suffix in the order $\CC<\B<\A$;
the longest index in the support is likewise unique (sort within the two blocks of the initial support, keeping the first letter fixed, as in
Lemma~\ref{lem:word-double-cosets}). Let $\lambda$ be the length of the maximum and $g$ its length gap from the longest index in the support. By
Lemma~\ref{lem:word-length}:
\begin{center}
\begin{tabular}{@{}clll@{}}
\toprule
First letter & Range & $\lambda$ & $g$\\\midrule
$\B$ & $0\leq u\leq j$ & $2u(N-u)-u$ & $u(N-u-1)$\\
$\A$ & $1\leq u\leq j$ & $2u(N-u)+u-N$ & $u(N-u)-j$\\\bottomrule
\end{tabular}
\end{center}
For instance, in the first row the longest word in the support is $\B^{N-j-u}\A^u\CC^u\B^{j-u}$, of length $u(N-u)$, and sorting its suffix gives the
displayed $g$. In the second row the first $\A$ is fixed and the longest word is $\A\B^{N-j-u}\A^{u-1}\CC^u\B^{j-u}$, of length $u(N-u)+u+j-N$.

The normalizing length difference is $\ell(w_Q)-\ell(w_{P_{N-j}\cap Q})=j(N-j-1)$, so Corollary~\ref{cor:lowest} (with $A=j(N-j)$) gives lowest exponents
$j-\lambda+2g$, that is $j-u$ and $N-j-u$ at the two kinds of maxima. The latter is positive since $N-2j\geq1$. The unique constant coefficient is
therefore at the first-row maximum with $u=j$, whose word is $\B\CC^j\B^{N-2j-1}\A^j$, namely $\gamma_j$. All other coefficients acquire a positive
power of $v$ by \eqref{eq:double-coset-coefficient-ratios}. Integrality, bar-invariance (Lemma~\ref{lem:overlaps}) and
Lemma~\ref{lem:criterion} now give $\D(P_{N-j},Q)=\C_{\gamma_j}$.
%
Finally, \eqref{eq:eps-pi} and Lemma~\ref{lem:initial-supports} give
\[
 h_{\gamma_j}=v^{j(N-j)}\frac{\Pi_{N-1}}{\Pi_{N-j-1}\Pi_j}=v^j\gb{N-1}{j},
\]
as required.
\end{proof}


















































\section{Common tools for the two chains}
\label{sec:corrections}
This section collects the tools used in the multiplication formulas for both chains (\Cref{prop:first-multiplication,prop:second-multiplication}).
After the general KL rule and descent restrictions, we record a positivity
bound for standard coefficients along finite increasing steps.
We then classify the possible correction indices and determine the cover
contributions. The expansions in Subsection~\ref{sec:candidate-expansions}
separate these known contributions from the remaining multiplicities.
Finally, affine cancellation reduces the required coefficients by deleting
two common positions at a time. Section~\ref{sec:first-multiplication-proof}
uses these tools in an induction along the first chain, with the rank fixed;
Section~\ref{sec:second-multiplication-proof} then treats the second chain
without induction on the rank. Both sections also invoke results in smaller
rank, but only ones already established for every~$r$.

\subsection{General tools for multiplication and Bruhat order}
\label{sec:standard-multiplication-rule}
For $w\in W$ and $x\in{}^I W$, write
\[
 R(w)=\{s\in S:ws<w\},\qquad R_I(x)=R(w_Ix).
\]

The following is the standard KL multiplication rule
\cite[(2.3.a)--(2.3.c)]{KL}, expressed in the spherical module
via the identification in \cite[Proposition~3.4(1)]{Soergel}.

\begin{proposition}\label{prop:KLrule}
Let $s\in S$ and $x,xs\in{}^I W$ with $xs>x$. Then
\begin{equation}\label{eq:KLrule}
 \C_xC_s=\C_{xs}
 +\sum_{\substack{c\in{}^I W\\c<x,\ s\in R_I(c)}}\mus(c,x)\C_c.
\end{equation}
\end{proposition}

For $c,x\in{}^I W$, write $c\triangleleft x$ if $c<x$ and
$\ell(x)=\ell(c)+1$. Such a cover satisfies $\mus(c,x)=1$,
since $w_Ic\triangleleft w_Ix$. We will use the following precise
form of the descent restriction, including its consequence for non-covers.

\begin{lemma}\label{lem:descents}
Let $c,x\in{}^I W$ satisfy $c<x$ and $\mus(c,x)\ne0$.
If $s\in R_I(x)\setminus R_I(c)$, then
\begin{equation}\label{eq:strong-descents}
 x=cs.
\end{equation}
In particular, if $\ell(x)-\ell(c)>1$, then
\begin{equation}\label{eq:desc}
 R_I(x)\subseteq R_I(c).
\end{equation}
\end{lemma}
\begin{proof}
Apply the right-handed form of \cite[Corollary~5.3]{Casselman}
to $w_Ic<w_Ix$. It gives $w_Ix=w_Ics$ when the leading
coefficient is nonzero and the upper element has a right descent
missing from the lower one. Cancelling $w_I$ proves
\eqref{eq:strong-descents}. This equality has length difference
one, giving \eqref{eq:desc} for non-covers.
\end{proof}

We collect the nonzero correction indices in \eqref{eq:KLrule} as follows.

\begin{definition}\label{def:correction-candidates}
Let $s\in S$ and $x,xs\in{}^I W$ satisfy $xs>x$. Define
\begin{equation}\label{eq:correction-candidates}
\begin{aligned}
 \mathcal C(x,s)
 &=\{c\in{}^I W:c<x,\ s\in R_I(c),\ \mus(c,x)\ne0\},\\
 \mathcal C_{\mathrm{cov}}(x,s)
 &=\{c\in\mathcal C(x,s):c\triangleleft x\},\\
 \mathcal C_{\mathrm{nc}}(x,s)
 &=\mathcal C(x,s)\setminus\mathcal C_{\mathrm{cov}}(x,s).
\end{aligned}
\end{equation}
These are the correction candidates in \eqref{eq:KLrule},
separated into covers and non-covers.
\end{definition}

We also use the positivity of ordinary KL polynomials
\cite[Corollary~1.2(1)]{EW14}. The standard coefficients of
$\C_x$ are the ordinary coefficients $h_{w_Ic,w_Ix}$;
thus they, and the integers $\mus(c,x)$, are nonnegative.
For Laurent polynomials, write $f\preceq g$ when
$g-f\in\mathbb Z_{\geq0}[v,v^{-1}]$.
We write $\ord(f)$ for the smallest exponent of a nonzero Laurent
polynomial or series, and put $\ord(0)=+\infty$.

\begin{lemma}\label{lem:coefficient-positivity}
Let $s\in I$ and $c\in{}^I W$ satisfy $cs\notin{}^I W$.
For every $F\in\M$,
\begin{equation}\label{eq:singleton-coefficient}
 [M_c](FC_s)=(v+v^{-1})[M_c]F.
\end{equation}
If $x,xs\in{}^I W$ and $xs>x$, then
\begin{equation}\label{eq:coefficient-positivity}
\begin{aligned}
 0\preceq[M_c]\C_{xs}
 &\preceq(v+v^{-1})[M_c]\C_x,\\
 \ord\bigl([M_c]\C_{xs}\bigr)
 &\geq\ord\bigl([M_c]\C_x\bigr)-1.
\end{aligned}
\end{equation}
\end{lemma}
\begin{proof}
Since $c$ is minimal and $cs$ is not, $cs=tc$ for some $t\in I$.
The standard spherical multiplication rule gives
$M_cC_s=(v+v^{-1})M_c$.
No distinct minimal index can contribute to $M_c$: its only possible
partner under right multiplication by $s$ would be $cs$, which is not
minimal. This proves \eqref{eq:singleton-coefficient} for arbitrary $F$.

By \eqref{eq:KLrule}, $\C_xC_s-\C_{xs}$ is a sum of canonical
elements with nonnegative multiplicities. Their standard coefficients
are nonnegative by KL positivity. Taking the coefficient of $M_c$ and
using \eqref{eq:singleton-coefficient} gives the first line of
\eqref{eq:coefficient-positivity}. Multiplication by $v+v^{-1}$
lowers the smallest possible exponent by one, proving the second line.
\end{proof}

\begin{corollary}\label{cor:weak-valuation}
Let $u_0,\ldots,u_L\in{}^I W$ satisfy
$u_i=u_{i-1}t_i>u_{i-1}$ with $t_i\in I$. Then
\begin{equation}\label{eq:weak-valuation}
 \ord(h_{u_L})\geq\ord(h_{u_0})-L.
\end{equation}
\end{corollary}
\begin{proof}
Take $c=e$ in Lemma~\ref{lem:coefficient-positivity} at each step.
Since $t_i\in I$, the element $et_i$ is not a minimal left
representative. Iteration gives \eqref{eq:weak-valuation}.
\end{proof}

We test Bruhat comparisons using the following standard criterion.

\begin{lemma}\label{lem:affine-rank-criterion}
For $w\in W$ and $m,M\in\mathbb Z$, define
\begin{equation}\label{eq:affine-rank-functions}
\begin{aligned}
 R_w(m,M)&=\#\{i\in\mathbb Z:i\leq m,\ w(i)\geq M\},\\
 \widetilde R_w(m,M)&=\#\{i\in\mathbb Z:i>m,\ w(i)\leq M\}.
\end{aligned}
\end{equation}
These numbers are finite. For $y,w\in W$,
\begin{equation}\label{eq:affine-rank-criterion}
\begin{aligned}
 y\leq w
 &\quad\Longleftrightarrow\quad
 R_y(m,M)\leq R_w(m,M)
 \quad\text{for all }m,M\in\mathbb Z\\
 &\quad\Longleftrightarrow\quad
 \widetilde R_y(m,M)\leq\widetilde R_w(m,M)
 \quad\text{for all }m,M\in\mathbb Z.
\end{aligned}
\end{equation}
\end{lemma}
\begin{proof}
 The criterion involving
$R$ is \cite[Theorem~8.3.7]{BB}; see also
\cite[Proposition~4.1]{Henderson}. Since inversion preserves Bruhat
order, applying this criterion to inverses and using
\[
 \widetilde R_w(m,M)=R_{w^{-1}}(M,m+1)
\]
gives the equivalent criterion involving $\widetilde R$.
\end{proof}

\subsection{Possible correction indices for the two chains}
We now apply the general tools above to our two chains. We first
introduce the elements $\Theta_j$, defined below, associated to the
fully sorted admissible words. The element $\Theta_r$ will serve
as a common Bruhat upper bound for both chains, as verified in
the proof of Lemma~\ref{lem:candidates}. The family $\Theta_j$
will also appear among the possible correction indices. The next
two lemmas describe the minimal representatives below $\Theta_r$
by admissible words and give the descent and prefix-count tests
used to restrict those words.

\begin{definition}\label{def:theta}
For $0\leq j\leq r$, let $\Theta_j$ be the element associated to
the admissible word of type $j$
\[
 \CC^j\B^{N-2j}\A^j.
\]
Its window is
\[
 \Theta_j=[N+1,\ldots,N+j,j+1,\ldots,N-j,1-j,\ldots,0].
\]
In particular, $\Theta_0=e$.
\end{definition}

\begin{lemma}\label{lem:word-ideal}
If $c\in{}^I W$ satisfies $c\leq\Theta_r$, then $c=c_\omega$
for an admissible word of some type $0\leq t\leq r$.
\end{lemma}
\begin{proof}
We have
\[
 R_{\Theta_r}(N,2N+1)=\widetilde R_{\Theta_r}(0,-N)=0.
\]
Indeed, if $i\leq N$, write $i=a+qN$ with
$1\leq a\leq N$ and $q\leq0$. Since the largest entry of
the window of $\Theta_r$ is $N+r$, periodicity gives
\[
 \Theta_r(i)=\Theta_r(a)+qN\leq N+r<2N+1.
\]
Thus  $\{ i \in \Z \mid i\leq N, \,\Theta_r(i)\geq 2N+1 \}$ is empty and   $R_{\Theta_r}(N,2N+1)=0$.
The other equality follows similarly.



Since $c\leq\Theta_r$,
Lemma~\ref{lem:affine-rank-criterion} now gives
\[
 R_c(N,2N+1)=\widetilde R_c(0,-N)=0.
\]
Consequently, $1-N\leq c(i)\leq 2N$ for $1\leq i\leq N$.
In other words, all the entries of the window of $c$ belong to the integer interval $[1-N,2N]$.


Since $c\in{}^IW$, the converse direction of Lemma~\ref{lem:words}
now gives a unique admissible word $\omega$ such that $c=c_\omega$.
Its type $t$ satisfies $0\leq t\leq r$ by Definition~\ref{def:words}.
\end{proof}

\begin{lemma}\label{lem:word-properties}
Let $\omega=\omega_1\cdots\omega_N$ be an admissible word.
For $1\leq i<N$,
\begin{equation}\label{eq:word-descents}
 s_i\in R_I(c_\omega)
 \quad\Longleftrightarrow\quad
 \omega_i\leq\omega_{i+1}\qquad(\CC<\B<\A).
\end{equation}
If $\eta = \eta_1 \cdots \eta_N$ is another admissible word with $c_\eta\leq c_\omega$,
then, for $1\leq m\leq N$,
\begin{equation}\label{eq:word-prefix-bound}
 \#\{i\leq m:\eta_i=\CC\}
 \leq\#\{i\leq m:\omega_i=\CC\}.
\end{equation}
\end{lemma}
\begin{proof}
We first prove \eqref{eq:word-descents}. Fix $1\leq i<N$.
If $\omega_i\ne\omega_{i+1}$, let $\eta$ be obtained by interchanging
these two letters. Then $c_\eta=c_\omega s_i$, and
Lemma~\ref{lem:word-length} gives
\[
 \ell(c_\eta)-\ell(c_\omega)
 =\iota(\omega)-\iota(\eta)
 =\begin{cases}
   -1,&\omega_i<\omega_{i+1},\\
    1,&\omega_i>\omega_{i+1},
  \end{cases}
 \qquad(\CC<\B<\A).
\]
Since $c_\omega,c_\eta\in{}^I W$, length additivity after left
multiplication by $w_I$ shows that $s_i\in R_I(c_\omega)$ exactly
when $\omega_i<\omega_{i+1}$.

If $\omega_i=\omega_{i+1}$, the corresponding window entries are
consecutive entries of the same list in \eqref{eq:lists}. They lie
in an interval $[qN+1,(q+1)N]$ for some $q\in\mathbb Z$, so
$c_\omega s_i=s_jc_\omega$ for some $s_j\in I$. Hence
\[
 \ell(w_Ic_\omega s_i)
 =\ell(w_Is_j)+\ell(c_\omega)
 =\ell(w_Ic_\omega)-1.
\]
Thus $s_i\in R_I(c_\omega)$ also when the letters are equal, proving
\eqref{eq:word-descents}.

Finally, let $\omega$ be an admissible word of type $t$. Its largest
window entry is $N+t$, so periodicity gives
$c_\omega(i)\leq t<N+1$ for every $i\leq0$. Hence, for
$1\leq m\leq N$,
\[
 R_{c_\omega}(m,N+1)
 =\#\{1\leq i\leq m:\omega_i=\CC\}.
\]
Applying this identity to $\eta$ and $\omega$, the inequality
$c_\eta\leq c_\omega$ and
Lemma~\ref{lem:affine-rank-criterion} give
\eqref{eq:word-prefix-bound}.
\end{proof}

The first-chain candidates include the following family.

\begin{definition}\label{def:kappa}
For $1\leq j\leq r$, let $\kappa_j=c_\omega\in{}^IW$ for the admissible
word of type $j$
\[
 \omega=\A\CC^j\B^{N-2j}\A^{j-1}.
\]
Its window is
\[
 \kappa_j=[1-j,N+1,\ldots,N+j,j+1,\ldots,N-j,2-j,\ldots,0].
\]
\end{definition}

{\color{black}
We record four properties of $\kappa_1$ that will be used to prove the vanishing of the $\kappa_j$ corrections in \Cref{lem:kappa-vanishing}.
Recall the element $z(1,1,0)$ from \Cref{def:final-maxima}.

\begin{lemma}
\label{lem:kappa-basic-properties}
For $r\geq2$, we have
\[
\kappa_1s_2=z(1,1,0)\in{}^IW,\qquad
\kappa_1s_2<\kappa_1,\qquad
\kappa_1\not\leq X_0,\qquad
\kappa_1s_i\notin{}^I W\quad(3\leq i\leq r).
\]
\end{lemma}

\begin{proof}
Right multiplication by $s_2$ exchanges the second and third
entries of the window of $\kappa_1$, giving
\[
\kappa_1s_2=[0,2,N+1,3,\ldots,N-1]=z(1,1,0)\in{}^IW.
\]
Since $\kappa_1(2)>\kappa_1(3)$, we get $\kappa_1s_2<\kappa_1$.
The admissible words of $\kappa_1$ and $X_0$ begin with $\A\CC$ and $\A\B$, respectively, so the Bruhat prefix bound \eqref{eq:word-prefix-bound} at $m=2$ gives $\kappa_1\not\leq X_0$.
Finally, right multiplication by $s_i$ exchanges the entries in
positions $i,i+1$, which in the window of $\kappa_1$ are $i-1,i$.
This is also the effect of left multiplication by $s_{i-1}$,
which exchanges these two values. Hence
$\kappa_1s_i=s_{i-1}\kappa_1$ for $3\leq i\leq r$.
Since $s_{i-1}\in I$ and $\kappa_1\in{}^I W$,
it follows that $\kappa_1s_i\notin{}^I W$.
\end{proof}
}

\begin{lemma}\label{lem:candidates}
Assume $r\geq2$.
\begin{enumerate}[label=\textup{(\roman*)}]
\item For $0\leq k\leq r-2$, the non-cover candidates for the
first-chain step $x_k\to x_{k+1}$ satisfy
\[
 \mathcal C_{\mathrm{nc}}(x_k,s_{k+2})
 \subseteq \{e\}\cup\{\gamma_j,\kappa_j:1\leq j\leq k\}.
\]
\item For $0\leq t\leq r-2$, put $b=r-1-t$.
Every element of $\mathcal C_{\mathrm{nc}}(y_{t+1},s_b)$
has the form $c_\omega$, where $\omega=\omega_1\cdots\omega_N$
is an admissible word satisfying
\begin{equation}\label{eq:outer-block-order}
 \omega_1\leq\cdots\leq\omega_{r+1},\qquad
 \omega_{r+2}\leq\cdots\leq\omega_N,
\end{equation}
and
\begin{equation}\label{eq:outer-prefix-letter}
 \omega_b\ne\CC,
\end{equation}
where the letters are ordered by $\CC<\B<\A$.
\end{enumerate}
\end{lemma}
\begin{proof}
For either chain, let $u$ be a non-cover candidate. Then
$u\in{}^I W$ and $u$ lies below the corresponding chain element.
The chain words are
\[
 x_k:\ \A\CC^k\B\CC^{r-k}\A^{r-1},\qquad
 y_{t+1}:\ \CC^{b-1}\A\CC^{t+1}\B\CC\A^{r-1},
 \quad b=r-1-t.
\]
We arrange the letters in each displayed word in the order
$\CC<\B<\A$ by repeatedly applying the adjacent interchanges
\[
 \A\B\longmapsto\B\A,\qquad
 \A\CC\longmapsto\CC\A,\qquad
 \B\CC\longmapsto\CC\B.
\]
By Lemma~\ref{lem:word-properties} and length additivity for minimal
representatives, each of these interchanges gives a Bruhat cover.
Both displayed words have type $r$, so the process ends at
$\CC^r\B\A^r$, the word of $\Theta_r$. Hence
$x_k\leq\Theta_r$ and $y_{t+1}\leq\Theta_r$.
It follows that $u\leq\Theta_r$, and
Lemma~\ref{lem:word-ideal} gives an admissible word $\omega$ with
$u=c_\omega$. We now apply Lemmas~\ref{lem:descents} and
\ref{lem:word-properties} to restrict this word.

\smallskip\noindent
\textup{(i)} Applying \eqref{eq:word-descents} to the word for $x_k$
displayed above shows that the finite generators missing from
$R_I(x_k)$ are $s_1$ and $s_{k+2}$. 
Lemma~\ref{lem:descents} gives
$R_I(x_k)\subseteq R_I(c_\omega)$, and the candidate condition adds
$s_{k+2}\in R_I(c_\omega)$. 
We therefore have $I\setminus\{s_1\}\subseteq R_I(c_\omega)$.
Equation~\eqref{eq:word-descents} gives
\[
 \omega_2\leq\omega_3\leq\cdots\leq\omega_N
 \qquad(\CC<\B<\A).
\]
These descent conditions leave the first letter $\omega_1$ unrestricted.
The suffix $\omega_2\cdots\omega_N$ consists of three (possibly empty)
blocks of $\CC$, $\B$, and $\A$, in that order.

Let $j$ be the type of $\omega$.
If $j=0$, then $\omega=\B^N$ and $c_\omega=e$.

Suppose $j>0$. 
We notice that first letter and the prescribed letter counts
determine the sorted suffix, giving the following three cases.

If $\omega_1=\CC$, the suffix has $j-1$ letters $\CC$,
$N-2j$ letters $\B$, and $j$ letters $\A$. Thus
\[
 \omega=\CC^j\B^{N-2j}\A^j,\qquad c_\omega=\Theta_j.
\]
Thus $u=c_\omega=\Theta_j$. Since $u$ is a candidate for the step
$x_k\to x_{k+1}$, we must also have $u<x_k$.
Applying the prefix bound \eqref{eq:word-prefix-bound} to
$\Theta_j<x_k$ at $m=1$ would give $1\leq0$: the word of
$\Theta_j$ begins with $\CC$, whereas the word of $x_k$ begins
with $\A$. This contradiction excludes the case $\omega_1=\CC$.

If $\omega_1=\B$, the suffix has $j$ letters $\CC$,
$N-2j-1$ letters $\B$, and $j$ letters $\A$. Hence
\[
 \omega=\B\CC^j\B^{N-2j-1}\A^j,\qquad c_\omega=\gamma_j.
\]

If $\omega_1=\A$, the suffix has $j$ letters $\CC$,
$N-2j$ letters $\B$, and $j-1$ letters $\A$. Hence
\[
 \omega=\A\CC^j\B^{N-2j}\A^{j-1},\qquad c_\omega=\kappa_j.
\]

For either of the last two words, the first $k+2$ positions contain
$\min(j,k+1)$ letters $\CC$, whereas the word of $x_k$ has $k$
such letters. Applying \eqref{eq:word-prefix-bound} at $m=k+2$ gives
\[
 \min(j,k+1)\leq k,
\]
and therefore $j\leq k$.
This proves \textup{(i)}.

\smallskip\noindent
\textup{(ii)} Applying \eqref{eq:word-descents} to the word for
$y_{t+1}$ displayed above shows that the finite generators missing
from $R_I(y_{t+1})$ are $s_b,s_{r+1}$. Lemma~\ref{lem:descents} gives
$R_I(y_{t+1})\subseteq R_I(c_\omega)$, and the candidate condition
adds $s_b\in R_I(c_\omega)$.
Equation~\eqref{eq:word-descents} therefore gives
\eqref{eq:outer-block-order}.
The word of $y_{t+1}$ begins with $\CC^{b-1}\A$, so the prefix
bound \eqref{eq:word-prefix-bound} at $m=b$ gives
\[
 \#\{1\leq i\leq b:\omega_i=\CC\}\leq b-1.
\]
If $\omega_b=\CC$, the first ordered block would begin with at
least $b$ letters $\CC$, contradicting this inequality.
This proves \eqref{eq:outer-prefix-letter} and completes
\textup{(ii)}.
\end{proof}

\subsection{Covers contribute only the preceding chain element}
We first record the local changes that can produce a cover between
minimal representatives. 

For integers $p<q$ with $p\not \equiv q \mod N$ we write $t_{p,q}\in W$ for the transposition that exchanges positions $p+kN$ and $q+kN$ for every $k \in \Z$. 

\begin{lemma}\label{lem:word-covers}
Let $x=c_\omega\leq\Theta_r$, and suppose that
$z\in{}^I W$ satisfies $z\triangleleft x$.
Then $z=c_{\omega'}$ for an admissible word $\omega'$ obtained from
$\omega$ by one of the following operations:
\begin{enumerate}[label=\textup{(\roman*)}]
\item Interchange $\CC$ at position $p$ with $\B$ at $q>p$,
with only letters $\A$ between them.
\item Interchange $\B$ at $p$ with $\A$ at $q>p$,
with only letters $\CC$ between them.
\item Interchange an adjacent pair $\CC\A$.
\item Replace the first $\A$ and the last $\CC$ of $\omega$ by $\B$,
provided the first $\A$ precedes the last $\CC$ and every occurrence
of $\B$ lies strictly between them.
\end{enumerate}
\end{lemma}
\begin{proof}
By Lemma~\ref{lem:word-ideal}, we have $z=c_{\omega'}$ for an
admissible word $\omega'$.
Since $z\triangleleft x$, we can write $z=xt_{p,q}$ for integers
$p<q$ with $1\leq p\leq N$ and $p\not\equiv q\pmod N$.
The cover criterion in \cite[Lemma~2.2(2)]{LLMS} gives
\begin{equation}\label{eq:affine-cover-test}
 \begin{gathered}
 x(p)>x(q),\\
 x(i)\notin(x(q),x(p))\quad\text{for every integer }p<i<q.
 \end{gathered}
\end{equation}
Write $q=q_0+dN$ with $1\leq q_0\leq N$.
Since $1\leq p\leq N$ and $p<q$, we have $d\geq0$.
Let $0\leq t\leq r$ be the type of $\omega$.
The window entries of $x$ lie between $1-t$ and $N+t$, and
$2t\leq2r<N$. Thus
\begin{equation}
 dN<x(p)-x(q_0)\leq(N+t)-(1-t)=N+2t-1<2N.
\end{equation}
We conclude that $d=0$ or $d=1$.

If $d=0$, both positions $p$ and $q$ lie in the window.
The exchange preserves the set of window entries, so $\omega'$ also has type $t$.
Since $x(p)>x(q)$, the letters $\omega_p$ and $\omega_q$ must be
distinct and satisfy $\omega_p<\omega_q$ in the letter order
$\CC<\B<\A$.
An occurrence of either letter $\omega_p$ or $\omega_q$ strictly
between $p$ and $q$ would put entries of that list out of increasing
order after the exchange, contradicting the admissibility of
$\omega'$. Thus only the third letter can occur between positions
$p$ and $q$.
Consequently, the subword $\omega_p\cdots\omega_q$ must be one of
$\CC\A^k\B$, $\B\CC^k\A$, or $\CC\B^k\A$, for some $k\geq0$.
The first two options are \textup{(i)} and \textup{(ii)}, respectively.
The option $\CC\B^k\A$ with $k\geq1$ is excluded by
\eqref{eq:affine-cover-test}: any occurrence of $\B$ at a position
$p<i<q$ gives $x(p)>x(i)>x(q)$.
The case $k=0$ corresponds to \textup{(iii)}.

If $d=1$, then $q=q_0+N$. By periodicity, for $1\leq i\leq N$ we have
\begin{equation}\label{eq:cover-shifted-window}
 z(i)=
 \begin{cases}
 x(i),&\text{if }i\notin\{p,q_0\},\\
 x(q_0)+N,&\text{if }i=p,\\
 x(p)-N,&\text{if }i=q_0.
 \end{cases}
\end{equation}
Since $x(p)>x(q)=x(q_0)+N$, we have $x(p)-x(q_0)>N$.
The ranges in \eqref{eq:lists} therefore force
$\omega_p=\CC$ and $\omega_{q_0}=\A$.

It follows from \eqref{eq:cover-shifted-window} that
\[
 z(q_0)=x(p)-N\in\{1,\ldots,t\},\qquad
 z(p)=x(q_0)+N\in\{N+1-t,\ldots,N\}.
\]
Thus the $\A$-entry at $q_0$ and the $\CC$-entry at $p$ are both
replaced by entries in $\{1,\ldots,N\}$.
All other window entries remain unchanged, so $z$ has exactly
$t-1$ entries at most $0$ and $t-1$ entries greater than $N$.
Since $z=c_{\omega'}$ and $\omega'$ is admissible, its type is $t-1$.
Its $\A$-entries are therefore $2-t,\ldots,0$, and its
$\CC$-entries are $N+1,\ldots,N+t-1$.
Compared with the lists for $\omega$, the missing entries are
precisely $1-t$ and $N+t$. Hence
\begin{equation}\label{eq: values of x}
   x(q_0)=1-t,\qquad x(p)=N+t.  
\end{equation}

Since each list occurs in increasing order in $x$, position $q_0$
contains the first $\A$ of $\omega$, and position $p$ contains its
last $\CC$. By \eqref{eq:cover-shifted-window} and
\eqref{eq: values of x}, the two new entries are
\[
 z(q_0)=t,\qquad z(p)=N+1-t.
\]
These are the smallest and largest entries of the $\B$-list for
type $t-1$. The original $\B$-entries, with values
$t+1,\ldots,N-t$, remain in their original positions.
Since the $\B$-entries of $z$ occur in increasing order, we obtain
\[
 q_0<i<p\qquad\text{whenever }\omega_i=\B.
\]
Since $\omega$ contains $N-2t\geq1$ occurrences of $\B$, we also
have $q_0<p$. Thus every occurrence of $\B$ in $\omega$ lies strictly
between its first $\A$, at $q_0$, and its last $\CC$, at $p$.
Replacing these two letters by $\B$ gives $\omega'$, as in
\textup{(iv)}.
\end{proof}

\begin{example}\label{ex:word-covers}
Let $N=7$, $r=3$, and consider the admissible word of type $3$
\[
 \omega=\CC\A\CC\A\B\CC\A.
\]
The three lists are $A_3=(-2,-1,0)$, $B_3=(4)$ and $C_3=(8,9,10)$,
so
\[
 x=c_\omega=[8,-2,9,-1,4,10,0].
\]
Sorting the letters into the order $\CC<\B<\A$ by adjacent
interchanges increases length at each step and ends at
$\CC^3\B\A^3$, the word of $\Theta_3$. Thus $x\leq\Theta_3$.
The following table gives one cover of each type in
Lemma~\ref{lem:word-covers}; the positions refer to the original word
$\omega$.
\begin{center}\small
\begin{tabular}{@{}ccll@{}}
\toprule
Case & Positions & $\omega'$ & $z=c_{\omega'}$ \\
\midrule
\textup{(i)} & $3,5$ & $\CC\A\B\A\CC\CC\A$
 & $[8,-2,4,-1,9,10,0]$ \\
\textup{(ii)} & $5,7$ & $\CC\A\CC\A\A\CC\B$
 & $[8,-2,9,-1,0,10,4]$ \\
\textup{(iii)} & $3,4$ & $\CC\A\A\CC\B\CC\A$
 & $[8,-2,-1,9,4,10,0]$ \\
\textup{(iv)} & $2,6$ & $\CC\B\CC\A\B\B\A$
 & $[8,3,9,-1,4,5,0]$ \\
\bottomrule
\end{tabular}
\end{center}
The reflections in the first three rows are $t_{3,5}$, $t_{5,7}$
and $s_3$, respectively. In the fourth, the only $\B$ lies between
the first $\A$ and the last $\CC$, at positions $2$ and $6$.
Here $z=xt_{6,9}$, with $z(2)=x(6)-7=3$ and
$z(6)=x(2)+7=5$, and the type drops from $3$ to $2$.

The inversion numbers of $\omega$ and the four resulting words are
$6$ and $7,7,7,3$, respectively. Using their types,
Lemma~\ref{lem:word-length} gives $\ell(x)=18$ and $\ell(z)=17$
in every row. Thus all four satisfy $z\triangleleft x$.
\end{example}

\begin{lemma}\label{lem:covers}
Assume $r\geq2$.
\begin{enumerate}[label=\textup{(\roman*)}]
\item For $0\leq k\leq r-2$, the cover candidates for the first
chain are
\[
 \mathcal C_{\mathrm{cov}}(x_k,s_{k+2})=
 \begin{cases}
  \varnothing,&k=0,\\
  \{x_{k-1}\},&1\leq k\leq r-2.
 \end{cases}
\]
\item For $0\leq t\leq r-2$, the cover candidates for the second
chain are
\[
 \mathcal C_{\mathrm{cov}}(y_{t+1},s_{r-1-t})=
 \begin{cases}
  \{y_{t+2}\},&0\leq t\leq r-3,\\
  \varnothing,&t=r-2.
 \end{cases}
\]
\end{enumerate}
Every candidate displayed here contributes with coefficient one.
\end{lemma}
\begin{proof}
For a step $u\to us_j>u$ in either chain, a cover
$c=c_\eta\triangleleft u$ belongs to
$\mathcal C_{\mathrm{cov}}(u,s_j)$ exactly when $s_j\in R_I(c)$:
its coefficient $\mus(c,u)$ is already one.
By the descent test \eqref{eq:word-descents},
\begin{equation}\label{eq:cover-descent-pairs}
 s_j\in R_I(c_\eta)
 \quad\Longleftrightarrow\quad
 \eta_j\eta_{j+1}\in
 \{\A\A,\B\B,\CC\CC,\CC\B,\CC\A,\B\A\}.
\end{equation}
Thus equal letters and the pairs $\CC\B$, $\CC\A$, and $\B\A$
give the required descent, while the pairs $\A\B$, $\A\CC$, and $\B\CC$
do not. The proof of Lemma~\ref{lem:candidates} shows that the
chain elements lie below $\Theta_r$. We may therefore apply
Lemma~\ref{lem:word-covers} to their covers. We treat its four
cases for each chain and test the resulting consecutive letters.

\smallskip\noindent
\textup{(i)} The step is $x_k\to x_{k+1}=x_ks_{k+2}$.
The word of $x_k$ is
\[
 \A\CC^k\B\CC^{r-k}\A^{r-1}.
\]
Its letters at positions $k+2,k+3$ are $\B\CC$; the step
interchanges them to $\CC\B$. For a lower cover $c_\eta$ to
contribute to this multiplication, its letters at those same
positions must satisfy \eqref{eq:cover-descent-pairs}.

There is only one $\B$. In cases \textup{(i)} and \textup{(ii)}
of Lemma~\ref{lem:word-covers}, it is exchanged with the $\CC$
immediately to its left, when $k\geq1$, or with the first $\A$
to its right, respectively. The only adjacent $\CC\A$ is at
positions $r+2,r+3$, and case \textup{(iv)} replaces the letters
at positions $1,r+2$. Hence every possible cover word is listed
below. The last two columns record its pair at $k+2,k+3$ and
whether it gives the required right descent.
\begin{center}
\begin{tabular}{@{}clcc@{}}
\toprule
Case & Word $\eta$ & Pair & $s_{k+2}\in R_I(c_\eta)$ \\
\midrule
\textup{(i)}, $k\geq1$
 & $\A\CC^{k-1}\B\CC^{r-k+1}\A^{r-1}$ & $\CC\CC$ & yes \\
\textup{(ii)}
 & $\A\CC^k\A\CC^{r-k}\B\A^{r-2}$ & $\A\CC$ & no \\
\textup{(iii)}
 & $\A\CC^k\B\CC^{r-k-1}\A\CC\A^{r-2}$ & $\B\CC$ & no \\
\textup{(iv)}
 & $\B\CC^k\B\CC^{r-k-1}\B\A^{r-1}$ & $\B\CC$ & no \\
\bottomrule
\end{tabular}
\end{center}
Here $k\leq r-2$ gives $r-k-1\geq1$, including at the boundary.
Only the first row satisfies the descent test, and its word is
that of $x_{k-1}$. This element is indeed a cover:
\[
 x_{k-1}=x_ks_{k+1}\triangleleft x_k\qquad(k\geq1),
\]
since the adjacent interchange $\CC\B\mapsto\B\CC$ decreases
length by one, by Lemma~\ref{lem:word-length}.
When $k=0$, the first row is absent and there are no cover candidates.
This proves \textup{(i)}.

\smallskip\noindent
\textup{(ii)} Put $b=r-1-t$, so $1\leq b\leq r-1$.
The step is $y_{t+1}\to y_t=y_{t+1}s_b$, and the word of
$y_{t+1}$ is
\[
 \CC^{b-1}\A\CC^{r-b}\B\CC\A^{r-1}.
\]
The letters at positions $b,b+1$ are $\A\CC$; the step
interchanges them to $\CC\A$. A cover contributes precisely
when its letters at $b,b+1$ satisfy
\eqref{eq:cover-descent-pairs}.

In cases \textup{(i)} and \textup{(ii)} of
Lemma~\ref{lem:word-covers}, the exchanged positions are
$r,r+1$ and $r+1,r+3$, respectively.
In case \textup{(iii)}, the adjacent $\CC\A$ pairs are at
$b-1,b$, when $b\geq2$, and at $r+2,r+3$.
Case \textup{(iv)} replaces the letters at $b,r+2$.
The possible cover words and their pairs at $b,b+1$ are therefore
as follows.
\begin{center}
\begin{tabular}{@{}clcc@{}}
\toprule
Case & Word $\eta$ & Pair & $s_b\in R_I(c_\eta)$ \\
\midrule
\textup{(i)}
 & $\CC^{b-1}\A\CC^{r-b-1}\B\CC^2\A^{r-1}$
 & $\A\CC$ or $\A\B$ & no \\
\textup{(ii)}
 & $\CC^{b-1}\A\CC^{r-b}\A\CC\B\A^{r-2}$ & $\A\CC$ & no \\
\textup{(iii)}, $b\geq2$
 & $\CC^{b-2}\A\CC^{r-b+1}\B\CC\A^{r-1}$ & $\CC\CC$ & yes \\
\textup{(iii)}
 & $\CC^{b-1}\A\CC^{r-b}\B\A\CC\A^{r-2}$ & $\A\CC$ & no \\
\textup{(iv)}
 & $\CC^{b-1}\B\CC^{r-b}\B^2\A^{r-1}$ & $\B\CC$ & no \\
\bottomrule
\end{tabular}
\end{center}
In the first row, the pair is $\A\CC$ when $b\leq r-2$ and
$\A\B$ when $b=r-1$; neither gives a descent.
The two rows from case \textup{(iii)} correspond, in order,
to the adjacent pairs at $b-1,b$ and $r+2,r+3$.
Only the first of these rows passes the test. Its word is that
of $y_{t+2}$, and
\[
 y_{t+2}=y_{t+1}s_{b-1}\triangleleft y_{t+1}\qquad(b\geq2),
\]
since the adjacent interchange $\CC\A\mapsto\A\CC$ decreases
length by one. This row exists exactly when $b\geq2$, or
equivalently $t\leq r-3$. When $b=1$, all remaining pairs fail
the test, so there are no cover candidates. This proves
\textup{(ii)}.
\end{proof}

\subsection{The multiplication formulas at this stage}
\label{sec:candidate-expansions}
Lemmas~\ref{lem:candidates} and \ref{lem:covers} locate every
correction that can be nonzero in \eqref{eq:KLrule}. The cover
contributions are now completely determined, each with coefficient one.
For the non-covers, however, we have only necessary restrictions on the
indices: an element satisfying these restrictions may still have
coefficient $\mus$ equal to zero.

For $k=0$, the first-chain formula \eqref{eq:KLrule} becomes
\begin{equation}\label{eq:first-chain-initial-candidate-expansion}
 \C_{x_0}C_{s_2}=\C_{x_1}+\mus(e,x_0)\C_e.
\end{equation}
For $1\leq k\leq r-2$, it becomes
\begin{equation}\label{eq:first-chain-candidate-expansion}
 \C_{x_k}C_{s_{k+2}}
 =\C_{x_{k+1}}+\C_{x_{k-1}}
 +\mus(e,x_k)\C_e
 +\sum_{j=1}^{k}\bigl(
   \mus(\gamma_j,x_k)\C_{\gamma_j}
   +\mus(\kappa_j,x_k)\C_{\kappa_j}\bigr).
\end{equation}
The coefficient-one term in the second formula is the cover
contribution; all terms multiplied by $\mus$ are possible non-cover
contributions.

For the second chain, fix $0\leq t\leq r-2$, put $b=r-1-t$, and define
\begin{equation}\label{eq:remaining-second-chain-indices}
 \mathcal N_t=
 \left\{c_\omega:
 \begin{array}{l}
  \omega\text{ is admissible and satisfies }
  \eqref{eq:outer-block-order}\text{ and }\eqref{eq:outer-prefix-letter},\\
  c_\omega<y_{t+1},\quad
  \ell(y_{t+1})-\ell(c_\omega)>1
 \end{array}
 \right\}.
\end{equation}
The order within the first block gives $s_b\in R_I(c)$ for every
$c\in\mathcal N_t$, by \eqref{eq:word-descents}.
Thus \eqref{eq:KLrule} becomes
\begin{equation}\label{eq:second-chain-candidate-expansion}
\begin{aligned}
 \C_{y_{t+1}}C_{s_{r-1-t}}
 ={}&\C_{y_t}
 +\begin{cases}
\C_{y_{t+2}},&0\leq t\leq r-3,\\
  0,&t=r-2,
 \end{cases}\\
 &+\sum_{c\in\mathcal N_t}\mus(c,y_{t+1})\C_c.
\end{aligned}
\end{equation}
Here $\mathcal C_{\mathrm{nc}}(y_{t+1},s_{r-1-t})\subseteq
\mathcal N_t$, but membership in $\mathcal N_t$ does not assert
that the corresponding coefficient is nonzero.

% The cover coefficients in these identities are already known.
% Section~\ref{sec:first-multiplication-proof} will exclude the
% $\kappa_j$ and determine the $\gamma_j$ multiplicities while computing
% the successive scalar coefficients. For the second chain, positivity
% will bound the pure contributions, and a rank-two argument will exclude
% the mixed ones. The cancellation result below supplies the remaining
% coefficient identities. Only the uniform initial estimate and the
% first-chain calculation are needed; no induction on the rank is required.







\subsection{Cancellation and reduction to smaller rank}

We now use affine cancellation to express certain correction coefficients
in terms of Kazhdan--Lusztig polynomials in smaller rank. 






\begin{notation}
For $r\geq2$, a superscript $(r)$ indicates that the symbol is taken
with rank parameter $r$, so that $N=2r+1$. In particular, $W^{(r)}$
is the affine Weyl group of type $\widetilde A_{2r}$, realized by
affine permutations of period $2r+1$, whose windows have $2r+1$ entries.
The elements $x_k^{(r)}$, $\gamma_j^{(r)}$, $\kappa_j^{(r)}$ and
$\Theta_j^{(r)}$ are defined by the same formulas as before, with
this value of $N$. 

The same convention applies to the associated basis elements and
coefficients. Thus
\[
 h_{x_k}^{(r)}=[M_e^{(r)}]\C_{x_k}^{(r)},
\]
and $[M_{\gamma_j}^{(r)}]\C_{x_k}^{(r)}$ is computed in the spherical
module for $W^{(r)}$, with indices $\gamma_j^{(r)}$ and $x_k^{(r)}$.
In particular, $h_{x_{k-j}}^{(r-j)}$ uses the index
$x_{k-j}^{(r-j)}$ in period $2(r-j)+1$.
We omit the superscript when the rank parameter is clear.
\end{notation}

\begin{example}\label{ex:rank-notation}
We illustrate  the elements $x_1$ and $\gamma_1$ in rank parameters $2$ and $3$:
\begin{center}
\begin{tabular}{@{}clc@{}}
\toprule
Index & Admissible word & Window\\\midrule
$x_1^{(2)}$ & $\mathsf{ACBCA}$ & $[-1,6,3,7,0]$\\
$x_1^{(3)}$ & $\mathsf{ACBCCAA}$ & $[-2,8,4,9,10,-1,0]$\\
$\gamma_1^{(2)}$ & $\mathsf{BCBBA}$ & $[2,6,3,4,0]$\\
$\gamma_1^{(3)}$ & $\mathsf{BCBBBBA}$ & $[2,8,3,4,5,6,0]$\\\bottomrule
\end{tabular}
\end{center}

\end{example}

\begin{definition}\label{def:paired-deletion}
Let $N\geq5$, and let $w$ be an affine permutation of period $N$
satisfying
\[
 w(a)=2N,\qquad w(N)=1-N
 \qquad\text{for some }1\leq a<N.
\]
We define $\operatorname{del}_a(w)$ by deleting the position classes
$a+N\mathbb Z$ and $N+N\mathbb Z$, together with their image classes (that are $1+N\Z$ and $N+N\Z$),
and relabelling the remaining positions and values increasingly.

Explicitly, list the surviving window positions as
\[
 \{p_1<\cdots<p_{N-2}\}=\{1,\ldots,N\}\setminus\{a,N\},
\]
and write $w(p_i)=b_i+q_iN$, where $2\leq b_i\leq N-1$ and
$q_i\in\mathbb Z$. Set
\[
 \operatorname{del}_a(w)(i)=b_i-1+q_i(N-2)
 \qquad(1\leq i\leq N-2),
\]
and extend periodically with period $N-2$.
The removed residue classes have shifts $+1$ and $-1$, so
$\sum_iq_i=0$. Thus $\operatorname{del}_a(w)$ is an affine permutation
of period $N-2$, with window sum $(N-2)(N-1)/2$.
\end{definition}

\begin{proposition}\label{prop:affine-cancel}
Let $N\geq5$, and let $Z\leq Y$ be affine permutations of period $N$.
Suppose that
\[
 1-N\leq Z(i),Y(i)\leq2N\qquad(1\leq i\leq N),
\]
and that, for some $1\leq a<N$,
\[
 Z(a)=Y(a)=2N,\qquad Z(N)=Y(N)=1-N.
\]
Then
\[
 \operatorname{del}_a(Z)\leq\operatorname{del}_a(Y),\qquad
 h_{Z,Y}=h_{\operatorname{del}_a(Z),\operatorname{del}_a(Y)}.
\]
\end{proposition}
\begin{proof}
This follows by two applications of 
\cite[Proposition~4.2]{Henderson}. 
\end{proof}

\begin{example}\label{ex:affine-cancellation}
For $\gamma_1^{(3)}$ and $x_1^{(3)}$, consider the ordinary endpoints
\[
 \begin{aligned}
 Z=w_I^{(3)}\gamma_1^{(3)}&=[6,\mathbf{14},5,4,3,2,\mathbf{-6}],\\
 Y=w_I^{(3)}x_1^{(3)}&=[-4,\mathbf{14},4,13,12,-5,\mathbf{-6}].
 \end{aligned}
\]
The rank criterion gives $Z\leq Y$. Their windows lie in $[-6,14]$
and share the extreme entries $14$ and $-6$ at positions $2$ and $7$.
Thus Proposition~\ref{prop:affine-cancel} applies with $N=7$ and $a=2$.
The surviving positions are $1,3,4,5,6$, and
Definition~\ref{def:paired-deletion} gives
\[
 \begin{aligned}
 \operatorname{del}_2(Z)&=[5,4,3,2,1]=w_I^{(2)},\\
 \operatorname{del}_2(Y)&=[-3,3,10,9,-4]=w_I^{(2)}x_0^{(2)}.
 \end{aligned}
\]
At the word level, deleting the second and last letters gives
\[
 \mathsf{BCBBBBA}\longmapsto\mathsf{BBBBB},\qquad
 \mathsf{ACBCCAA}\longmapsto\mathsf{ABCCA}.
\]
Consequently,
\[
 [M_{\gamma_1}^{(3)}]\C_{x_1}^{(3)}
 =h_{Z,Y}
 =h_{\operatorname{del}_2(Z),\operatorname{del}_2(Y)}
 =h_{x_0}^{(2)}.
\]
Thus a coefficient indexed by $\gamma_1^{(3)}$ in period $7$ becomes
an identity coefficient in period $5$, with upper element $x_0^{(2)}$.
This is the case $r=3$, $j=k=1$ of the first-chain reduction in
Lemma~\ref{lem:rank-reduction} below.
\end{example}

We apply cancellation to common extreme entries. For
$w\in{}^I W$ associated to an admissible word, the entries of $w_Iw$
lie in $[1-N,2N]$. The following reductions therefore require only
two common positions at each stage.



\begin{lemma}\label{lem:rank-reduction}
Assume $r\geq3$. Then:
\begin{enumerate}[label=\textup{(\roman*)}]
\item For $1\leq j\leq k\leq r-2$,
\begin{equation}\label{eq:rank-reduction-first-chain}
 [M_{\gamma_j}^{(r)}]\C_{x_k}^{(r)}=h_{x_{k-j}}^{(r-j)}.
\end{equation}
\item For $0\leq t\leq r-3$ and $1\leq j\leq r-t-2$,
\begin{equation}\label{eq:rank-reduction-second-chain}
 [M_{\Theta_j}^{(r)}]\C_{y_{t+1}}^{(r)}=h_{y_{t+1}}^{(r-j)}.
\end{equation}
\item If $1\leq j\leq k\leq r-2$ and $\kappa_j<x_k$, put $m=r-j+1$.
Then
\begin{equation}\label{eq:rank-reduction-kappa}
 [M_{\kappa_j}^{(r)}]\C_{x_k}^{(r)}
 =\bigl[M_{\kappa_1}^{(m)}\bigr]\C_{x_{k-j+1}}^{(m)}.
\end{equation}
\end{enumerate}
\end{lemma}

\begin{proof}
We only prove \textup{(i)}; the other two identities are proved
in a similar fashion.
For $u,v\in{}^I W$, we have
\[
 [M_u]\C_v=h_{w_Iu,w_Iv}.
\]
Thus \eqref{eq:rank-reduction-first-chain} is equivalent to
\begin{equation}\label{eq: equiv deleting}
 h^{(r)}_{w_I^{(r)}\gamma_j^{(r)},\,w_I^{(r)}x_k^{(r)}}
 =h^{(r-j)}_{w_I^{(r-j)},\,w_I^{(r-j)}x_{k-j}^{(r-j)}}.
\end{equation}
Put $\gamma_0=e$ and set
\[
 Z=w_I^{(r)}\gamma_j^{(r)},\qquad
 Y=w_I^{(r)}x_k^{(r)}.
\]

By \Cref{def:chains,def:gamma}, the windows of $x_k^{(r)}$ and
$\gamma_j^{(r)}$ both have $N+1$ at position $2$ and $0$ at
position $N$. Since $w_I^{(r)}(N+1)=2N$ and
$w_I^{(r)}(0)=1-N$, we obtain
\[
 Z(2)=Y(2)=2N,\qquad Z(N)=Y(N)=1-N.
\]
The lists in \eqref{eq:lists} also show that both endpoint windows
lie in $[1-N,2N]$.
The admissible words of the two minimal indices are
\[
 \B\CC^j\B^{N-2j-1}\A^j,\qquad
 \A\CC^k\B\CC^{r-k}\A^{r-1}.
\]
Deleting the letters in positions $2$ and $N$ gives the words of
$\gamma_{j-1}^{(r-1)}$ and $x_{k-1}^{(r-1)}$.
Using the value relabelling in Definition~\ref{def:paired-deletion},
we therefore obtain
\[
\begin{aligned}
 \widehat Z:=\operatorname{del}_2(Z)
 &=w_I^{(r-1)}\gamma_{j-1}^{(r-1)},\\
 \widehat Y:=\operatorname{del}_2(Y)
 &=w_I^{(r-1)}x_{k-1}^{(r-1)}.
\end{aligned}
\]

The deleted periodic strands contribute equally to the rank functions
of $Z$ and $Y$, and the increasing relabellings in
Definition~\ref{def:paired-deletion} identify the remaining cutoffs.
Hence Lemma~\ref{lem:affine-rank-criterion} gives $Z\leq Y$ if and
only if $\widehat Z\leq\widehat Y$.
Iterating this comparison $j$ times reduces it to
\[
 w_I^{(r-j)}=w_I^{(r-j)}\gamma_0^{(r-j)}
 \leq w_I^{(r-j)}x_{k-j}^{(r-j)},
\]
which follows from length additivity for the minimal representative
$x_{k-j}^{(r-j)}$. The parameter bounds ensure that all intermediate
indices are defined and $r-j\geq2$.

We can now apply Proposition~\ref{prop:affine-cancel} to obtain
\begin{equation}
 h_{Z,Y}^{(r)}=h_{\widehat Z,\widehat Y}^{(r-1)}.
\end{equation}
For $j=1$, this proves \eqref{eq: equiv deleting}, since
$\gamma_0=e$. For general $j$, replacing the pair $(Z,Y)$ by
$(\widehat Z,\widehat Y)$ and repeating the argument $j-1$ more
times yields \eqref{eq: equiv deleting}. This proves \textup{(i)}.
\end{proof}


\section{Proof of the first multiplication formulas}
\label{sec:first-multiplication-proof}

In this section we simultaneously prove the multiplication formulas in \Cref{prop:first-multiplication} and the coefficient estimates in \Cref{lem:scalar-comparison} for the first chain $x_0=X_0,x_1,\ldots,x_{r-1}$ from \Cref{def:chains}.
Throughout this section, we let $r\geq2$ and put $\eta=v+v^{-1}$.

Recall the elements $\gamma_j$ and $\kappa_j$ from \Cref{def:gamma,def:kappa}.
To obtain Proposition~\ref{prop:first-multiplication} from the candidate
expansions \eqref{eq:first-chain-initial-candidate-expansion}
and \eqref{eq:first-chain-candidate-expansion}, it remains to prove
the following identities:
\begin{equation}\label{eq:first-chain-required-multiplicities}
\begin{aligned}
 \mus(e,x_0)&=1,\\
 \mus(e,x_k)&=0
   &&(1\leq k\leq r-2),\\
 \mus(\kappa_j,x_k)&=0
   &&(1\leq j\leq k\leq r-2),\\
 \mus(\gamma_j,x_k)&=
 \begin{cases}
  0,&j<k,\\
  1,&j=k,
 \end{cases}
   &&(1\leq j\leq k\leq r-2).
\end{aligned}
\end{equation}
The first identity follows immediately from the initial estimate in \Cref{prop:seed}.
We prove the third identity in \Cref{lem:kappa-vanishing}, and then establish a recurrence
for the numbers $A_{a,b}$ in \Cref{lem:partition-step}.
With these two lemmas in hand, we prove \Cref{prop:first-multiplication} and \Cref{lem:scalar-comparison} simultaneously, establishing the remaining identities along the way.


\begin{lemma}\label{lem:kappa-vanishing}
For $1\leq j\leq k\leq r-2$,
\begin{equation}\label{eq:kappa-vanishing}
 \mus(\kappa_j,x_k)=0.
\end{equation}
\end{lemma}
\begin{proof}
We first prove that
\begin{equation}\label{eq:kappa-one-coefficient-bound}
\ord\bigl([M_{\kappa_1}]\C_{x_k}\bigr)
\geq N-k-2
\qquad(1\leq k\leq r-1).
\end{equation}
% We obtain the initial bound from
% \Cref{lem:kappa-seed-coefficient} and propagate it along
% the first chain using positivity.
Combining \Cref{lem:kappa-basic-properties,lem:kappa-seed-coefficient}, we get
\[
[M_{\kappa_1s_2}]\C_{X_0}=O(v^{N-3}),
\qquad
[M_{\kappa_1}]\C_{X_0}=0.
\]
Since $\kappa_1,\kappa_1s_2\in {}^IW$, the standard multiplication rules give, for every $F\in\M$,
\[
[M_{\kappa_1}](F C_{s_2})
=[M_{\kappa_1s_2}]F+v^{-1}[M_{\kappa_1}]F.
\]
Applying this to
\eqref{eq:first-chain-initial-candidate-expansion},
and using $[M_{\kappa_1}]\C_e=0$, yields
\[
[M_{\kappa_1}]\C_{x_1}
=[M_{\kappa_1s_2}]\C_{X_0}
=O(v^{N-3}).
\]

For $1\leq k\leq r-2$, \Cref{lem:kappa-basic-properties} gives $\kappa_1s_{k+2}\notin{}^I W$.
Thus \Cref{lem:coefficient-positivity} implies
\[
\ord\bigl([M_{\kappa_1}]\C_{x_{k+1}}\bigr)
\geq
\ord\bigl([M_{\kappa_1}]\C_{x_k}\bigr)-1.
\]
Iterating proves the claimed bound \eqref{eq:kappa-one-coefficient-bound}.
Since $N-k-2\geq r\geq2$, we get $\mus(\kappa_1,x_k)=0$ by \Cref{lem:mu}.

For $2\leq j\leq k\leq r-2$ suppose that $\mus(\kappa_j,x_k)\neq0$, so that $\kappa_j<x_k$.
Then part (iii) of \Cref{lem:rank-reduction} gives
\[
[M_{\kappa_j}]\C_{x_k}
=
[M_{\kappa_1}^{(m)}]\C_{x_a}^{(m)},
\qquad
m=r-j+1,\quad a=k-j+1.
\]
Here $1\leq a\leq m-2$ and $m\geq3$.
Applying the bound \eqref{eq:kappa-one-coefficient-bound} with rank parameter $m$ gives
\[
\ord
\bigl([M_{\kappa_1}^{(m)}]\C_{x_a}^{(m)}\bigr)
\geq 2m-a-1\geq m+1\geq4.
\]
The coefficient of $v$ therefore vanishes, contradicting
$\mus(\kappa_j,x_k)\neq0$.
\end{proof}

We now record the recurrence satisfied by the partition counts $A_{a,b}$ defined in \eqref{eq:A-definition}.

\begin{lemma}\label{lem:partition-step}
% For $a,b\geq0$, let $d_{a,b}$ and $A_{a,b}$ be as defined in
% Subsection~\ref{sec:first-chain}, with $A_{-1,b}=0$.
Write $\Part_{\leq a}(n)$ for the number of partitions of $n$ with at most $a$ parts.
Then $A_{0,b}=\Part(b)$, $A_{a,0}=a+1$ and
\begin{equation}\label{eq:partition-step}
 A_{a+1,b}=A_{a,b}+A_{a,b+1}-A_{a-1,b+1}           -\Part_{\leq a}(b+1).
\end{equation}
\end{lemma}
\begin{proof}
The first two identities are immediate.
The difference $d_{a,b+1}-d_{a+1,b}$ counts partitions of $a+b+1$ with exactly $a$ parts.
Subtracting one from each part of these partitions and discarding zeros gives a bijection with the partitions of $b+1$ having at most $a$ parts.
Its inverse consists of padding the partition with zeros to obtain exactly $a$ entries, then adding one to each entry.
Hence
\begin{equation}\label{eq: ddd}
  d_{a,b+1}=d_{a+1,b}+\Part_{\leq a}(b+1).   
\end{equation}
On the other hand,  $d_{a,b}=A_{a,b}-A_{a-1,b}$ by definition. 
By replacing this in \eqref{eq: ddd} and rearranging the terms we obtain  \eqref{eq:partition-step}. 
\end{proof}

We are now in position to complete the proof of \Cref{prop:first-multiplication} and \Cref{lem:scalar-comparison}.

\begin{proof}[Proof of Proposition~\ref{prop:first-multiplication} and
Lemma~\ref{lem:scalar-comparison}]
Fix $r\geq2$. We prove the multiplication formulas and the triangular
estimate \eqref{eq:pascal-triangle} simultaneously, by induction on $k$.
The scalar calculations use the initial estimate in \Cref{prop:seed}, the
partition recurrence in \Cref{lem:partition-step}, and the formula for
$h_{\gamma_j}$ in \Cref{lem:gaussian-products}.

Denote the multiplication identity at step $k$ by $(\mathsf M_k)$ and
the approximation to $h_{x_k}$ by $(\mathsf N_k)$:
\begin{equation+}\tag{\(\mathsf M_k\)}
 \C_{x_k}C_{s_{k+2}}
 =\begin{cases}
   \C_{x_{k+1}}+\C_e,&k=0,\\
   \C_{x_{k+1}}+\C_{x_{k-1}}+\C_{\gamma_k},&k>0,
  \end{cases}
 \qquad(0\leq k\leq r-2),
\end{equation+}
\begin{equation+}\tag{\(\mathsf N_k\)}
 h_{x_k}=\sum_{b=0}^{r-1-k}A_{k,b}v^{k+1+2b}
          +O(v^{2r+1-k})
 \qquad(0\leq k\leq r-1).
\end{equation+}
Starting from $(\mathsf N_0)$ we prove $(\mathsf M_0)$ and then $(\mathsf N_1)$.
At each subsequent step, $(\mathsf N_{k-1})$ and $(\mathsf N_k)$ will give $(\mathsf M_k)$; applying $\eps$ to that multiplication identity will then give $(\mathsf N_{k+1})$.

\Cref{prop:seed} and the identity $A_{0,b}=\Part(b)$ give $(\mathsf N_0)$.
This implies that $[v]h_{x_0}=1$, which in turn gives $(\mathsf M_0)$ by formula \eqref{eq:first-chain-initial-candidate-expansion}.
Applying $\eps$ and \Cref{lem:scalar} gives
\[
 h_{x_1}=\eta h_{x_0}-1, \qquad \eta=v+v^{-1}.
\]
To prove $(\mathsf N_1)$, we substitute $(\mathsf N_0)$ into
this scalar recurrence, obtaining
\[
 h_{x_1}
 =\sum_{d=0}^{r-1}\Part(d)v^{2d+2}
  +\sum_{d=0}^{r-1}\Part(d)v^{2d}-1+O(v^{2r}).
\]
The constant term $\Part(0)=1$ cancels with the $-1$.
For $0\leq b\leq r-2$, the coefficient of $v^{2+2b}$ receives $\Part(b)$ from the term with $d=b$ in the first sum and $\Part(b+1)$ from the term with $d=b+1$ in the second sum.
The error term does not contribute, since $2+2b<2r$.
By definition, $d_{0,b}=\Part(b)$ and $d_{1,b}=\Part(b+1)$.
Hence,
\[
[v^{2+2b}]h_{x_1}= \Part(b)+\Part(b+1)=d_{0,b}+d_{1,b}=A_{1,b}.
\]
Absorbing the last term $\Part(r-1)v^{2r}$ of the first sum into the error $O(v^{2r})$ gives $(\mathsf N_1)$:
\[
 h_{x_1}=\sum_{b=0}^{r-2}A_{1,b}v^{2+2b}+O(v^{2r}).
\]

Now let $1\leq k\leq r-2$ and assume $(\mathsf N_{k-1})$ and $(\mathsf N_k)$.
We first prove $(\mathsf M_k)$.
We have
\[
h_{x_{k-1}}=kv^k+O(v^{k+2}), \qquad
h_{x_k}=(k+1)v^{k+1}+O(v^{k+3}). 
\]
Applying $\eps$ to \eqref{eq:first-chain-candidate-expansion} and using \Cref{lem:kappa-vanishing} gives
\begin{equation}\label{eq:first-chain-positive-corrections}
 \eta h_{x_k}-h_{x_{k-1}}
 =h_{x_{k+1}}+\mus(e,x_k)
   +\sum_{j=1}^{k}\mus(\gamma_j,x_k)h_{\gamma_j}.
\end{equation}
The left-hand side has no terms below degree $k$.
Every summand on the right is coefficient-wise nonnegative by KL positivity, and $h_{\gamma_j}$ begins with $v^j$ by \Cref{lem:gaussian-products}.
Consequently
\[
 \mus(e,x_k)=0,\qquad
 \mus(\gamma_j,x_k)=0\quad(1\leq j<k).
\]
For $j=k$, part (i) of \Cref{lem:rank-reduction} combined with \Cref{prop:seed}, \Cref{lem:mu} and $r-k\geq 2$, yields
\begin{equation}\label{eq:first-chain-correction-multiplicities}
 \mus(\gamma_k,x_k)=[v]([M_{\gamma_k}]\C_{x_k})= [v]h_{X_0}^{(r-k)}=1.
\end{equation}
Substituting these $\mu_I$-values into \eqref{eq:first-chain-candidate-expansion} gives
\[
\C_{x_k}C_{s_{k+2}}
=\C_{x_{k+1}}+\C_{x_{k-1}}+\C_{\gamma_k},
\]
which proves $(\mathsf M_k)$.
Applying $\eps$ and \Cref{lem:scalar} to this identity gives
\begin{equation}\label{eq:inner}
 h_{x_{k+1}}
 =\eta h_{x_k}-h_{x_{k-1}}-h_{\gamma_k}.
\end{equation}

We now prove $(\mathsf N_{k+1})$ by computing \eqref{eq:inner} modulo
$v^{2r-k}$, the error order required in $(\mathsf N_{k+1})$.
The partition interpretation in \cite[Proposition~1.7.3]{StanleyEC1} implies
\[
[q^n]\left(\gb{N-1}{k}\right)=\Part_{\leq k}(n)
\qquad(0\leq n\leq N-1-k).
\]
Since $q=v^2$, \Cref{lem:gaussian-products} then shows that
\[
h_{\gamma_k} =\sum_{n=0}^{r-1-k}\Part_{\leq k}(n)v^{2n+k}+O(v^{2r-k}).
\]
Substituting these expansions into \eqref{eq:inner}, we compute the remaining coefficients.
By $(\mathsf N_k)$ and $(\mathsf N_{k-1})$, the error terms of $\eta h_{x_k}$ and $h_{x_{k-1}}$ have orders at least $2r-k$ and $2r+2-k$, respectively.
Since no terms of degree less than $k$ occur, it remains to compute $[v^i]h_{x_{k+1}}$ for $k\leq i<2r-k$.
Moreover, these coefficients vanish unless $i$ has parity $k$, by \Cref{lem:mu} together with \Cref{lem:endpoints} and \Cref{lem:endpoint-lengths}.
The coefficient in degree $k$ vanishes:
\[
[v^k]h_{x_{k+1}}
=A_{k,0}-A_{k-1,0}-\Part_{\leq k}(0)
=(k+1)-k-1=0.
\]
For $0\leq b\leq r-k-2$, we obtain
\[
[v^{k+2+2b}]h_{x_{k+1}}=
A_{k,b}+A_{k,b+1}-A_{k-1,b+1}-\Part_{\leq k}(b+1)=
A_{k+1,b},
\]
where the last equality follows from \Cref{lem:partition-step}.
This proves $(\mathsf N_{k+1})$ and completes the induction.

We have proved Proposition~\ref{prop:first-multiplication} and \cref{eq:pascal-triangle} of \Cref{lem:scalar-comparison}.
Now let $0\leq a\leq r-1$.
Since $A_{a,0}=a+1$ and $a+3\leq 2r+1-a$, $(\mathsf N_{a})$ gives
\[
 h_{x_a}=(a+1)v^{a+1}+O(v^{a+3}).
\]
In particular, $h_{x_{r-1}}=rv^r+O(v^{r+2})$, completing the proof of \Cref{lem:scalar-comparison}.
\end{proof}




\section{Proof of the second multiplication formulas}
\label{sec:second-multiplication-proof}
In this section we prove the multiplication formulas in \Cref{prop:second-multiplication} for the second chain.
Throughout, $r\geq2$.
Comparing Proposition~\ref{prop:second-multiplication} with the candidate expansion \eqref{eq:second-chain-candidate-expansion}, it suffices to prove the following order estimate:
\begin{equation}\label{eq:second-chain-required-order}
 \ord\left(\sum_{z\in\mathcal N_t}
       \mus(z,y_{t+1})h_z\right)\geq t+3
 \qquad(0\leq t\leq r-2).
\end{equation}
Fix $0\leq t\leq r-2$ and put $b=r-1-t$.
Recall from \eqref{eq:remaining-second-chain-indices} that $z\in\mathcal N_t$ if and only if $z=c_\omega$ for an admissible word $\omega$ whose first $r+1$ and last $r$ positions are separately ordered by $\CC<\B<\A$, with $\omega_b\ne\CC$, $z<y_{t+1}$, and
$\ell(y_{t+1})-\ell(z)>1$.

Given $c_\omega\in \mathcal N_t$ with $\omega$ of type $j$, let $c$ and $a$ be the numbers of $\CC$'s and $\A$'s in the first block, respectively.
We have
\[
 \omega=\CC^c\B^{r+1-c-a}\A^a\,\bigm|\,
         \CC^{j-c}\B^{r-2j+c+a}\A^{j-a},
\]
where the bar separates positions $r+1$ and $r+2$.
All exponents are nonnegative.
Since each block is ordered, inversions can only occur between the two blocks.
We distinguish the following cases:
\begin{enumerate}[label=\textup{(\roman*)}]
\item \emph{Ordered case}: $a=0$ and $c=j$. Here, $c_\omega=\Theta_j$.
\item \emph{Pure $\B\CC$ case}: $a=0$ and $c<j$. 
\item \emph{Pure $\A\B$ case}: $a>0$ and $c=j$. 
\item \emph{Mixed $\A\CC$ case}: $a>0$ and $c<j$.
\end{enumerate}
The labels $\B\CC$, $\A\B$, and $\A\CC$ indicate the pair of letters adjacent to the bar.
The terminology refers to inversions for $\CC<\B<\A$: in the ordered case there are no inversions; in each pure case, all inversions have the same type, $\B\CC$ or $\A\B$, respectively; and in the mixed case, $\A\CC$ inversions occur together with at least one of the other two types, since $N-2j\geq1$.
In all four cases, the condition $\omega_b\ne\CC$ is equivalent to $c\leq b-1=r-t-2$.

For $z\in\mathcal N_t$, the coefficients $\mus(z,y_{t+1})$ may still vanish.
We prove that they do vanish in the ordered and mixed cases in \Cref{lem:second-coarse,lem:mixed}, respectively.
In the pure cases, we prove that $\ord(h_z)\geq t+3$ in \Cref{lem:pure}.

\begin{lemma}\label{lem:second-coarse}
Let $z\in \mathcal N_t$ belong to the ordered case.
Then $\mus(z,y_{t+1})=0$.
\end{lemma}
\begin{proof}
Let $0\leq s\leq r-1$.
By \Cref{lem:scalar-comparison}, we have $\ord(h_{x_{r-1}})=r$.
The factorization $y_s = x_{r-1}(s_1\cdots s_{r-1-s})$ from \eqref{eq:second-chain-product} gives a path from
$x_{r-1}$ to $y_s$ with $r-1-s$ finite ascents, all of whose elements lie in $ {}^IW$ by \Cref{lem:endpoints}.
\Cref{cor:weak-valuation} therefore gives
\begin{equation}\label{eq:second-coarse}
\ord\bigl(h_{y_s}\bigr)\geq r-(r-1-s)=s+1.  
\end{equation}
Since $z$ belongs to the ordered case, we have $z=\Theta_j$ with $0\leq j\leq b-1=r-t-2$.
If $j=0$, then $z=e$, and \Cref{lem:mu} together with \cref{eq:second-coarse} give
\[
 \mus(e,y_{t+1})=[v]h_{y_{t+1}}=0.
\]
Suppose now that $j\geq1$.
Then part (ii) of \Cref{lem:rank-reduction} yields
\[
\mus(\Theta_j,y_{t+1})=
[v]\left([M_{\Theta_j}]\C_{y_{t+1}}\right)=
[v]h_{y_{t+1}}^{(r-j)}.
\]
Since $r-j\geq t+2$, the element $y_{t+1}^{(r-j)}$ is well-defined, and thus $\ord\bigl(h_{y_{t+1}}^{(r-j)}\bigr)\geq2$, by \cref{eq:second-coarse}.
Its coefficient of $v$ therefore vanishes, proving $\mus(z,y_{t+1})=0$.
\end{proof}

We now treat the pure cases.

\begin{lemma}\label{lem:pure}
Let $z\in\mathcal N_t$ belong to one of the pure cases.
Then $\ord(h_z)\geq t+3$.
\end{lemma}
\begin{proof}
Let $z=c_\omega$ with $\omega$ of type $j$, and put $n=N-2j$.
In either pure case, $1\leq j\leq r$.
Our strategy is to find a length-additive factorization
\[
z=ds_{i_1}\cdots s_{i_L}, \qquad s_{i_k}\in I,
\]
with $d\in{}^IW$ for which both $h_{d}$ and $L$ can be computed explicitly.
Since $z\in{}^IW$, all intermediate elements also lie in ${}^IW$, and \Cref{cor:weak-valuation} will then give
\begin{equation}\label{eq:purecase-step}
 \ord(h_z)\geq\ord(h_{d})-L. 
\end{equation}

We first construct the starting elements.
For $1\leq s<N$, let
\[
 d_s=w_{I\cap P_s}w_{P_s}, \qquad P_s=S\setminus\{s_s\}.
\]
This is the unique longest element of $ {}^{I\cap P_s}W_{P_s}$.
We have $\ell(d_s)=s(N-s)$, since the factorization $w_{P_s}=w_{I\cap P_s}\,d_s$ is length-additive.
By \cref{eq:BSP}, we have the expansion
\[
\D(P_s)=M_{d_s}+ \sum_{\substack{u\in{}^{I\cap P_s}W_{P_s}\\u\ne d_s}} v^{s(N-s)-\ell(u)}M_u,
\]
where every exponent in the sum is positive.
\Cref{lem:overlaps} gives bar-invariance, so the canonical basis criterion in Lemma~\ref{lem:criterion} yields
\[
 \D(P_s)=\C_{d_s},\qquad h_{d_s}=[M_e]\D(P_s)=v^{s(N-s)}.
\]

Now consider the words
\[
\B^n\A^j\CC^j,\qquad \A^j\CC^j\B^n.
\]
Both are admissible of type $j$, and have inversion number $j(N-j)$ for the order $\CC<\B<\A$.
By \Cref{lem:initial-supports}, the corresponding elements belong to the supports of $\D(P_{N-j})$ and $\D(P_j)$, respectively.
By \Cref{lem:word-length}, both have length $j(N-j)$, so they are the unique longest elements $d_{N-j}$ and $d_j$ of these supports.

We take $d=d_{N-j}$ in the pure $\B\CC$ case and $d=d_j$ in the pure $\A\B$ case.
In either case, $h_{d}=v^{j(N-j)}$.
It remains to construct the path from $d$ to $z$ and determine its length $L$.
Each adjacent interchange
\[
 \A\B\longmapsto\B\A,\qquad
 \A\CC\longmapsto\CC\A,\qquad
 \B\CC\longmapsto\CC\B
\]
reduces the inversion number by one.
By \Cref{lem:words,lem:word-length,lem:word-properties}, it corresponds to right multiplication by a finite simple reflection and increases length by one.
We construct the path using these interchanges.
Recall that $c\leq b-1=r-t-2$.

\smallskip\noindent
\textbf{Pure $\B\CC$ case.}
Here $a=0$ and $c<j$.
We have
\[
\B^n\A^j\CC^j\ \longmapsto\
\B^n\CC^j\A^j\ \longmapsto\
\B^{r+1-c}\CC^{j}\B^{r-2j+c}\A^j\ \longmapsto\
\CC^c\B^{r+1-c}\,\bigm|\,
 \CC^{j-c}\B^{r-2j+c}\A^j,
\]
where the first and last words correspond to $d=d_{N-j}$ and $z$, respectively.
Each arrow is a succession of adjacent interchanges of the types listed above, so we get a length-additive factorization $z=ds_{i_1}\cdots s_{i_L}$ with $s_{i_k}\in I$ and
\[
L=j^2+j(r-2j+c)+c(r+1-c)
=j(N-j)-(r+1-c)(j-c).
\]
In particular,
\[
j(N-j)-L\geq  r+1-c\geq t+3.
\]

\smallskip\noindent
\textbf{Pure $\A\B$ case.}
Here $a>0$ and $c=j$.
Similarly, we have
\[
\A^j\CC^j\B^n\ \longmapsto\
\CC^j\A^j\B^n\ \longmapsto\
\CC^j\B^{r+1-j-a}\A^j\B^{r-j+a}\ \longmapsto\
\CC^j\B^{r+1-j-a}\A^a\,\bigm|\,
\B^{r-j+a}\A^{j-a},
\]
where the first and last words correspond to $d=d_{j}$ and $z$, respectively.
We thus get a length-additive factorization $z=ds_{i_1}\cdots s_{i_L}$ with $s_{i_k}\in I$ and
\[
L=j^2+j(r+1-j-a)+(j-a)(r-j+a)
 =j(N-j)-a(r-j+a).
\]
In particular,
\[
j(N-j)-L\geq r-j+a\geq r-j+1\geq t+3.
\]

\smallskip
Finally, combining $h_{d}=v^{j(N-j)}$ and \cref{eq:purecase-step}, we get
\[
 \ord(h_z)  \geq j(N-j)-L \geq t+3,
\]
as required.
\end{proof}

To treat the mixed case, we recall right star operations and the invariance property needed below.

\begin{definition}\label{def:star}
For distinct simple reflections $s,t$ with $sts=tst$, let
\[
D_R(s,t)=\{w\in W:|R(w)\cap\{s,t\}|=1\}.
\]
For $w\in D_R(s,t)$, exactly one of $ws,wt$ belongs to
$D_R(s,t)$; that element is denoted by $w^*$.
Equivalently, in the right parabolic factor from $W_{\{s,t\}}$
this involution exchanges $s\leftrightarrow st$ and
$t\leftrightarrow ts$.
\end{definition}

\begin{proposition}\label{prop:star}
Let $Z,Y\in D_R(s,t)$ be such that $Z<Y$ in the Bruhat order.
Suppose that $\ell(Z^*)=\ell(Z)+1$, $\ell(Y^*)=\ell(Y)-1$, and $\ell(Y)-\ell(Z)\geq3$.
Then
\[
 \mu(Z,Y)=\mu(Z^*,Y^*),
\]
where the right side is zero if the starred pair is not Bruhat-comparable.
\end{proposition}
\begin{proof}
Since $\ell(Y)-\ell(Z)\geq3$, $Z,Y$ lie in different right $W_{\{s,t\}}$-cosets: within one such coset elements of $D_R(s,t)$ have parabolic factors of length one or two, so their lengths differ by at most one.
By \cite[Theorem~6.2]{Casselman}, we have $Z\prec Y$ if and only if $Z^*\prec Y^*$, where $u\prec v$ means $u<v$ in Bruhat order and $\mu(u,v)\ne0$.
When these conditions hold, the same theorem gives $\mu(Z,Y)=\mu(Z^*,Y^*)$.
Otherwise, both sides vanish, with the stated convention, since
\[
 \ell(Y^*)-\ell(Z^*)
 =\ell(Y)-\ell(Z)-2\geq1
\]
rules out $Y^*<Z^*$.
\end{proof}

\begin{lemma}\label{lem:mixed}
Let $z\in\mathcal N_t$ belong to the mixed case.
Then $\mus(z,y_{t+1})=0$.
\end{lemma}
\begin{proof}
Suppose that $\mus(z,y_{t+1})\ne0$, and write $z=c_\omega$.
Since $z$ belongs to the mixed case, we have
$\omega_{r+1}\omega_{r+2}=\A\CC$.
Let $j\leq r$ be the type of $\omega$.
We first check that the letters at positions $r,\ldots,r+3$ have the form
\[
 X\A\CC T,\qquad X\in\{\A,\B\},\quad T\in\{\B,\CC\}.
\]
If $X=\CC$, the first ordered block contains at least $r$ letters $\CC$, with another at position $r+2$.
This would force $j\geq r+1$.
If $T=\A$, the last ordered block is $\CC\A^{r-1}$. Together with the $\A$ at position $r+1$, this forces $j=r$.
Therefore,
\[
 \omega=\CC^{r-1}\B\A\,\bigm|\,\CC\A^{r-1}.
\]
This contradicts $\omega_b\ne\CC$, since $1\leq b\leq r-1$.

Set $Z=w_Iz$ and $Y=w_Iy_{t+1}$.
We will apply a right star operation and then use the descent restriction to force an equality incompatible with the ordered first block of $\omega$.

By \Cref{def:second-chain}, the word of $y_{t+1}$ is
\[
 \CC^{b-1}\A\CC^{t+1}\B\CC\A^{r-1}.
\]
Since
\[
 w_I(k+qN)=N+1-k+qN
 \qquad(1\leq k\leq N,\ q\in\mathbb Z),
\]
left multiplication by $w_I$ reverses the numerical order within each list in \eqref{eq:lists}.
Entries corresponding to $\A$ remain smaller than those corresponding to $\B$, which remain smaller than those corresponding to $\CC$.
Since the letters of $z$ and $y_{t+1}$ at positions
$r,\ldots,r+3$ are, respectively,
\[
 X\A\CC T
 \qquad\text{and}\qquad
 \CC\B\CC\A,
\]
we must have
\[
Z(r+1)<Z(r)<Z(r+2),
\qquad
Y(r+1)<Y(r+2)<Y(r).
\]
The first inequality holds also when $X=\A$, since successive occurrences of $\A$ have decreasing entries in $Z$.
The triples of entries at positions $r,r+1,r+2$ have relative order $213$ and $312$, respectively.
Thus $Z,Y\in D_R(s_r,s_{r+1})$ and
\[
 Z^*=Zs_{r+1},\qquad Y^*=Ys_r,
 \qquad
 \ell(Z^*)=\ell(Z)+1,\quad
 \ell(Y^*)=\ell(Y)-1.
\]
We have $Z<Y$ and $\mu(Z,Y)=\mus(z,y_{t+1})\ne0$,
so $\ell(Y)-\ell(Z)$ is odd.
Moreover,
\[
 \ell(Y)-\ell(Z)=\ell(y_{t+1})-\ell(z)>1
\]
because $z\in\mathcal N_t$.
Thus this difference is at least three.
We may therefore apply \Cref{prop:star}, obtaining
\[
 Z^*<Y^*,\qquad \mu(Z^*,Y^*)\ne0.
\]

Put $z'=zs_{r+1}$ and $y'=y_{t+1}s_r$.
In the words of $z$ and $y_{t+1}$, these swaps exchange adjacent distinct letters, so the resulting admissible words correspond to the elements $z'$ and $y'$.
Thus, $z',y'\in{}^IW$ by \Cref{lem:words}.
Since $Z^*=w_Iz'$ and $Y^*=w_Iy'$, we have
\[
 z'<y',\qquad \mus(z',y')\ne0.
\]
Their words at positions $r,\ldots,r+3$ are now
\[
 X\CC\A T
 \qquad\text{and}\qquad
 \B\CC\CC\A,
\]
respectively.
Since $T\in\{\B,\CC\}$, the descent test \Cref{lem:word-properties} gives
\[
 s_{r+2}\in R_I(y')\setminus R_I(z').
\]
Then, \Cref{lem:descents} yields
\[
 y'=z's_{r+2}=zs_{r+1}s_{r+2}.
\]

The two swaps on the right exchange distinct letters and
leave the first $r$ letters of $\omega$ unchanged.
Thus the word of $zs_{r+1}s_{r+2}$ has an ordered prefix
of length $r$.
However, the word of $y'$ has prefix
\[
 \CC^{b-1}\A\CC^t\B,
\]
which has length $r$ and is not ordered, since its $\A$ precedes its $\B$.
This contradicts the uniqueness of the admissible word in \Cref{lem:words}.
\end{proof}

We are now in a position to complete the proof of \Cref{prop:second-multiplication}.

\begin{proof}[Proof of \Cref{prop:second-multiplication}]
Fix $r\geq2$ and $0\leq t\leq r-2$, and define
\[
 \mathcal E_t=\sum_{z\in\mathcal N_t}
       \mus(z,y_{t+1})\C_z.
\]
The candidate expansion \eqref{eq:second-chain-candidate-expansion}
gives \eqref{eq:second-multiplication}, with the cover term
$\C_{y_{t+2}}$ omitted when $t=r-2$.
It remains to prove that $\ord(\eps(\mathcal E_t))\geq t+3$, as explained in the beginning of this section.

By \Cref{lem:second-coarse,lem:mixed}, the coefficients $\mus(z,y_{t+1})$ vanish in the ordered and mixed cases.
Hence
\[
 \eps(\mathcal E_t)
 =\sum_{\substack{z\in\mathcal N_t\\
                  z\text{ belongs to a pure case}}}
       \mus(z,y_{t+1})h_z.
\]
By \Cref{lem:pure}, every $h_z$ in this sum has order at least $t+3$. Since the coefficients $\mus(z,y_{t+1})$ are integers, it follows that
\[
 \ord\bigl(\eps(\mathcal E_t)\bigr)\geq t+3,
\]
as required.
\end{proof}

% {\color{blue}
% \begin{proof}[Proof of \Cref{prop:second-multiplication}]
% Fix $r\geq2$ and $0\leq t\leq r-2$, and put $b=r-1-t$. Define
% \[
%  \mathcal E_t=\sum_{z\in\mathcal N_t}
%        \mus(z,y_{t+1})\C_z.
% \]
% The candidate expansion \eqref{eq:second-chain-candidate-expansion}
% already gives \eqref{eq:second-multiplication}, with the cover
% $y_{t+2}$ omitted when $t=r-2$. It remains to estimate
% $\eps(\mathcal E_t)$, as in \eqref{eq:second-chain-required-order}.
% We treat the four cases listed above.

% First suppose that the word is ordered, so $z=\Theta_j$.
% The prefix condition gives $0\leq j\leq r-t-2$.
% For $j=0$, we have $\Theta_0=e$, and Lemmas~\ref{lem:mu}
% and~\ref{lem:second-coarse} give
% \[
%  \mus(e,y_{t+1})=[v]h_{y_{t+1}}=0,
% \]
% since $\ord(h_{y_{t+1}})\geq t+2\geq2$.
% If $j\geq1$, then $t\leq r-3$. Put $m=r-j$. We have
% $m\geq t+2$, so $0\leq t+1\leq m-1$ and $y_{t+1}^{(m)}$ is
% defined. Applying the rank-reduction identity
% \eqref{eq:rank-reduction-second-chain} and then
% Lemma~\ref{lem:second-coarse} with rank parameter $m$ yields
% \begin{equation}\label{eq:second-chain-correction-multiplicities}
%  \mus(\Theta_j,y_{t+1})
%  =[v]h_{y_{t+1}}^{(m)}=0,
% \end{equation}
% because the latter polynomial has order at least $t+2$.

% For every $z\in\mathcal N_t$, positions $b,b+1$ lie in the first
% ordered block, since $1\leq b\leq r-1$. Thus
% $s_b\in R_I(z)$ by \eqref{eq:word-descents}.
% Together with \eqref{eq:remaining-second-chain-indices}, this shows
% that every $z\in\mathcal N_t$ with $\mus(z,y_{t+1})\ne0$ belongs to
% $\mathcal C_{\mathrm{nc}}(y_{t+1},s_b)$.
% Lemma~\ref{lem:mixed} therefore makes all mixed contributions vanish. The only remaining terms come from
% the two pure cases. By Lemma~\ref{lem:pure}, each such $h_z$ has
% order at least $t+3$. Since $\mus(z,y_{t+1})$ is an integer,
% \[
%  \eps(\mathcal E_t)
%  =\sum_{\substack{z=c_\omega\in\mathcal N_t\\
%         \omega_{r+1}\omega_{r+2}\in\{\B\CC,\A\B\}}}
%        \mus(z,y_{t+1})h_z
%  =O(v^{t+3}).
% \]
% This proves the required remainder estimate.
% The estimates used in smaller ranks come from
% Lemma~\ref{lem:second-coarse}, whose proof is independent of the
% second-chain formulas, so no induction on the rank is needed.
% \end{proof}
%}


% \newpage
% \input{fund_box}

% Bibliographic records are maintained in bibliography.bib.
% Only references cited in the manuscript are included below.
\bibliographystyle{plain}
\bibliography{bibliography}

\end{document}
